\documentclass[12pt,a4paper]{article}
\usepackage[utf8]{inputenc}
\usepackage[T1]{fontenc}
\usepackage[french]{babel}
\usepackage{geometry}
\geometry{left=2.5cm, right=2.5cm, top=2.5cm, bottom=2.5cm}
\usepackage{amsmath, amssymb, amsthm}
\usepackage{graphicx}
\usepackage{array}
\usepackage{multirow}
\usepackage{colortbl}
\usepackage{xcolor}
\usepackage{tikz}
\usetikzlibrary{shapes, arrows, positioning, calc, decorations.pathreplacing}
\usepackage{hyperref}
\hypersetup{
    colorlinks=true,
    linkcolor=blue,
    urlcolor=blue,
}
\usepackage{enumitem}
\usepackage{float}
\usepackage{booktabs}
\usepackage{caption}
\usepackage{subcaption}

% Définition des couleurs pour les sections
\definecolor{titlecolor}{RGB}{0, 79, 144}
\definecolor{sectioncolor}{RGB}{0, 112, 192}
\definecolor{methodcolor}{RGB}{0, 150, 130}
\definecolor{leancolor}{RGB}{200, 70, 0}

% Style des titres
\usepackage{titlesec}
\titleformat{\section}
  {\color{sectioncolor}\normalfont\Large\bfseries}
  {\thesection}{1em}{}
\titleformat{\subsection}
  {\color{methodcolor}\normalfont\large\bfseries}
  {\thesubsection}{1em}{}
\titleformat{\subsubsection}
  {\color{leancolor}\normalfont\normalsize\bfseries}
  {\thesubsubsection}{1em}{}

% Environnement pour les notes importantes
\usepackage{tcolorbox}
\tcbuselibrary{skins, breakable}
\newenvironment{importantnote}
  {\begin{tcolorbox}[colback=blue!5!white, colframe=blue!50!black, title=\textbf{Point Important}]\small}
  {\end{tcolorbox}}

% Environnement pour les méthodes
\newenvironment{methodbox}[1]
  {\begin{tcolorbox}[colback=blue!5!white, colframe=methodcolor, title=#1]\small}
  {\end{tcolorbox}}

% Commande pour les boîtes de conseil
\newenvironment{conseil}
  {\begin{tcolorbox}[colback=green!5!white, colframe=green!50!black, title=\textbf{Conseil}]\small}
  {\end{tcolorbox}}

\title{\Huge\textbf{Équilibrage des Postes de Travail} \\
        \Large Méthodologie, outils et lien avec le Lean Manufacturing}
\author{Guide Méthodologique}
\date{\today}

\begin{document}

\maketitle
\tableofcontents
\newpage

% ============================================================
% PARTIE 1: INTRODUCTION À L'ÉQUILIBRAGE DES LIGNES DE PRODUCTION
% ============================================================

\section{Introduction à l'équilibrage des lignes de production}

L'équilibrage des lignes de production (Line Balancing) est une méthode fondamentale du Lean Manufacturing qui vise à répartir équitablement la charge de travail entre les différents postes d'une chaîne de production. Cette approche permet de minimiser les temps morts, d'optimiser l'utilisation des ressources et d'améliorer le flux de production.

\subsection{Définition et objectifs}

L'équilibrage de ligne consiste à regrouper les opérations élémentaires (tâches) en postes de travail de telle sorte que :

\begin{itemize}
    \item Chaque poste ait une durée d'activité la plus proche possible du \textbf{Takt Time}
    \item Les contraintes de précédence entre les tâches soient respectées
    \item Le nombre de postes soit minimisé
    \item Les temps morts (déséquilibre) soient réduits au maximum
\end{itemize}

\begin{importantnote}
Un équilibrage parfait est atteint lorsque la durée de chaque poste est exactement égale au Takt Time, éliminant ainsi toute attente. Cette situation idéale est rarement atteignable dans la pratique.
\end{importantnote}

\subsubsection{Le concept de goulot d'étranglement}

Un goulot d'étranglement (bottleneck) est un poste de travail qui reçoit plus de travail que ce qu'il peut traiter à sa capacité de production maximale. Cela provoque une interruption du flux de travail et des retards dans le processus de production. L'équilibrage vise précisément à identifier et à résoudre ces goulots.

\subsection{Différence entre Cycle Time, Lead Time et Takt Time}

Pour bien comprendre l'équilibrage, il est essentiel de distinguer ces trois concepts fondamentaux qui sont souvent confondus.

\subsubsection{Takt Time ($T_T$)}

Le Takt Time est le rythme de production dicté par la demande client. Il représente le temps maximum autorisé pour produire une unité afin de répondre à la demande.

\[
\boxed{T_T = \frac{\text{Temps de production disponible par période}}{\text{Demande client par période}}}
\]

\textbf{Caractéristiques :}
\begin{itemize}
    \item C'est une donnée objective basée sur la demande client
    \item Il sert de référence pour l'équilibrage
    \item Il peut varier selon les fluctuations de la demande
    \item Exemple : $T_T = \dfrac{7h \times 60}{600 \text{ pièces}} = 0,7 \text{ min/pièce}$
\end{itemize}

\subsubsection{Cycle Time ($T_C$)}

Le Cycle Time est le temps réel entre deux pièces produites sur un poste donné. Après équilibrage, l'objectif est que le cycle time de chaque poste soit proche du Takt Time.

\[
\boxed{T_C = \frac{\text{Temps de production d'une pièce sur un poste}}{\text{Nombre de pièces produites}}}
\]

\textbf{Caractéristiques :}
\begin{itemize}
    \item C'est une mesure réelle de la performance
    \item Si $T_C < T_T$ : Le poste a une surcapacité (temps morts)
    \item Si $T_C > T_T$ : Le poste est un goulot (incapacité à suivre la demande)
    \item Si $T_C = T_T$ : Équilibre parfait (situation idéale)
\end{itemize}

\subsubsection{Lead Time ($LT$)}

Le Lead Time est le temps total entre l'entrée de la matière première et la sortie du produit fini, incluant tous les temps de transformation, d'attente et de transport.

\[
\boxed{LT = \sum (\text{Temps de transformation}) + \sum (\text{Temps d'attente}) + \sum (\text{Temps de transport})}
\]

\textbf{Caractéristiques :}
\begin{itemize}
    \item C'est le temps que le client perçoit
    \item Un bon équilibrage contribue à réduire le lead time
    \item Relation : $LT = N \times T_C + \text{temps d'attente}$ (où $N$ est le nombre de postes)
\end{itemize}

\subsubsection{Tableau récapitulatif}

\begin{table}[H]
\centering
\begin{tabular}{|p{2.5cm}|p{3.5cm}|p{3.5cm}|p{4cm}|}
\hline
\textbf{Concept} & \textbf{Définition} & \textbf{Formule} & \textbf{Rôle dans l'équilibrage} \\
\hline
Takt Time & Rythme imposé par le client & $T_T = \frac{T_{dispo}}{D}$ & Objectif à atteindre \\
\hline
Cycle Time & Temps réel de production & $T_C = \frac{T_{prod}}{Q}$ & Indicateur d'écart \\
\hline
Lead Time & Temps total de traversée & $LT = \sum T_C + \sum \text{attente}$ & Indicateur global \\
\hline
\end{tabular}
\caption{Différences fondamentales entre Takt Time, Cycle Time et Lead Time}
\end{table}

\subsubsection{Illustration graphique des concepts}

\begin{center}
\begin{tikzpicture}[node distance=0.8cm, auto]
    % Timeline
    \draw[thick, ->] (0,0) -- (12,0) node[right] {Temps};
    
    % Lead Time
    \draw[decorate, decoration={brace, amplitude=10pt, mirror}] (0.5,-0.5) -- (11.5,-0.5) node[midway, below] {Lead Time};
    
    % Postes
    \fill[blue!20] (1,0) rectangle (3,0.5) node[midway, above] {Poste 1\\$T_C=0.62$};
    \fill[blue!20] (3.5,0) rectangle (5.5,0.5) node[midway, above] {Poste 2\\$T_C=0.66$};
    \fill[blue!20] (6,0) rectangle (8,0.5) node[midway, above] {Poste 3\\$T_C=0.70$};
    \fill[blue!20] (8.5,0) rectangle (10.5,0.5) node[midway, above] {Poste 4\\$T_C=0.63$};
    
    % Attentes
    \draw[<->, red] (3,0.2) -- (3.5,0.2) node[midway, above] {Attente};
    \draw[<->, red] (5.5,0.2) -- (6,0.2) node[midway, above] {Attente};
    \draw[<->, red] (8,0.2) -- (8.5,0.2) node[midway, above] {Attente};
    
    % Takt Time line
    \draw[dashed, green!60!black, thick] (0,0.8) -- (12,0.8) node[right] {Takt Time = 0,7 min};
    
    \node[text width=3cm, align=center] at (6,-1.2) {Le Cycle Time du poste le plus lent\\détermine le rythme de production};
\end{tikzpicture}
\captionof{figure}{Relation entre Cycle Time, Lead Time et Takt Time}
\end{center}

\subsection{Pourquoi équilibrer une ligne ?}

L'équilibrage des lignes de production répond à des enjeux majeurs de performance industrielle.

\subsubsection{Les bénéfices d'un bon équilibrage}

\begin{enumerate}
    \item \textbf{Réduction des temps morts}
    \begin{itemize}
        \item Diminution des attentes entre les postes
        \item Meilleure utilisation du temps disponible
        \item Gain de productivité significatif
    \end{itemize}
    
    \item \textbf{Amélioration du flux de production}
    \begin{itemize}
        \item Fluidification des opérations
        \item Réduction des ruptures de flux
        \item Meilleure synchronisation des postes
    \end{itemize}
    
    \item \textbf{Optimisation des ressources}
    \begin{itemize}
        \item Meilleure utilisation des opérateurs
        \item Optimisation du parc machines
        \item Réduction des besoins en surfaces
    \end{itemize}
    
    \item \textbf{Réduction des stocks intermédiaires (WIP)}
    \begin{itemize}
        \item Moins d'en-cours entre les postes
        \item Réduction du capital immobilisé
        \item Meilleure visibilité des flux
    \end{itemize}
    
    \item \textbf{Augmentation de la productivité}
    \begin{itemize}
        \item Production plus efficace
        \item Meilleur rendement global
        \item Augmentation du taux de sortie
    \end{itemize}
    
    \item \textbf{Meilleure visibilité des problèmes}
    \begin{itemize}
        \item Identification claire des goulots
        \item Détection rapide des anomalies
        \item Base pour l'amélioration continue
    \end{itemize}
\end{enumerate}

\subsubsection{Les coûts du déséquilibre}

Un déséquilibre de ligne génère des pertes significatives :

\begin{table}[H]
\centering
\begin{tabular}{|p{4cm}|p{8cm}|}
\hline
\textbf{Type de perte} & \textbf{Conséquences} \\
\hline
Temps morts & Sous-utilisation des opérateurs et des machines \\
\hline
Stocks WIP & Capital immobilisé, risque d'obsolescence, surfaces occupées \\
\hline
Lead time long & Manque de réactivité, délais clients allongés \\
\hline
Goulots & Production limitée par le poste le plus lent \\
\hline
Surcapacité & Investissements sous-optimisés \\
\hline
\end{tabular}
\caption{Conséquences financières et opérationnelles du déséquilibre}
\end{table}

\subsubsection{L'impact sur la performance globale}

Un équilibrage efficace permet d'atteindre des objectifs stratégiques :

\begin{center}
\begin{tikzpicture}[node distance=1.2cm]
    \node[draw, rectangle, fill=blue!10, minimum width=3cm, minimum height=0.8cm] (lean) at (0,0) {Lean Manufacturing};
    \node[draw, rectangle, fill=green!10, minimum width=3cm, minimum height=0.8cm, below of=lean] (equi) {Équilibrage};
    \node[draw, rectangle, fill=yellow!10, minimum width=2.5cm, minimum height=0.8cm, below left of=equi, xshift=-1.5cm] (prod) {Productivité};
    \node[draw, rectangle, fill=yellow!10, minimum width=2.5cm, minimum height=0.8cm, below of=equi] (qual) {Qualité};
    \node[draw, rectangle, fill=yellow!10, minimum width=2.5cm, minimum height=0.8cm, below right of=equi, xshift=1.5cm] (delai) {Délais};
    \node[draw, rectangle, fill=red!10, minimum width=4cm, minimum height=0.8cm, below of=prod, xshift=2cm] (client) {Satisfaction Client};
    
    \draw[->, thick] (lean) -- (equi);
    \draw[->, thick] (equi) -- (prod);
    \draw[->, thick] (equi) -- (qual);
    \draw[->, thick] (equi) -- (delai);
    \draw[->, thick] (prod) -- (client);
    \draw[->, thick] (qual) -- (client);
    \draw[->, thick] (delai) -- (client);
\end{tikzpicture}
\captionof{figure}{Impact de l'équilibrage sur la performance globale}
\end{center}

\subsubsection{Quand faut-il équilibrer ?}

L'équilibrage doit être envisagé dans les situations suivantes :
\begin{itemize}
    \item Création d'une nouvelle ligne de production
    \item Changement significatif de la demande client
    \item Introduction d'un nouveau produit
    \item Identification de goulots récurrents
    \item Taux d'efficacité inférieur à 80\%
    \item Stocks intermédiaires anormalement élevés
    \item Après une amélioration significative (SMED, automatisation)
\end{itemize}

\begin{importantnote}
L'équilibrage n'est pas une action ponctuelle mais une démarche d'amélioration continue. Les lignes de production doivent être régulièrement rééquilibrées pour s'adapter aux évolutions de la demande et aux améliorations terrain.
\end{importantnote}
% ============================================================
% PARTIE 2: INFORMATIONS NÉCESSAIRES POUR RÉALISER UN ÉQUILIBRAGE
% ============================================================

\newpage
\section{Informations nécessaires pour réaliser un équilibrage}

Avant de commencer l'équilibrage, il est crucial de collecter des données précises et complètes. La qualité de l'équilibrage dépend directement de la fiabilité des informations collectées.

\subsection{Données de production}

\subsubsection{La demande client}

La demande client est le point de départ de tout équilibrage. Elle permet de calculer le Takt Time et de dimensionner la ligne.

\begin{methodbox}{Données à collecter sur la demande}
\begin{itemize}
    \item Quantité journalière/hebdomadaire/mensuelle à produire
    \item Variations saisonnières et tendances
    \item Mix produit (si plusieurs références)
    \item Niveau de service attendu (taux de disponibilité)
    \item Prévisions à moyen et long terme
\end{itemize}
\end{methodbox}

\subsubsection{Temps de production disponible}

Le temps disponible est le temps pendant lequel la production peut effectivement fonctionner.

\[
\boxed{T_{dispo} = \text{Temps ouvré} - \text{Pauses} - \text{Réunions} - \text{Maintenance préventive}}
\]

\begin{methodbox}{Éléments à prendre en compte}
\begin{itemize}
    \item Nombre d'équipes et horaires de travail
    \item Durée des pauses (déjeuner, pauses café)
    \item Temps de réunion d'équipe (5S, briefs quotidiens)
    \item Temps de maintenance préventive programmée
    \item Taux de disponibilité des équipements
\end{itemize}
\end{methodbox}

\subsubsection{Temps de cycle des opérations}

Le temps de cycle de chaque opération est la donnée fondamentale pour l'équilibrage.

\begin{importantnote}
La mesure des temps de cycle doit être réalisée dans des conditions réelles de production, en tenant compte des variations normales. Il est recommandé de prendre au moins 10 mesures par opération et d'utiliser la moyenne ou le temps médian.
\end{importantnote}

\textbf{Méthodes de mesure :}
\begin{itemize}
    \item \textbf{Chronométrage direct} : Mesure avec un chronomètre sur le terrain
    \item \textbf{Étude des temps} : Analyse des données historiques
    \item \textbf{Échantillonnage de travail} : Observation aléatoire pour déterminer les temps
    \item \textbf{MTM (Methods-Time Measurement)} : Temps prédéterminés selon les mouvements
\end{itemize}

\subsubsection{Contraintes de précédence}

Les contraintes de précédence définissent l'ordre obligatoire des opérations. Certaines tâches ne peuvent commencer qu'après la fin d'autres tâches.

\begin{methodbox}{Types de contraintes de précédence}
\begin{itemize}
    \item \textbf{Séquence obligatoire} : La tâche B ne peut commencer qu'après la tâche A
    \item \textbf{Contrainte technique} : Séchage, refroidissement, etc.
    \item \textbf{Contrainte de montage} : Ordre d'assemblage imposé par la conception
    \item \textbf{Contrainte de sécurité} : Certaines opérations nécessitent des prérequis
    \item \textbf{Contrainte d'outillage} : Partage d'outils ou de machines spécifiques
\end{itemize}
\end{methodbox}

\subsection{Méthodes de collecte des données}

\subsubsection{Observation directe sur le terrain (Gemba)}

L'observation directe est la méthode la plus fiable pour collecter des données réelles.

\begin{conseil}
\noindent\textbf{Conseils pour l'observation :}
\begin{enumerate}
    \item Se rendre sur le terrain (Gemba) aux heures de production normale
    \item Observer plusieurs cycles pour capter les variations
    \item Prendre en compte les temps de déplacement et de manipulation
    \item Noter les anomalies et interruptions
    \item Discuter avec les opérateurs pour comprendre les difficultés
    \item Photographier ou filmer pour analyse ultérieure
\end{enumerate}
\end{conseil}

\subsubsection{Analyse des gammes de fabrication}

Les gammes de fabrication sont des documents techniques qui décrivent :
\begin{itemize}
    \item La séquence des opérations
    \item Les machines utilisées
    \item Les temps standards (souvent optimistes)
    \item Les outillages nécessaires
    \item Les compétences requises
\end{itemize}

\subsubsection{Entretiens avec les opérateurs}

Les opérateurs sont les experts de leur poste de travail. Leur expérience est précieuse.

\begin{methodbox}{Questions à poser aux opérateurs}
\begin{itemize}
    \item Quelles sont les difficultés rencontrées ?
    \item Y a-t-il des temps d'attente réguliers ?
    \item Quels sont les goulots identifiés ?
    \item Des améliorations ont-elles déjà été testées ?
    \item Quelles seraient vos propositions d'amélioration ?
\end{itemize}
\end{methodbox}

\subsubsection{Analyse des données historiques}

Les données de production historiques permettent de valider les mesures ponctuelles :
\begin{itemize}
    \item Productivité réelle par opérateur
    \item Taux de rendement synthétique (TRS / OEE)
    \item Temps d'arrêt et pannes
    \item Niveaux de stocks intermédiaires
    \item Délais de production observés
\end{itemize}

\subsection{Préparation du diagramme d'antériorité}

Le diagramme d'antériorité (ou graphe de précédence) est un outil essentiel qui représente visuellement :
\begin{itemize}
    \item Chaque tâche avec son temps de cycle
    \item Les relations de précédence (flèches)
    \item La séquence logique de production
    \item Les tâches pouvant être réalisées en parallèle
\end{itemize}

\subsubsection{Construction du diagramme d'antériorité}

\begin{enumerate}
    \item \textbf{Lister toutes les tâches} avec leur temps de cycle
    \item \textbf{Identifier les prédécesseurs} pour chaque tâche
    \item \textbf{Dessiner les nœuds} représentant chaque tâche
    \item \textbf{Tracer les flèches} reliant les tâches selon les contraintes
    \item \textbf{Vérifier la cohérence} (absence de cycles)
\end{enumerate}

\subsubsection{Exemple de diagramme d'antériorité}

\begin{center}
\begin{tikzpicture}[node distance=1.5cm, auto, task/.style={rectangle, draw, minimum width=1.2cm, minimum height=0.8cm, align=center, fill=blue!10}]
    % Nœuds
    \node[task] (A) at (0,0) {A\\0,2 min};
    \node[task] (B) at (2,0) {B\\0,4 min};
    \node[task] (C) at (1,-1.5) {C\\0,7 min};
    \node[task] (D) at (-0.5,-3) {D\\0,1 min};
    \node[task] (E) at (2.5,-3) {E\\0,3 min};
    \node[task] (F) at (1,-4.5) {F\\0,2 min};
    
    % Flèches
    \draw[->, thick] (A) -- (C);
    \draw[->, thick] (B) -- (C);
    \draw[->, thick] (C) -- (D);
    \draw[->, thick] (C) -- (E);
    \draw[->, thick] (D) -- (F);
    \draw[->, thick] (E) -- (F);
    
    % Légende
    \node[text width=4cm, align=center, below] at (1,-5.8) {Figure : Exemple de diagramme d'antériorité\\Les flèches indiquent l'ordre des opérations};
\end{tikzpicture}
\end{center}

\subsubsection{Règles de construction}

\begin{importantnote}
\textbf{Règles fondamentales :}
\begin{enumerate}
    \item Une tâche ne peut commencer que lorsque toutes ses tâches prédécesseurs sont terminées
    \item Le diagramme ne doit pas contenir de cycle (une tâche ne peut pas être antérieure à elle-même)
    \item Les tâches sans prédécesseur peuvent commencer au début du processus
    \item Le temps total de traitement est la somme de tous les temps de cycle
\end{enumerate}
\end{importantnote}

\subsection{Formulaire de collecte des données}

Voici un formulaire type à utiliser pour la collecte des données sur le terrain :

\begin{table}[H]
\centering
\begin{tabular}{|c|p{3cm}|p{2cm}|p{2.5cm}|p{3cm}|}
\hline
\textbf{Réf.} & \textbf{Description tâche} & \textbf{Temps (min)} & \textbf{Prédécesseurs} & \textbf{Observations} \\
\hline
A & & & & \\
\hline
B & & & & \\
\hline
C & & & & \\
\hline
D & & & & \\
\hline
E & & & & \\
\hline
F & & & & \\
\hline
G & & & & \\
\hline
H & & & & \\
\hline
\end{tabular}
\caption{Formulaire de collecte des données d'équilibrage}
\end{table}

\subsection{Check-list pré-équilibrage}

Avant de commencer l'équilibrage, vérifiez les points suivants :

\begin{table}[H]
\centering
\begin{tabular}{|p{8cm}|c|}
\hline
\textbf{Point à vérifier} & \textbf{OK} \\
\hline
Demande client validée pour la période & $\square$ \\
\hline
Temps de production disponible calculé & $\square$ \\
\hline
Takt Time calculé & $\square$ \\
\hline
Toutes les tâches sont identifiées & $\square$ \\
\hline
Temps de cycle mesurés dans conditions réelles & $\square$ \\
\hline
Contraintes de précédence validées & $\square$ \\
\hline
Diagramme d'antériorité construit & $\square$ \\
\hline
Opérateurs consultés & $\square$ \\
\hline
Capacités machines connues & $\square$ \\
\hline
Temps de changement (SMED) intégrés & $\square$ \\
\hline
\end{tabular}
\caption{Check-list de préparation à l'équilibrage}
\end{table}

\subsection{Erreurs fréquentes à éviter}

\begin{methodbox}{Pièges à éviter lors de la collecte des données}
\begin{itemize}
    \item \textbf{Utiliser des temps standards trop optimistes} : les temps réels incluent les variations
    \item \textbf{Oublier les temps de déplacement} : entre les postes, pour la manutention
    \item \textbf{Négliger les temps de réglage} : changements de série, nettoyage
    \item \textbf{Ignorer les aléas} : micro-pannes, qualité, approvisionnement
    \item \textbf{Ne pas impliquer les opérateurs} : ils connaissent les vrais problèmes
    \item \textbf{Se baser sur des données anciennes} : les processus évoluent
    \item \textbf{Oublier les contraintes spécifiques} : partage d'équipements, compétences
\end{itemize}
\end{methodbox}

\begin{conseil}
\noindent\textbf{Bonne pratique :} Réalisez un VSM (Value Stream Mapping) avant l'équilibrage. Cette cartographie de la chaîne de valeur permet d'avoir une vision globale des flux, d'identifier les goulots et de collecter toutes les données nécessaires en une seule analyse.
\end{conseil}
\newpage
% ============================================================
% PARTIE 3: JUSTIFICATION DE LA NÉCESSITÉ D'UN ÉQUILIBRAGE
% ============================================================

\section{Justification de la nécessité d'un équilibrage}

Avant de lancer un projet d'équilibrage, il est essentiel de démontrer objectivement sa nécessité. Cette justification repose sur des indicateurs quantifiables et une analyse méthodique des dysfonctionnements observés.

\subsection{Indicateurs de déséquilibre}

Plusieurs signes visibles et mesurables indiquent qu'un équilibrage est nécessaire. La présence d'un ou plusieurs de ces indicateurs justifie une intervention.

\subsubsection{Stocks intermédiaires importants (WIP)}

L'accumulation de produits en cours (Work In Progress) entre les postes est le signe le plus visible d'un déséquilibre.

\begin{methodbox}{Observation des stocks WIP}
\begin{itemize}
    \item \textbf{Avant un poste} : indique que le poste est plus lent que le précédent (goulot)
    \item \textbf{Après un poste} : indique que le poste est plus rapide que le suivant (surcapacité)
    \item \textbf{Stocks importants et permanents} : déséquilibre structurel
    \item \textbf{Stocks fluctuants} : variabilité des cadences ou problèmes ponctuels
\end{itemize}
\end{methodbox}

\subsubsection{Temps d'attente observés}

Les temps d'attente des opérateurs ou des machines sont une perte directe de productivité.

\begin{itemize}
    \item \textbf{Opérateurs inactifs} : attente de pièces venant du poste précédent
    \item \textbf{Machines à l'arrêt} : en attente de chargement/déchargement
    \item \textbf{Taux d'occupation faible} : inférieur à 80\% sur certains postes
\end{itemize}

\subsubsection{Variabilité des cadences}

Lorsque les cadences de production sont très différentes d'un poste à l'autre :

\[
\text{Écart de cadence} = \frac{T_{Cmax} - T_{Cmin}}{T_{Cmax}} \times 100\%
\]

\begin{importantnote}
Un écart supérieur à 20\% entre le poste le plus rapide et le poste le plus lent indique un déséquilibre critique nécessitant une intervention.
\end{importantnote}

\subsubsection{Goulots d'étranglement identifiés}

Un goulot est un poste dont la cadence limite la production globale de la ligne.

\begin{methodbox}{Caractéristiques d'un goulot}
\begin{itemize}
    \item Le poste a le temps de cycle le plus élevé ($T_C > T_T$)
    \item Un stock important s'accumule en amont
    \item Une pénurie se crée en aval
    \item L'arrêt de ce poste arrête toute la ligne
    \item Son taux d'utilisation est proche de 100\%
\end{itemize}
\end{methodbox}

\subsubsection{Taux d'efficacité faible}

L'efficacité d'équilibrage est un indicateur synthétique qui mesure la qualité de la répartition des charges.

\[
\boxed{Eff = \frac{\sum T_i}{N \times T_{max}}}
\]

\begin{center}
\begin{tabular}{|c|c|}
\hline
\textbf{Efficacité} & \textbf{Niveau d'équilibrage} \\
\hline
$> 90\%$ & Excellent - Maintenir \\
\hline
$80\% - 90\%$ & Acceptable - Améliorations possibles \\
\hline
$70\% - 80\%$ & Médiocre - Équilibrage nécessaire \\
\hline
$< 70\%$ & Critique - Équilibrage urgent \\
\hline
\end{tabular}
\captionof{table}{Grille d'évaluation de l'efficacité d'équilibrage}
\end{center}

\subsubsection{Temps de cycle non alignés avec le Takt Time}

L'alignement des temps de cycle sur le Takt Time est l'objectif fondamental de l'équilibrage.

\[
\text{Écart relatif} = \frac{|T_C - T_T|}{T_T} \times 100\%
\]

\begin{importantnote}
Un écart supérieur à 10\% sur un ou plusieurs postes justifie un projet d'équilibrage. Si l'écart est supérieur à 20\% sur un poste, celui-ci constitue un goulot critique.
\end{importantnote}

\subsection{Calcul de l'efficacité d'équilibrage}

L'efficacité d'équilibrage peut être calculée de deux manières complémentaires.

\subsubsection{Efficacité théorique}

L'efficacité théorique compare le nombre théorique d'opérateurs au nombre réel utilisé.

\[
\boxed{Eff_{th} = \frac{N_{th}}{N_{r\acute{e}el}}}
\]

Avec :
\[
N_{th} = \frac{\sum T_i}{T_T}
\]

Où :
\begin{itemize}
    \item $N_{th}$ : Nombre théorique d'opérateurs (ou postes)
    \item $\sum T_i$ : Somme des temps de toutes les tâches
    \item $T_T$ : Takt Time
    \item $N_{r\acute{e}el}$ : Nombre réel d'opérateurs utilisé
\end{itemize}

\subsubsection{Efficacité réelle}

L'efficacité réelle mesure la performance effective de l'équilibrage réalisé.

\[
\boxed{Eff_{r\acute{e}} = \frac{\sum T_i}{N \times T_{max}}}
\]

Où :
\begin{itemize}
    \item $\sum T_i$ : Somme des temps de toutes les tâches
    \item $N$ : Nombre de postes
    \item $T_{max}$ : Temps du poste le plus lent (poste goulot)
\end{itemize}

\subsubsection{Exemple de calcul}

Prenons l'exemple d'une ligne avec les caractéristiques suivantes :

\begin{table}[H]
\centering
\begin{tabular}{|c|c|}
\hline
\textbf{Donnée} & \textbf{Valeur} \\
\hline
Somme des temps de tâches & 2,61 min \\
\hline
Takt Time & 0,7 min \\
\hline
Nombre réel de postes & 4 \\
\hline
Temps du poste le plus lent & 0,7 min \\
\hline
\end{tabular}
\end{table}

\begin{itemize}
    \item $N_{th} = \dfrac{2,61}{0,7} = 3,73$ opérateurs théoriques
    \item $Eff_{th} = \dfrac{3,73}{4} = 0,9325$ soit $93,25\%$
    \item $Eff_{r\acute{e}} = \dfrac{2,61}{4 \times 0,7} = \dfrac{2,61}{2,8} = 0,9325$ soit $93,25\%$
\end{itemize}

Dans cet exemple, l'efficacité théorique et réelle sont égales car le poste goulot est exactement au Takt Time.

\subsection{Efficacité par poste de travail}

L'analyse par poste permet d'identifier plus finement les déséquilibres et les goulots potentiels.

\[
\boxed{Eff_{poste} = \frac{T_{poste}}{T_{max}}}
\]

Ou plus couramment par rapport au Takt Time :

\[
\boxed{Eff_{poste} = \frac{T_{poste}}{T_T}}
\]

\begin{center}
\begin{tikzpicture}[scale=0.9]
    % Axes
    \draw[->] (0,0) -- (5,0) node[right] {Postes};
    \draw[->] (0,0) -- (0,5) node[above] {Temps (min)};
    
    % Barres
    \fill[blue!30] (0.5,0) rectangle (1.5,3.1);
    \node at (1,3.2) {0,62};
    \fill[blue!30] (2,0) rectangle (3,3.3);
    \node at (2.5,3.4) {0,66};
    \fill[blue!30] (3.5,0) rectangle (4.5,3.5);
    \node at (4,3.6) {0,70};
    \fill[blue!30] (5,0) rectangle (6,3.15);
    \node at (5.5,3.25) {0,63};
    
    % Takt Time
    \draw[dashed, red, thick] (0,3.5) -- (6.5,3.5) node[right] {Takt Time = 0,70 min};
    
    % Étiquettes
    \node at (1,-0.3) {Poste 1};
    \node at (2.5,-0.3) {Poste 2};
    \node at (4,-0.3) {Poste 3};
    \node at (5.5,-0.3) {Poste 4};
\end{tikzpicture}
\captionof{figure}{Diagramme d'équilibrage avec identification des écarts}
\end{center}

Les efficacités par poste sont alors :

\begin{table}[H]
\centering
\begin{tabular}{|c|c|c|c|}
\hline
\textbf{Poste} & \textbf{Temps (min)} & \textbf{Efficacité par poste} & \textbf{Écart au TT} \\
\hline
Poste 1 & 0,62 & 88,6\% & -0,08 min \\
\hline
Poste 2 & 0,66 & 94,3\% & -0,04 min \\
\hline
Poste 3 & 0,70 & 100\% & 0 min \\
\hline
Poste 4 & 0,63 & 90,0\% & -0,07 min \\
\hline
\end{tabular}
\caption{Analyse d'efficacité par poste}
\end{table}

\begin{importantnote}
Le poste 3 est le goulot de la ligne. C'est le poste le plus critique car le moindre aléa réduira sa production et impactera toute la ligne. Il doit être priorisé pour les actions d'amélioration (SMED, maintenance, soutien).
\end{importantnote}

\subsection{Tolérance et seuils d'intervention}

Pour définir une stratégie d'équilibrage, il est utile de disposer de seuils d'intervention clairs.

\begin{table}[H]
\centering
\begin{tabular}{|p{3.5cm}|p{2.5cm}|p{2.5cm}|p{3.5cm}|}
\hline
\textbf{Indicateur} & \textbf{Seuil d'alerte} & \textbf{Seuil critique} & \textbf{Action recommandée} \\
\hline
Efficacité globale & $< 85\%$ & $< 75\%$ & Équilibrage complet \\
\hline
Écart max au TT & $> 15\%$ & $> 25\%$ & Réaffectation tâches \\
\hline
Écart type des postes & $> 0,05$ min & $> 0,10$ min & Optimisation \\
\hline
WIP total & $> 1,5 \times \text{normal}$ & $> 3 \times \text{normal}$ & Rééquilibrage \\
\hline
Nombre de goulots & 1 & $\geq 2$ & Réorganisation \\
\hline
\end{tabular}
\caption{Seuils d'intervention pour l'équilibrage}
\end{table}

\subsection{Analyse des causes du déséquilibre}

Avant d'équilibrer, il est important de comprendre les causes profondes du déséquilibre.

\begin{methodbox}{Causes fréquentes de déséquilibre}
\begin{itemize}
    \item \textbf{Variabilité de la demande} : Changements non pris en compte
    \item \textbf{Évolution des produits} : Nouvelles opérations non intégrées
    \item \textbf{Différences de productivité} : Compétences hétérogènes
    \item \textbf{Temps de changement variables} : SMED non appliqué
    \item \textbf{Pannes récurrentes} : Maintenance insuffisante
    \item \textbf{Problèmes qualité} : Taux de rebut variable
    \item \textbf{Défaut de conception} : Séquences non optimisées
    \item \textbf{Absentéisme} : Variation des effectifs disponibles
\end{itemize}
\end{methodbox}

\subsection{Justification économique de l'équilibrage}

Pour convaincre de la nécessité d'un équilibrage, il est souvent utile de quantifier les gains potentiels.

\subsubsection{Coûts du déséquilibre}

\begin{table}[H]
\centering
\begin{tabular}{|p{4cm}|p{3cm}|p{5cm}|}
\hline
\textbf{Type de coût} & \textbf{Formule} & \textbf{Exemple} \\
\hline
Temps morts & $N_{postes} \times (T_T - T_{poste})$ & 4 postes $\times$ (0,7 - 0,65) = 0,2 min/cycle \\
\hline
Stocks WIP & $Q_{WIP} \times \text{Coût unitaire}$ & 500 pièces $\times$ 10 € = 5000 € \\
\hline
Surcapacité & $(N_{réel} - N_{th}) \times \text{Coût horaire}$ & 0,3 ETP $\times$ 35 k€ = 10,5 k€/an \\
\hline
Lead time & $\Delta LT \times \text{Coût de stockage}$ & +2 jours $\times$ 1000 €/jour = 2000 € \\
\hline
\end{tabular}
\caption{Évaluation des coûts liés au déséquilibre}
\end{table}

\subsubsection{Calcul du retour sur investissement}

\[
\boxed{ROI = \frac{\text{Gains annuels} - \text{Coût du projet}}{\text{Coût du projet}} \times 100\%}
\]

\begin{conseil}
Un projet d'équilibrage est généralement rentable si :
\begin{itemize}
    \item Le ROI est supérieur à 50\% sur la première année
    \item Le temps de retour sur investissement est inférieur à 18 mois
    \item Les gains en productivité dépassent 5\% de la masse salariale des opérateurs concernés
\end{itemize}
\end{conseil}

\subsection{Étude de cas : Diagnostic pré-équilibrage}

\textbf{Situation observée :}
Une ligne d'assemblage de 5 postes produit 400 pièces par jour sur 7h. Les temps mesurés sont :

\begin{table}[H]
\centering
\begin{tabular}{|c|c|}
\hline
\textbf{Poste} & \textbf{Temps de cycle (min)} \\
\hline
Poste 1 & 0,85 \\
\hline
Poste 2 & 0,72 \\
\hline
Poste 3 & 0,68 \\
\hline
Poste 4 & 0,91 \\
\hline
Poste 5 & 0,64 \\
\hline
\end{tabular}
\end{table}

\textbf{Diagnostic :}

\begin{enumerate}
    \item \textbf{Takt Time} : $T_T = \dfrac{7 \times 60}{400} = 1,05$ min/pièce
    \item \textbf{Temps total} : $\sum T_i = 0,85 + 0,72 + 0,68 + 0,91 + 0,64 = 3,8$ min
    \item \textbf{Nombre théorique} : $N_{th} = \dfrac{3,8}{1,05} = 3,62$ postes
    \item \textbf{Efficacité actuelle} : $Eff = \dfrac{3,8}{5 \times 0,91} = \dfrac{3,8}{4,55} = 83,5\%$
    \item \textbf{Goulot identifié} : Poste 4 avec $T_C = 0,91$ min (écart de -13\% par rapport au TT)
\end{enumerate}

\textbf{Conclusion :} L'efficacité est acceptable (83,5\%) mais un goulot est identifié au poste 4. Un rééquilibrage ciblé sur ce poste est recommandé.

\begin{importantnote}
Ce diagnostic montre que l'équilibrage n'est pas toujours une refonte complète. Une intervention ciblée sur le goulot peut suffire à améliorer significativement la performance de la ligne.
\end{importantnote}
\newpage
% ============================================================
% PARTIE 4: LES MÉTHODES D'ÉQUILIBRAGE
% ============================================================

\section{Les méthodes d'équilibrage}

L'équilibrage d'une ligne de production peut être réalisé selon différentes méthodes, allant des approches manuelles intuitives aux algorithmes numériques sophistiqués. Le choix de la méthode dépend de la complexité de la ligne, du nombre de tâches, et des contraintes de précédence.

\subsection{Classification des méthodes d'équilibrage}

\begin{center}
\begin{tikzpicture}[node distance=0.8cm, auto]
    \tikzset{
        methodnode/.style={rectangle, draw, minimum width=3cm, minimum height=0.6cm, align=center, rounded corners, fill=blue!5},
        category/.style={rectangle, draw, minimum width=4cm, minimum height=0.8cm, align=center, rounded corners, fill=blue!20, font=\bfseries}
    }
    
    \node[category] (manual) at (0,0) {Méthodes Manuelles};
    \node[methodnode] (intuitive) at (-3,-1.2) {Méthode Intuitive};
    \node[methodnode] (largest) at (0,-1.2) {Candidat le plus grand};
    \node[methodnode] (position) at (3,-1.2) {Poids Positionnel};
    
    \node[category] (num) at (0,-2.8) {Méthodes Numériques};
    \node[methodnode] (comsoal) at (-2.5,-4) {COMSOAL};
    \node[methodnode] (alpaca) at (0,-4) {ALPACA};
    \node[methodnode] (calb) at (2.5,-4) {CALB};
    \node[methodnode] (lp) at (5,-4) {Programmation Linéaire};
    
    \draw[->, thick] (manual) -- (intuitive);
    \draw[->, thick] (manual) -- (largest);
    \draw[->, thick] (manual) -- (position);
    \draw[->, thick] (num) -- (comsoal);
    \draw[->, thick] (num) -- (alpaca);
    \draw[->, thick] (num) -- (calb);
    \draw[->, thick] (num) -- (lp);
\end{tikzpicture}
\captionof{figure}{Classification des méthodes d'équilibrage}
\end{center}

\subsection{Méthodes manuelles}

Les méthodes manuelles sont particulièrement adaptées pour les lignes comportant un nombre limité de tâches (moins de 20) et pour les premières approches d'équilibrage.

\subsubsection{Méthode intuitive}

La méthode intuitive est la plus simple à mettre en œuvre. Elle repose sur le jugement et l'expérience de l'équilibreur.

\begin{methodbox}{Démarche de la méthode intuitive}
\begin{enumerate}
    \item Calculer le Takt Time et le nombre théorique de postes
    \item Regrouper les tâches en respectant les contraintes de précédence
    \item Remplir chaque poste jusqu'à approcher le Takt Time sans le dépasser
    \item Ajuster les regroupements pour améliorer l'équilibre
    \item Évaluer l'efficacité obtenue
    \item Itérer jusqu'à satisfaction
\end{enumerate}
\end{methodbox}

\textbf{Avantages :}
\begin{itemize}
    \item Rapide à mettre en œuvre
    \item Ne nécessite pas d'outils particuliers
    \item Favorise l'appropriation par les équipes terrain
\end{itemize}

\textbf{Inconvénients :}
\begin{itemize}
    \item Ne garantit pas la solution optimale
    \item Dépend fortement de l'expérience de l'équilibreur
    \item Difficile à formaliser
\end{itemize}

\subsubsection{Méthode du candidat le plus grand}

Cette méthode, également appelée "méthode du plus grand temps de cycle", assigne en priorité les tâches les plus longues.

\begin{methodbox}{Règle d'application}
\begin{enumerate}
    \item Lister toutes les tâches par ordre décroissant de temps de cycle
    \item Créer le premier poste de travail
    \item Prendre la première tâche de la liste (la plus longue) et l'assigner au poste si les contraintes de précédence sont respectées et si le temps total ne dépasse pas le Takt Time
    \item Continuer avec les tâches suivantes dans l'ordre de la liste
    \item Lorsqu'aucune tâche ne peut être ajoutée sans dépasser le Takt Time, créer un nouveau poste
    \item Répéter jusqu'à assignation de toutes les tâches
\end{enumerate}
\end{methodbox}

\textbf{Exemple d'application :}

Soit les tâches suivantes avec leurs temps de cycle :

\begin{table}[H]
\centering
\begin{tabular}{|c|c|c|}
\hline
\textbf{Tâche} & \textbf{Temps (min)} & \textbf{Prédécesseurs} \\
\hline
A & 0,4 & - \\
\hline
B & 0,6 & A \\
\hline
C & 0,5 & A \\
\hline
D & 0,3 & B, C \\
\hline
E & 0,4 & D \\
\hline
F & 0,2 & E \\
\hline
\end{tabular}
\caption{Données pour l'exemple}
\end{table}

Avec un Takt Time de 1,0 min.

\textbf{Liste ordonnée :} B(0,6), C(0,5), A(0,4), E(0,4), D(0,3), F(0,2)

\textbf{Assignation :}

\begin{enumerate}
    \item \textbf{Poste 1 :} B (0,6) → total 0,6 ; puis A (0,4) → total 1,0 (TT atteint)
    \item \textbf{Poste 2 :} C (0,5) → total 0,5 ; puis E (0,4) → total 0,9 ; F (0,2) → total 1,1 > TT, donc non
    \item \textbf{Poste 2 suite :} D (0,3) → total 0,8 ; F (0,2) → total 1,0 (TT atteint)
    \item \textbf{Poste 3 :} restant E (0,4) → total 0,4
\end{enumerate}

\textbf{Résultat :} 3 postes avec une efficacité de $\dfrac{2,4}{3 \times 1,0} = 80\%$

\subsubsection{Méthode du poids positionnel (Ranked Positional Weight)}

Cette méthode, développée par Helgeson et Birnie, attribue à chaque tâche un poids correspondant à la somme de son temps et de tous les temps des tâches qui lui succèdent.

\begin{methodbox}{Calcul du poids positionnel}
Pour une tâche $i$, le poids positionnel $PP_i$ est :
\[
PP_i = T_i + \sum_{j \in Successeurs(i)} T_j
\]
Où $Successeurs(i)$ est l'ensemble des tâches qui doivent être réalisées après la tâche $i$.
\end{methodbox}

\textbf{Démarche :}

\begin{enumerate}
    \item Calculer le poids positionnel de chaque tâche
    \item Classer les tâches par ordre décroissant de poids positionnel
    \item Assigner les tâches aux postes dans cet ordre, en respectant :
    \begin{itemize}
        \item Les contraintes de précédence
        \item La limite du Takt Time ($\sum T_{poste} \leq T_T$)
    \end{itemize}
\end{enumerate}

\textbf{Exemple :}

Reprenons les données précédentes avec le diagramme d'antériorité suivant :

\begin{center}
\begin{tikzpicture}[node distance=1.2cm, auto]
    \tikzset{
        task/.style={rectangle, draw, minimum width=1cm, minimum height=0.8cm, align=center, fill=blue!10}
    }
    \node[task] (A) at (0,0) {A\\0,4};
    \node[task] (B) at (-1,-1.2) {B\\0,6};
    \node[task] (C) at (1,-1.2) {C\\0,5};
    \node[task] (D) at (0,-2.4) {D\\0,3};
    \node[task] (E) at (0,-3.6) {E\\0,4};
    \node[task] (F) at (0,-4.8) {F\\0,2};
    
    \draw[->, thick] (A) -- (B);
    \draw[->, thick] (A) -- (C);
    \draw[->, thick] (B) -- (D);
    \draw[->, thick] (C) -- (D);
    \draw[->, thick] (D) -- (E);
    \draw[->, thick] (E) -- (F);
\end{tikzpicture}
\captionof{figure}{Diagramme d'antériorité pour l'exemple}
\end{center}

\textbf{Calcul des poids positionnels :}

\begin{itemize}
    \item Tâche F : $PP_F = 0,2$ (dernière tâche)
    \item Tâche E : $PP_E = 0,4 + 0,2 = 0,6$
    \item Tâche D : $PP_D = 0,3 + 0,4 + 0,2 = 0,9$
    \item Tâche B : $PP_B = 0,6 + 0,3 + 0,4 + 0,2 = 1,5$
    \item Tâche C : $PP_C = 0,5 + 0,3 + 0,4 + 0,2 = 1,4$
    \item Tâche A : $PP_A = 0,4 + 0,6 + 0,5 + 0,3 + 0,4 + 0,2 = 2,4$
\end{itemize}

\textbf{Liste ordonnée :} A(2,4), B(1,5), C(1,4), D(0,9), E(0,6), F(0,2)

\textbf{Assignation avec TT = 1,0 min :}

\begin{enumerate}
    \item \textbf{Poste 1 :} A (0,4) → total 0,4 ; B (0,6) → total 1,0 (TT atteint)
    \item \textbf{Poste 2 :} C (0,5) → total 0,5 ; D (0,3) → total 0,8 ; E (0,4) → total 1,2 > TT, donc non ; F (0,2) → total 1,0 (TT atteint)
    \item \textbf{Poste 3 :} E (0,4) → total 0,4
\end{enumerate}

\textbf{Résultat :} 3 postes avec la même répartition que la méthode précédente.

\subsection{Tableau comparatif des méthodes manuelles}

\begin{table}[H]
\centering
\begin{tabular}{|p{3cm}|p{3cm}|p{3cm}|p{4cm}|}
\hline
\textbf{Méthode} & \textbf{Avantages} & \textbf{Inconvénients} & \textbf{Utilisation recommandée} \\
\hline
Intuitive & Rapide, simple, participative & Subjective, non optimale & Petites lignes (< 10 tâches), pré-étude \\
\hline
Candidat le plus grand & Logique, facile à comprendre & Ignore les successeurs & Lignes simples, tâches indépendantes \\
\hline
Poids positionnel & Prend en compte les impacts & Calcul plus complexe & Lignes complexes avec précédences fortes \\
\hline
\end{tabular}
\caption{Comparaison des méthodes manuelles d'équilibrage}
\end{table}

\subsection{Méthodes numériques}

Pour les lignes complexes comportant de nombreuses tâches (plus de 20), les méthodes numériques assistées par ordinateur sont préférables.

\subsubsection{COMSOAL (Computer Method of Sequencing Operations for Assembly Lines)}

Développée par l'Université de Michigan, cette méthode génère aléatoirement de nombreuses solutions et conserve la meilleure.

\begin{methodbox}{Principe de COMSOAL}
\begin{enumerate}
    \item Lister toutes les tâches disponibles (sans prédécesseurs non assignés)
    \item Sélectionner aléatoirement une tâche parmi celles disponibles
    \item L'assigner au poste courant si le temps le permet
    \item Répéter jusqu'à épuisement des tâches ou saturation du poste
    \item Créer un nouveau poste et continuer
    \item Évaluer la solution (nombre de postes, efficacité)
    \item Générer un grand nombre de solutions (souvent 1000 à 10000)
    \item Retenir la meilleure solution
\end{enumerate}
\end{methodbox}

\textbf{Avantages :}
\begin{itemize}
    \item Explore un grand nombre de configurations
    \item Peut trouver des solutions proches de l'optimum
    \item Implémentation relativement simple
\end{itemize}

\subsubsection{ALPACA (Assembly Line Planning and Control Activity)}

ALPACA est une approche intégrée qui combine équilibrage et planification.

\textbf{Caractéristiques :}
\begin{itemize}
    \item Prend en compte les contraintes de production réelles
    \item Intègre la gestion des compétences des opérateurs
    \item Permet la simulation de scénarios
    \item Gère les variations de demande
\end{itemize}

\subsubsection{CALB (Computer-Aided Line Balancing)}

CALB est un outil spécialisé pour l'équilibrage assisté par ordinateur.

\textbf{Fonctionnalités :}
\begin{itemize}
    \item Interface graphique pour la saisie des tâches
    \item Visualisation du diagramme d'antériorité
    \item Algorithmes d'optimisation intégrés
    \item Génération de rapports d'équilibrage
    \item Simulation des flux
\end{itemize}

\subsubsection{Programmation linéaire}

L'équilibrage peut être formulé comme un problème d'optimisation mathématique.

\textbf{Formulation classique :}

Minimiser $N$ (nombre de postes)

Sous les contraintes :
\[
\sum_{k=1}^{N} x_{ik} = 1 \quad \forall i \text{ (chaque tâche assignée à un poste)}
\]
\[
\sum_{i} T_i \cdot x_{ik} \leq T_T \quad \forall k \text{ (limite de temps par poste)}
\]
\[
\sum_{k} k \cdot x_{ik} \leq \sum_{k} k \cdot x_{jk} \quad \forall (i,j) \text{ où i précède j}
\]
Avec $x_{ik} = 1$ si la tâche $i$ est assignée au poste $k$, 0 sinon.

\subsection{Méthodes hybrides et avancées}

\begin{itemize}
    \item \textbf{Algorithmes génétiques} : Inspiration de l'évolution biologique
    \item \textbf{Recuit simulé} : Exploration de l'espace des solutions avec acceptation de solutions dégradées
    \item \textbf{Colonies de fourmis} : Optimisation par comportement collectif
    \item \textbf{Réseaux de neurones} : Apprentissage sur des solutions historiques
\end{itemize}

\subsection{Choix de la méthode selon le contexte}

\begin{table}[H]
\centering
\begin{tabular}{|p{3cm}|p{3cm}|p{4cm}|p{3cm}|}
\hline
\textbf{Nombre de tâches} & \textbf{Complexité} & \textbf{Méthode recommandée} & \textbf{Temps d'étude} \\
\hline
$< 10$ & Faible & Intuitive ou Candidat le plus grand & 1-2 heures \\
\hline
$10 - 20$ & Moyenne & Poids positionnel & 1/2 journée \\
\hline
$20 - 50$ & Élevée & COMSOAL ou CALB & 1-2 jours \\
\hline
$> 50$ & Très élevée & Algorithmes avancés + logiciel & 3-5 jours \\
\hline
\end{tabular}
\caption{Guide de choix de la méthode d'équilibrage}
\end{table}

\begin{conseil}
\textbf{Bonnes pratiques :}
\begin{itemize}
    \item Commencer toujours par une méthode manuelle simple pour comprendre la structure du problème
    \item Valider les résultats obtenus avec les équipes terrain
    \item Utiliser un logiciel spécialisé pour les lignes complexes
    \item Documenter la solution retenue (diagrammes, tableaux d'assignation)
    \item Prévoir des marges de manœuvre pour les aléas (polyvalence, formation croisée)
\end{itemize}
\end{conseil}

\subsection{Étude de cas : Comparaison des méthodes}

Soit une ligne d'assemblage de 9 tâches avec les données suivantes :

\begin{table}[H]
\centering
\begin{tabular}{|c|c|c|}
\hline
\textbf{Tâche} & \textbf{Temps (min)} & \textbf{Prédécesseurs} \\
\hline
1 & 0,2 & - \\
\hline
2 & 0,4 & - \\
\hline
3 & 0,7 & 1, 2 \\
\hline
4 & 0,1 & 3 \\
\hline
5 & 0,3 & 3 \\
\hline
6 & 0,1 & 5 \\
\hline
7 & 0,3 & 5 \\
\hline
8 & 0,6 & 3, 4 \\
\hline
9 & 0,2 & 6, 7, 8 \\
\hline
\end{tabular}
\end{table}

Takt Time = 1,0 min

\textbf{Résultats comparatifs :}

\begin{table}[H]
\centering
\begin{tabular}{|p{4cm}|p{2.5cm}|p{2.5cm}|p{2.5cm}|}
\hline
\textbf{Méthode} & \textbf{Nombre de postes} & \textbf{Efficacité} & \textbf{Temps d'étude} \\
\hline
Intuitive & 4 & 85\% & 15 min \\
\hline
Candidat le plus grand & 4 & 85\% & 20 min \\
\hline
Poids positionnel & 4 & 85\% & 30 min \\
\hline
COMSOAL (1000 itérations) & 4 & 85\% & 2 min \\
\hline
Optimale théorique & 4 & 85\% & - \\
\hline
\end{tabular}
\caption{Comparaison des résultats selon la méthode}
\end{table}

\begin{importantnote}
Dans cet exemple, toutes les méthodes convergent vers la même solution. Pour des problèmes plus complexes, les écarts peuvent être significatifs (jusqu'à 10-15\% d'écart d'efficacité entre la méthode intuitive et l'optimum théorique).
\end{importantnote}
\newpage
% ============================================================
% PARTIE 5: LIEN ENTRE ÉQUILIBRAGE DES POSTES ET VSM (VALUE STREAM MAPPING)
% ============================================================

\section{Lien entre équilibrage des postes et VSM (Value Stream Mapping)}

Le Value Stream Mapping (VSM) ou cartographie de la chaîne de valeur est un outil fondamental du Lean Manufacturing qui permet d'identifier les opportunités d'amélioration, y compris pour l'équilibrage des postes de travail. L'intégration de ces deux approches permet une vision globale et systémique de la performance industrielle.

\subsection{Le VSM comme outil de diagnostic préalable}

Avant de réaliser un équilibrage, le VSM permet d'avoir une vision d'ensemble des flux et d'identifier les points critiques.

\subsubsection{Cartographie de l'état actuel}

Le VSM de l'état actuel représente le flux de valeur tel qu'il est réellement, mettant en évidence :

\begin{itemize}
    \item Les temps de cycle de chaque poste
    \item Les temps d'attente entre les opérations
    \item Les niveaux de stocks intermédiaires (WIP)
    \item Les temps de changement (SMED)
    \item Les taux de disponibilité des équipements
    \item Le lead time total et le value-added time
\end{itemize}

\begin{center}
\begin{tikzpicture}[scale=0.9, node distance=1cm]
    \tikzset{
        process/.style={rectangle, draw, minimum width=2cm, minimum height=1cm, align=center, fill=blue!10},
        data/.style={rectangle, draw, minimum width=1.5cm, minimum height=0.6cm, align=center, fill=yellow!10},
        arrow/.style={->, thick, -{Stealth[scale=1.2]}}
    }
    
    % Processus
    \node[process] (P1) at (0,0) {Poste 1\\$T_C=62s$\\$T_R=10min$\\Dispo=95\%};
    \node[process] (P2) at (4,0) {Poste 2\\$T_C=66s$\\$T_R=15min$\\Dispo=98\%};
    \node[process] (P3) at (8,0) {Poste 3\\$T_C=70s$\\$T_R=5min$\\Dispo=85\%};
    \node[process] (P4) at (12,0) {Poste 4\\$T_C=63s$\\$T_R=8min$\\Dispo=92\%};
    
    % Stocks
    \node[data] (WIP1) at (2,0.8) {WIP\\120 pcs};
    \node[data] (WIP2) at (6,0.8) {WIP\\85 pcs};
    \node[data] (WIP3) at (10,0.8) {WIP\\210 pcs};
    
    % Flèches
    \draw[arrow] (P1) -- (WIP1) -- (P2);
    \draw[arrow] (P2) -- (WIP2) -- (P3);
    \draw[arrow] (P3) -- (WIP3) -- (P4);
    
    % Temps
    \node[below] at (2,-0.5) {Attente 4s};
    \node[below] at (6,-0.5) {Attente 4s};
    \node[below] at (10,-0.5) {Attente 7s};
    
    % Ligne de temps
    \draw[decorate, decoration={brace, amplitude=8pt}] (0,-1.2) -- (12,-1.2) node[midway, below] {Lead Time = 45 min};
    \draw[decorate, decoration={brace, amplitude=8pt, mirror}] (0,-1.8) -- (12,-1.8) node[midway, below] {Value Added Time = 4,35 min};
\end{tikzpicture}
\captionof{figure}{Exemple de VSM simplifié avec identification des écarts}
\end{center}

\subsubsection{Indicateurs VSM pour l'équilibrage}

\begin{table}[H]
\centering
\begin{tabular}{|p{3cm}|p{3cm}|p{4cm}|p{3cm}|}
\hline
\textbf{Indicateur} & \textbf{Symbole} & \textbf{Signification} & \textbf{Lien avec équilibrage} \\
\hline
Cycle Time & $T_C$ & Temps de traitement d'un poste & Compare au Takt Time \\
\hline
Changeover Time & $T_R$ & Temps de changement de série & Impacte la charge disponible \\
\hline
Up Time & $U$ & Taux de disponibilité & Affecte le temps de cycle réel \\
\hline
Work In Progress & $WIP$ & Stocks entre postes & Indicateur de déséquilibre \\
\hline
Lead Time & $LT$ & Temps total de traversée & Impacté par l'équilibrage \\
\hline
Takt Time & $T_T$ & Rythme client & Objectif d'équilibrage \\
\hline
\end{tabular}
\caption{Indicateurs VSM clés pour l'équilibrage}
\end{table}

\subsection{Identification des goulots via le VSM}

Le VSM permet d'identifier visuellement les goulots d'étranglement qui limitent le flux de production.

\subsubsection{Le concept de goulot dans le VSM}

Un goulot est caractérisé par :

\begin{methodbox}{Signes d'un goulot dans le VSM}
\begin{itemize}
    \item $T_C > T_T$ : Le poste est plus lent que le rythme client
    \item $WIP$ important en amont : Accumulation devant le poste
    \item $WIP$ faible ou nul en aval : Pénurie après le poste
    \item $U \approx 100\%$ : Le poste fonctionne en continu
    \item $LT$ long : Le goulot allonge le lead time total
\end{itemize}
\end{methodbox}

\subsubsection{Analyse des flux matières}

Le VSM permet de visualiser le flux des matières et d'identifier les zones de stockage excessif.

\[
\text{Taux de rotation des stocks} = \frac{\text{Production journalière}}{\text{Stock moyen}}
\]

Un taux de rotation faible devant un poste indique un goulot potentiel.

\subsection{Synchronisation des flux avec le Takt Time}

L'un des objectifs principaux du VSM est d'aligner le flux sur le Takt Time.

\subsubsection{Calcul du Takt Time dans le VSM}

Le Takt Time est calculé à partir des données du VSM :

\[
\boxed{T_T = \frac{\text{Temps de production disponible par période}}{\text{Demande client par période}}}
\]

Dans un VSM complet, le Takt Time est affiché comme référence pour évaluer chaque poste.

\subsubsection{Alignement des temps de cycle}

Pour chaque poste, on calcule l'écart au Takt Time :

\[
\Delta_i = T_{C_i} - T_T
\]

\begin{center}
\begin{tikzpicture}
    \draw[->] (0,0) -- (10,0) node[right] {Postes};
    \draw[->] (0,0) -- (0,4) node[above] {Temps (s)};
    
    \draw[dashed, red, thick] (0,2.5) -- (10,2.5) node[right] {Takt Time = 70s};
    
    \fill[green!30] (0.5,0) rectangle (1.5,2.2);
    \fill[red!30] (2,0) rectangle (3,2.4);
    \fill[red!30] (3.5,0) rectangle (4.5,2.6);
    \fill[orange!30] (5,0) rectangle (6,2.9);
    \fill[green!30] (6.5,0) rectangle (7.5,2.1);
    
    \node at (1,2.5) {$\downarrow$};
    \node at (2.5,2.5) {$\downarrow$};
    \node at (4,2.5) {$\downarrow$};
    \node at (5.5,2.5) {$\uparrow$};
    \node at (7,2.5) {$\downarrow$};
    
    \node at (1,2.8) {-0,3};
    \node at (2.5,2.8) {-0,4};
    \node at (4,2.8) {-0,2};
    \node at (5.5,2.2) {+0,2};
    \node at (7,2.8) {-0,9};
\end{tikzpicture}
\captionof{figure}{Alignement des temps de cycle sur le Takt Time}
\end{center}

\begin{importantnote}
Un écart positif ($T_C > T_T$) indique un goulot qui nécessite une action prioritaire. Un écart négatif ($T_C < T_T$) indique une surcapacité qui peut être réaffectée.
\end{importantnote}

\subsection{Calcul du Ratio de Valeur Ajoutée}

Le Ratio de Valeur Ajoutée (Value Added Ratio) est un indicateur clé du VSM qui mesure l'efficacité globale du flux.

\[
\boxed{VAR = \frac{\text{Value Added Time}}{\text{Lead Time}} \times 100\%}
\]

Avant équilibrage :
\[
VAR_{avant} = \frac{4,35 \text{ min}}{45 \text{ min}} = 9,7\%
\]

Après équilibrage :
\[
VAR_{après} = \frac{4,35 \text{ min}}{12 \text{ min}} = 36,3\%
\]

\begin{conseil}
L'objectif de l'équilibrage est de réduire le Lead Time en diminuant les temps d'attente entre les postes, augmentant ainsi significativement le Ratio de Valeur Ajoutée.
\end{conseil}

\subsection{De l'état actuel à l'état futur}

Le VSM permet de définir un état futur cible après équilibrage.

\subsubsection{Étapes de transition}

\begin{enumerate}
    \item \textbf{Cartographier l'état actuel} : Relever tous les indicateurs
    \item \textbf{Identifier les opportunités} : Goulots, stocks excessifs, temps morts
    \item \textbf{Définir l'état futur} : Objectifs d'équilibrage et de flux
    \item \textbf{Réaliser l'équilibrage} : Appliquer les méthodes appropriées
    \item \textbf{Valider le nouvel état} : Mesurer et cartographier le flux amélioré
\end{enumerate}

\begin{center}
\begin{tikzpicture}[node distance=1.5cm, auto]
    \tikzset{
        state/.style={rectangle, draw, minimum width=3cm, minimum height=1.2cm, align=center, rounded corners, fill=blue!10}
    }
    
    \node[state] (current) at (0,0) {VSM\\État Actuel};
    \node[state] (analysis) at (3.5,0) {Analyse\\Goulots + WIP};
    \node[state] (balance) at (7,0) {Équilibrage\\des postes};
    \node[state] (future) at (10.5,0) {VSM\\État Futur};
    
    \draw[->, thick] (current) -- (analysis);
    \draw[->, thick] (analysis) -- (balance);
    \draw[->, thick] (balance) -- (future);
    
    \node[below] at (1.75,-0.5) {Identification};
    \node[below] at (5.25,-0.5) {Répartition};
    \node[below] at (8.75,-0.5) {Validation};
\end{tikzpicture}
\captionof{figure}{Intégration VSM-Équilibrage dans la démarche Lean}
\end{center}

\subsection{Exemple d'intégration VSM-Équilibrage}

\textbf{Cas concret :} Ligne d'assemblage de composants électroniques

\subsubsection{État actuel (VSM)}

\begin{table}[H]
\centering
\begin{tabular}{|c|c|c|c|c|}
\hline
\textbf{Poste} & \textbf{$T_C$ (s)} & \textbf{WIP (pcs)} & \textbf{$T_R$ (min)} & \textbf{$U$ (\%)} \\
\hline
Insertion composants & 45 & 250 & 20 & 92 \\
\hline
Soudure & 52 & 180 & 15 & 88 \\
\hline
Contrôle optique & 48 & 320 & 10 & 95 \\
\hline
Test fonctionnel & 68 & 120 & 5 & 85 \\
\hline
Emballage & 35 & 80 & 8 & 98 \\
\hline
\end{tabular}
\end{table}

\textbf{Calcul du Takt Time :} $T_T = \dfrac{7h \times 3600}{500 \text{ pièces}} = 50,4$ s/pièce

\textbf{Diagnostic VSM :}
\begin{itemize}
    \item Poste 4 (Test fonctionnel) : $T_C = 68s > T_T$ → Goulot
    \item WIP important devant poste 4 (320 pcs) → Accumulation
    \item Lead Time total : 45 min
    \item Value Added Time : 248 s = 4,13 min
    \item VAR = 9,2\%
\end{itemize}

\subsubsection{Équilibrage réalisé}

Application de la méthode du poids positionnel avec $T_T = 50,4$ s :

\begin{table}[H]
\centering
\begin{tabular}{|c|c|c|}
\hline
\textbf{Nouveau poste} & \textbf{Tâches regroupées} & \textbf{$T_C$ (s)} \\
\hline
Poste A & Insertion + Contrôle optique & 45 + 48 = 93 \\
\hline
Poste B & Soudure + Test partiel & 52 + 16 = 68 \\
\hline
Poste C & Test fonctionnel (dédoublé) & 68 / 2 = 34 \\
\hline
Poste D & Emballage + Test final & 35 + 12 = 47 \\
\hline
\end{tabular}
\end{table}

\subsubsection{État futur après équilibrage}

\begin{table}[H]
\centering
\begin{tabular}{|c|c|c|c|}
\hline
\textbf{Poste} & \textbf{$T_C$ (s)} & \textbf{WIP (pcs)} & \textbf{Observations} \\
\hline
Poste A & 45 & 50 & Poste rapide \\
\hline
Poste B & 52 & 45 & Proche du TT \\
\hline
Poste C & 34 & 30 & Poste dédoublé \\
\hline
Poste D & 47 & 40 & Poste rapide \\
\hline
\end{tabular}
\end{table}

\textbf{Résultats :}
\begin{itemize}
    \item Lead Time réduit de 45 min à 18 min (-60\%)
    \item VAR augmenté de 9,2\% à 23\% (+150\%)
    \item WIP total réduit de 950 à 165 pcs (-83\%)
    \item Efficacité d'équilibrage : 87\%
\end{itemize}

\subsection{Outils d'analyse complémentaires}

\subsubsection{Matrice de flux (Flow Matrix)}

La matrice de flux permet de visualiser les flux entre les postes avant et après équilibrage.

\begin{table}[H]
\centering
\begin{tabular}{|c|c|c|c|c|}
\hline
\textbf{De / Vers} & \textbf{Poste 1} & \textbf{Poste 2} & \textbf{Poste 3} & \textbf{Poste 4} \\
\hline
Poste 1 & - & 150 & 20 & 10 \\
\hline
Poste 2 & 5 & - & 180 & 15 \\
\hline
Poste 3 & 0 & 10 & - & 320 \\
\hline
Poste 4 & 0 & 0 & 5 & - \\
\hline
\end{tabular}
\caption{Matrice des flux avant équilibrage (flux journaliers)}
\end{table}

\subsubsection{Spaghetti Diagram}

Le diagramme spaghetti visualise les déplacements des opérateurs et des pièces, permettant d'identifier les gaspillages de mouvement qui peuvent être corrigés par l'équilibrage.

\begin{center}
\begin{tikzpicture}[scale=0.7]
    % Zones de travail
    \draw[fill=blue!10] (0,0) rectangle (2,2) node[midway] {Poste 1};
    \draw[fill=blue!10] (3,0) rectangle (5,2) node[midway] {Poste 2};
    \draw[fill=blue!10] (0,3) rectangle (2,5) node[midway] {Poste 3};
    \draw[fill=blue!10] (3,3) rectangle (5,5) node[midway] {Poste 4};
    
    % Déplacements (spaghetti)
    \draw[red, thick, loosely dotted] (1,1) -- (4,1) -- (1,4) -- (4,4) -- (1,1);
    \draw[red, thick, densely dotted] (1,1) -- (1,4) -- (4,4) -- (4,1) -- (1,1);
    
    \node at (2.5,-0.5) {Déplacements avant équilibrage : 156 m/jour};
    \node at (2.5,6) {Après équilibrage : 42 m/jour};
\end{tikzpicture}
\captionof{figure}{Diagramme spaghetti avant et après équilibrage}
\end{center}

\subsection{Conclusion sur l'intégration VSM-Équilibrage}

\begin{importantnote}
L'intégration du VSM et de l'équilibrage des postes permet une approche systémique de l'amélioration des flux :
\begin{enumerate}
    \item Le VSM fournit une vision globale et identifie les opportunités
    \item L'équilibrage apporte une solution structurée pour harmoniser les cadences
    \item La combinaison des deux outils maximise la réduction du lead time
    \item L'état futur VSM valide et pérennise les résultats de l'équilibrage
\end{enumerate}
\end{importantnote}

\begin{conseil}
Pour une démarche d'amélioration durable, il est recommandé de :
\begin{itemize}
    \item Réaliser un VSM complet avant tout projet d'équilibrage
    \item Impliquer les opérateurs dans la cartographie des flux
    \item Mettre à jour le VSM après chaque équilibrage majeur
    \item Utiliser les indicateurs VSM pour suivre la performance dans le temps
\end{itemize}
\end{conseil}
\newpage
% ============================================================
% PARTIE 6: MÉTHODOLOGIE ÉTAPE PAR ÉTAPE POUR L'ÉQUILIBRAGE
% ============================================================

\section{Méthodologie étape par étape pour l'équilibrage}

L'équilibrage d'une ligne de production suit une méthodologie structurée en quatre phases principales. Cette approche systématique garantit la rigueur de l'analyse et la pérennité des résultats.

\subsection{Vue d'ensemble de la méthodologie}

\begin{center}
\begin{tikzpicture}[node distance=1.2cm, auto]
    \tikzset{
        phase/.style={rectangle, draw, minimum width=3.5cm, minimum height=1cm, align=center, rounded corners, fill=blue!10, font=\bfseries},
        etape/.style={rectangle, draw, minimum width=3cm, minimum height=0.8cm, align=center, rounded corners, fill=green!5},
        arrow/.style={->, thick, -{Stealth[scale=1.2]}}
    }
    
    \node[phase] (P1) at (0,0) {Phase 1\\Analyse et collecte};
    \node[phase] (P2) at (4,0) {Phase 2\\Regroupement des tâches};
    \node[phase] (P3) at (8,0) {Phase 3\\Évaluation et ajustements};
    \node[phase] (P4) at (12,0) {Phase 4\\Implémentation et suivi};
    
    \node[etape] (E1) at (0,-1.5) {• Demande client\\• Temps disponibles\\• Mesures terrain};
    \node[etape] (E2) at (4,-1.5) {• Choix méthode\\• Assignation\\• Contraintes précédence};
    \node[etape] (E3) at (8,-1.5) {• Calcul efficacité\\• Identification écarts\\• Ajustements};
    \node[etape] (E4) at (12,-1.5) {• Formation\\• Standards\\• Suivi indicateurs};
    
    \draw[arrow] (P1) -- (P2);
    \draw[arrow] (P2) -- (P3);
    \draw[arrow] (P3) -- (P4);
    
    \node[below] at (2,-2.5) {Durée : 1-2 jours};
    \node[below] at (6,-2.5) {Durée : 1/2-1 jour};
    \node[below] at (10,-2.5) {Durée : 1 jour};
    \node[below] at (14,-2.5) {Durée : 1 semaine};
\end{tikzpicture}
\captionof{figure}{Vue d'ensemble de la méthodologie d'équilibrage}
\end{center}

\subsection{Phase 1 : Analyse et collecte des données}

Cette phase fondamentale consiste à rassembler toutes les informations nécessaires à l'équilibrage.

\subsubsection{Étape 1.1 : Identifier la demande client et calculer le Takt Time}

\begin{methodbox}{Calcul du Takt Time}
\[
\boxed{TT = \frac{\text{Temps de production disponible par période}}{\text{Demande client par période}}}
\]

\textbf{Exemple :}
\begin{itemize}
    \item Temps disponible : 7h30 = 450 minutes par jour
    \item Demande : 600 pièces par jour
    \item $TT = \dfrac{450}{600} = 0,75$ min/pièce = 45 secondes/pièce
\end{itemize}
\end{methodbox}

\textbf{Points de vigilance :}
\begin{itemize}
    \item Ne pas oublier de déduire les pauses et réunions
    \item Prendre en compte la maintenance préventive
    \item Considérer les variations saisonnières de la demande
\end{itemize}

\subsubsection{Étape 1.2 : Décomposer le processus et mesurer les temps de cycle}

\begin{methodbox}{Fiche de mesure des temps}
\begin{table}[H]
\centering
\begin{tabular}{|c|c|c|c|c|c|c|}
\hline
\textbf{Tâche} & \textbf{Mesure 1} & \textbf{Mesure 2} & \textbf{Mesure 3} & \textbf{Moyenne} & \textbf{Min} & \textbf{Max} \\
\hline
A & 0,58 & 0,62 & 0,60 & 0,60 & 0,58 & 0,62 \\
\hline
B & 0,65 & 0,68 & 0,67 & 0,67 & 0,65 & 0,68 \\
\hline
C & 0,70 & 0,72 & 0,71 & 0,71 & 0,70 & 0,72 \\
\hline
\end{tabular}
\end{table}
\end{methodbox}

\textbf{Bonnes pratiques de mesure :}
\begin{enumerate}
    \item Effectuer au moins 10 mesures par tâche
    \item Mesurer dans des conditions réelles de production
    \item Inclure les temps de manipulation et déplacement
    \item Identifier et exclure les mesures aberrantes
    \item Prendre en compte la variabilité normale
\end{enumerate}

\subsubsection{Étape 1.3 : Établir le diagramme d'antériorité}

\begin{center}
\begin{tikzpicture}[node distance=1.2cm, auto]
    \tikzset{
        task/.style={rectangle, draw, minimum width=1.2cm, minimum height=0.8cm, align=center, fill=blue!10}
    }
    
    % Nœuds
    \node[task] (1) at (0,0) {A\\0,6};
    \node[task] (2) at (2,0) {B\\0,4};
    \node[task] (3) at (1,-1.2) {C\\0,7};
    \node[task] (4) at (0,-2.4) {D\\0,5};
    \node[task] (5) at (2,-2.4) {E\\0,3};
    \node[task] (6) at (1,-3.6) {F\\0,2};
    
    % Flèches
    \draw[->, thick] (1) -- (3);
    \draw[->, thick] (2) -- (3);
    \draw[->, thick] (3) -- (4);
    \draw[->, thick] (3) -- (5);
    \draw[->, thick] (4) -- (6);
    \draw[->, thick] (5) -- (6);
    
    % Légende
    \node[below] at (1,-4.5) {Temps en minutes};
\end{tikzpicture}
\captionof{figure}{Exemple de diagramme d'antériorité complété}
\end{center}

\subsubsection{Étape 1.4 : Calculer le nombre théorique de postes}

\[
\boxed{N_{th} = \left\lceil \frac{\sum T_i}{TT} \right\rceil}
\]

Où $\lceil x \rceil$ représente l'arrondi à l'entier supérieur.

\textbf{Exemple :}
\begin{itemize}
    \item $\sum T_i = 2,7$ minutes
    \item $TT = 0,75$ minute
    \item $N_{th} = \lceil 2,7 / 0,75 \rceil = \lceil 3,6 \rceil = 4$ postes
\end{itemize}

\subsection{Phase 2 : Regroupement des tâches}

Cette phase consiste à assigner les tâches aux postes de travail en respectant les contraintes.

\subsubsection{Étape 2.1 : Choisir la méthode d'équilibrage}

Le choix de la méthode dépend de :
\begin{itemize}
    \item Nombre de tâches
    \item Complexité des contraintes de précédence
    \item Objectif de précision recherché
    \item Outils disponibles
\end{itemize}

\begin{table}[H]
\centering
\begin{tabular}{|p{3cm}|p{3cm}|p{3cm}|p{3cm}|}
\hline
\textbf{Nb tâches} & \textbf{Complexité} & \textbf{Méthode recommandée} & \textbf{Temps} \\
\hline
$< 10$ & Faible & Intuitive / Candidat plus grand & 1-2h \\
\hline
$10 - 20$ & Moyenne & Poids positionnel & 4-6h \\
\hline
$> 20$ & Élevée & Logiciel spécialisé & 1-2j \\
\hline
\end{tabular}
\caption{Guide de choix de la méthode}
\end{table}

\subsubsection{Étape 2.2 : Assigner les tâches aux postes}

\textbf{Règles d'assignation :}
\begin{enumerate}
    \item Respecter les contraintes de précédence
    \item Ne pas dépasser le Takt Time par poste
    \item Maximiser le taux de remplissage de chaque poste
    \item Considérer la polyvalence des opérateurs
    \item Tenir compte des contraintes d'outillage
\end{enumerate}

\begin{methodbox}{Tableau d'assignation des tâches}
\begin{table}[H]
\centering
\begin{tabular}{|c|c|c|c|c|}
\hline
\textbf{Poste} & \textbf{Tâches} & \textbf{Temps cumulé} & \textbf{Temps restant} & \textbf{Observations} \\
\hline
1 & A, B & 0,6 + 0,4 = 1,0 & 0,0 & TT atteint \\
\hline
2 & C, D & 0,7 + 0,5 = 1,2 & 0,2 & Dépasse TT \\
\hline
2 & C & 0,7 & 0,3 & À compléter \\
\hline
2 & C, E & 0,7 + 0,3 = 1,0 & 0,0 & TT atteint \\
\hline
3 & D & 0,5 & 0,5 & \\
\hline
3 & D, F & 0,5 + 0,2 = 0,7 & 0,3 & \\
\hline
\end{tabular}
\end{table}
\end{methodbox}

\subsubsection{Étape 2.3 : Documenter les regroupements}

Pour chaque poste, documenter :
\begin{itemize}
    \item La liste des tâches assignées
    \item Le temps de cycle total
    \item Les outils et équipements nécessaires
    \item Les compétences requises
    \item Les contraintes particulières
\end{itemize}

\subsection{Phase 3 : Évaluation et ajustements}

Cette phase permet de valider la qualité de l'équilibrage et d'optimiser la solution.

\subsubsection{Étape 3.1 : Calculer l'efficacité d'équilibrage}

\[
\boxed{Eff = \frac{\sum T_i}{N \times T_{max}} \times 100\%}
\]

\textbf{Exemple de calcul :}
\begin{table}[H]
\centering
\begin{tabular}{|c|c|c|c|}
\hline
\textbf{Poste} & \textbf{Temps (min)} & \textbf{Écart au TT} & \textbf{Efficacité poste} \\
\hline
1 & 0,68 & -0,02 & 97,1\% \\
\hline
2 & 0,70 & 0,00 & 100\% \\
\hline
3 & 0,65 & -0,05 & 92,9\% \\
\hline
4 & 0,67 & -0,03 & 95,7\% \\
\hline
\end{tabular}
\caption{Analyse par poste}
\end{table}

\[
Eff = \frac{0,68 + 0,70 + 0,65 + 0,67}{4 \times 0,70} = \frac{2,70}{2,80} = 96,4\%
\]

\subsubsection{Étape 3.2 : Identifier les écarts et les goulots}

\begin{center}
\begin{tikzpicture}
    \draw[->] (0,0) -- (6,0) node[right] {Postes};
    \draw[->] (0,0) -- (0,4) node[above] {Temps (min)};
    
    \draw[dashed, red, thick] (0,2.5) -- (6,2.5) node[right] {TT = 0,70 min};
    
    \fill[blue!30] (0.5,0) rectangle (1.5,2.38) node[midway, above] {0,68};
    \fill[blue!30] (2,0) rectangle (3,2.5) node[midway, above] {0,70};
    \fill[blue!30] (3.5,0) rectangle (4.5,2.28) node[midway, above] {0,65};
    \fill[blue!30] (5,0) rectangle (6,2.35) node[midway, above] {0,67};
    
    \node at (1,-0.3) {P1};
    \node at (2.5,-0.3) {P2};
    \node at (4,-0.3) {P3};
    \node at (5.5,-0.3) {P4};
\end{tikzpicture}
\captionof{figure}{Diagramme d'équilibrage avec identification du goulot}
\end{center}

\subsubsection{Étape 3.3 : Ajuster la répartition}

Si l'efficacité est insuffisante ($< 85\%$) :
\begin{enumerate}
    \item Identifier les postes les plus chargés et les moins chargés
    \item Transférer des tâches entre postes en respectant les précédences
    \item Dédoubler un poste goulot si nécessaire
    \item Regrouper des tâches de postes sous-chargés
    \item Recalculer l'efficacité après chaque ajustement
\end{enumerate}

\subsection{Phase 4 : Implémentation et suivi}

La dernière phase assure le déploiement opérationnel et la pérennisation des résultats.

\subsubsection{Étape 4.1 : Former les opérateurs}

\begin{methodbox}{Plan de formation type}
\begin{itemize}
    \item \textbf{Présentation} : Objectifs et bénéfices du nouvel équilibrage
    \item \textbf{Formation technique} : Nouvelles séquences de travail
    \item \textbf{Formation croisée} : Polyvalence sur plusieurs postes
    \item \textbf{Simulation} : Mise en situation sur poste
    \item \textbf{Accompagnement} : Coaching sur le terrain
    \item \textbf{Évaluation} : Validation des compétences
\end{itemize}
\end{methodbox}

\subsubsection{Étape 4.2 : Mettre en place les standards}

\begin{table}[H]
\centering
\begin{tabular}{|p{4cm}|p{4cm}|p{4cm}|}
\hline
\textbf{Document} & \textbf{Contenu} & \textbf{Utilisation} \\
\hline
Instruction de travail & Séquences, temps, points de contrôle & Formation, référence \\
\hline
Tableau de bord & Indicateurs clés & Suivi quotidien \\
\hline
Audit 5S & Organisation poste & Maintien standards \\
\hline
Fiche d'anomalie & Détection problèmes & Amélioration continue \\
\hline
\end{tabular}
\caption{Standards à mettre en place}
\end{table}

\subsubsection{Étape 4.3 : Suivre les indicateurs de performance}

\begin{table}[H]
\centering
\begin{tabular}{|c|c|c|c|c|}
\hline
\textbf{Indicateur} & \textbf{Cible} & \textbf{J-30} & \textbf{J+30} & \textbf{J+90} \\
\hline
Efficacité équilibrage & $> 90\%$ & 83\% & 89\% & 92\% \\
\hline
Taux de service & $> 98\%$ & 94\% & 96\% & 98\% \\
\hline
WIP (nb pièces) & $< 50$ & 120 & 75 & 45 \\
\hline
Lead Time (heures) & $< 4$ & 6,5 & 5,0 & 3,8 \\
\hline
Productivité (pcs/h) & $> 45$ & 38 & 42 & 46 \\
\hline
\end{tabular}
\caption{Tableau de bord de suivi post-équilibrage}
\end{table}

\subsubsection{Étape 4.4 : Animer les routines de terrain}

\begin{methodbox}{Routines d'amélioration continue}
\begin{itemize}
    \item \textbf{Quotidien} : Brief 5 min, suivi indicateurs
    \item \textbf{Hebdomadaire} : Audit de standards, résolution problèmes
    \item \textbf{Mensuel} : Revue des performances, plan d'actions
    \item \textbf{Trimestriel} : Rééquilibrage si besoin, formation
\end{itemize}
\end{methodbox}

\subsection{Check-list complète de la méthodologie}

\begin{table}[H]
\centering
\begin{tabular}{|p{10cm}|c|}
\hline
\textbf{PHASE 1 : ANALYSE ET COLLECTE} & \textbf{OK} \\
\hline
Demande client validée & $\square$ \\
\hline
Temps de production disponible calculé & $\square$ \\
\hline
Takt Time calculé & $\square$ \\
\hline
Toutes les tâches identifiées & $\square$ \\
\hline
Temps de cycle mesurés (min 10 mesures) & $\square$ \\
\hline
Contraintes de précédence listées & $\square$ \\
\hline
Diagramme d'antériorité construit & $\square$ \\
\hline
Temps total de traitement calculé & $\square$ \\
\hline
Nombre théorique de postes calculé & $\square$ \\
\hline
\textbf{PHASE 2 : REGROUPEMENT DES TÂCHES} & \textbf{OK} \\
\hline
Méthode d'équilibrage choisie & $\square$ \\
\hline
Assignation des tâches au poste 1 & $\square$ \\
\hline
Assignation des tâches au poste 2 & $\square$ \\
\hline
Assignation des tâches au poste 3 & $\square$ \\
\hline
Assignation des tâches au poste 4 & $\square$ \\
\hline
Respect des contraintes de précédence vérifié & $\square$ \\
\hline
Respect du Takt Time vérifié & $\square$ \\
\hline
Documentation des regroupements réalisée & $\square$ \\
\hline
\textbf{PHASE 3 : ÉVALUATION ET AJUSTEMENTS} & \textbf{OK} \\
\hline
Calcul de l'efficacité d'équilibrage & $\square$ \\
\hline
Identification des écarts & $\square$ \\
\hline
Identification du goulot & $\square$ \\
\hline
Ajustements effectués & $\square$ \\
\hline
Nouveau calcul d'efficacité & $\square$ \\
\hline
Validation avec l'équipe & $\square$ \\
\hline
\textbf{PHASE 4 : IMPLÉMENTATION ET SUIVI} & \textbf{OK} \\
\hline
Plan de formation établi & $\square$ \\
\hline
Opérateurs formés & $\square$ \\
\hline
Instructions de travail créées & $\square$ \\
\hline
Tableau de bord mis en place & $\square$ \\
\hline
Standards déployés & $\square$ \\
\hline
Routines de suivi définies & $\square$ \\
\hline
Première revue réalisée & $\square$ \\
\hline
\end{tabular}
\caption{Check-list complète de la méthodologie d'équilibrage}
\end{table}

\subsection{Étude de cas : Application complète de la méthodologie}

\textbf{Contexte :} Ligne d'assemblage de produits électroménagers

\subsubsection{Phase 1 - Analyse}
\begin{itemize}
    \item Demande : 400 pièces/jour
    \item Temps disponible : 7h = 420 min
    \item Takt Time : $TT = 420 / 400 = 1,05$ min/pièce
    \item $\sum T_i = 4,2$ min
    \item $N_{th} = \lceil 4,2 / 1,05 \rceil = 4$ postes
\end{itemize}

\subsubsection{Phase 2 - Regroupement}
\begin{table}[H]
\centering
\begin{tabular}{|c|c|c|}
\hline
\textbf{Poste} & \textbf{Tâches} & \textbf{Temps total (min)} \\
\hline
1 & A, B, C & 1,02 \\
\hline
2 & D, E & 1,04 \\
\hline
3 & F, G & 1,03 \\
\hline
4 & H, I & 1,01 \\
\hline
\end{tabular}
\end{table}

\subsubsection{Phase 3 - Évaluation}
\begin{itemize}
    \item $Eff = \dfrac{4,2}{4 \times 1,05} = 100\%$ (équilibre parfait)
    \item Lead Time réduit de 25 min à 4,2 min
    \item WIP réduit de 180 à 40 pièces
\end{itemize}

\subsubsection{Phase 4 - Implémentation}
\begin{itemize}
    \item Formation des 8 opérateurs réalisée
    \item Instructions de travail standardisées
    \item Tableau de bord mis en place
    \item Gains : +15\% de productivité
\end{itemize}

\begin{importantnote}
La méthodologie étape par étape garantit une approche rigoureuse et reproductible. L'application systématique de ces phases permet d'atteindre des résultats durables et d'instaurer une culture d'amélioration continue au sein des équipes.
\end{importantnote}
\newpage
% ============================================================
% PARTIE 7: OUTILS DE SUIVI ET JUSTIFICATION CONTINUE
% ============================================================

\section{Outils de suivi et justification continue}

L'équilibrage n'est pas une action ponctuelle mais un processus d'amélioration continue. Pour assurer la pérennité des résultats et justifier les investissements, il est essentiel de mettre en place des outils de suivi adaptés.

\subsection{Le diagramme d'équilibrage (Balance Chart)}

Le diagramme d'équilibrage est l'outil visuel fondamental pour représenter la répartition des charges entre les postes de travail.

\subsubsection{Construction du diagramme d'équilibrage}

\begin{methodbox}{Étapes de construction}
\begin{enumerate}
    \item Définir l'échelle verticale (temps) jusqu'au Takt Time + 20\%
    \item Représenter chaque poste par une barre horizontale
    \item Colorer la partie productive et la partie temps mort
    \item Ajouter le Takt Time comme ligne de référence
    \item Annoter les temps de cycle et les écarts
\end{enumerate}
\end{methodbox}

\begin{center}
\begin{tikzpicture}[scale=0.9]
    % Axes
    \draw[->, thick] (0,0) -- (10,0) node[right] {Postes de travail};
    \draw[->, thick] (0,0) -- (0,5) node[above] {Temps (minutes)};
    
    % Takt Time
    \draw[dashed, red, thick] (0,3.5) -- (10,3.5) node[right] {Takt Time = 0,70 min};
    
    % Barres des postes
    \fill[green!40] (0.5,0) rectangle (1.5,3.1);
    \fill[red!30] (0.5,3.1) rectangle (1.5,3.5);
    \node at (1,3.25) {\small 0,62};
    \node at (1,1.55) {P1};
    
    \fill[green!40] (2,0) rectangle (3,3.3);
    \fill[red!30] (2,3.3) rectangle (3,3.5);
    \node at (2.5,3.4) {\small 0,66};
    \node at (2.5,1.65) {P2};
    
    \fill[green!40] (3.5,0) rectangle (4.5,3.5);
    \node at (4,3.6) {\small 0,70};
    \node at (4,1.75) {P3};
    
    \fill[green!40] (5,0) rectangle (6,3.15);
    \fill[red!30] (5,3.15) rectangle (6,3.5);
    \node at (5.5,3.3) {\small 0,63};
    \node at (5.5,1.57) {P4};
    
    \fill[green!40] (6.5,0) rectangle (7.5,3.25);
    \fill[red!30] (6.5,3.25) rectangle (7.5,3.5);
    \node at (7,3.4) {\small 0,65};
    \node at (7,1.62) {P5};
    
    % Légende
    \fill[green!40] (8.5,4) rectangle (9,4.5);
    \fill[red!30] (9,4) rectangle (9.5,4.5);
    \node at (9.75,4.25) {Temps productif / Temps mort};
    
\end{tikzpicture}
\captionof{figure}{Diagramme d'équilibrage avec identification des temps morts}
\end{center}

\subsubsection{Interprétation du diagramme}

\begin{table}[H]
\centering
\begin{tabular}{|p{4cm}|p{8cm}|}
\hline
\textbf{Observation} & \textbf{Interprétation} \\
\hline
Barre atteignant le Takt Time & Poste goulot - critique à surveiller \\
\hline
Barre dépassant le Takt Time & Poste surchargé - nécessite dédoublement \\
\hline
Barre très inférieure au Takt Time & Sous-utilisation - potentiel de réallocation \\
\hline
Temps morts importants & Déséquilibre marqué - opportunité d'amélioration \\
\hline
Barres homogènes & Bon équilibrage - maintenir \\
\hline
\end{tabular}
\caption{Interprétation du diagramme d'équilibrage}
\end{table}

\subsection{Tableau de bord de suivi des performances}

Un tableau de bord permet de suivre l'évolution des indicateurs clés dans le temps.

\subsubsection{Indicateurs clés à suivre}

\begin{methodbox}{Indicateurs essentiels}
\begin{itemize}
    \item \textbf{Efficacité d'équilibrage} : Mesure globale de la qualité de l'équilibrage
    \item \textbf{Écart maximal au Takt Time} : $\max|T_{C_i} - TT|$
    \item \textbf{Écart-type des temps de cycle} : $\sigma = \sqrt{\frac{1}{N}\sum(T_{C_i} - \bar{T})^2}$
    \item \textbf{Niveau de WIP} : Stock entre les postes
    \item \textbf{Lead Time} : Temps total de traversée
    \item \textbf{Taux de service} : Capacité à satisfaire la demande
\end{itemize}
\end{methodbox}

\subsubsection{Modèle de tableau de bord}

\begin{table}[H]
\centering
\begin{tabular}{|c|c|c|c|c|c|c|}
\hline
\textbf{Indicateur} & \textbf{Cible} & \textbf{S-4} & \textbf{S-2} & \textbf{S0} & \textbf{S+2} & \textbf{S+4} \\
\hline
Efficacité équilibrage & $> 90\%$ & 78\% & 82\% & 85\% & 89\% & 92\% \\
\hline
Écart max au TT & $< 0,10$ min & 0,22 & 0,18 & 0,14 & 0,09 & 0,07 \\
\hline
Écart-type postes & $< 0,05$ min & 0,12 & 0,10 & 0,08 & 0,05 & 0,04 \\
\hline
WIP total (pcs) & $< 100$ & 280 & 220 & 160 & 110 & 80 \\
\hline
Lead Time (min) & $< 15$ & 32 & 26 & 20 & 16 & 13 \\
\hline
Taux de service & $> 98\%$ & 92\% & 94\% & 96\% & 97\% & 98,5\% \\
\hline
Productivité (pcs/h) & $> 45$ & 38 & 40 & 42 & 44 & 46 \\
\hline
\end{tabular}
\caption{Tableau de bord de suivi post-équilibrage}
\end{table}

\subsection{Calcul de la charge et de l'efficacité par poste}

L'analyse par poste permet d'identifier précisément les points d'attention.

\subsubsection{Efficacité par poste}

\[
\boxed{Eff_{poste} = \frac{T_{poste}}{T_{max}} \times 100\%}
\]

Ou par rapport au Takt Time :
\[
\boxed{Eff_{poste} = \frac{T_{poste}}{TT} \times 100\%}
\]

\begin{center}
\begin{tikzpicture}
    \begin{axis}[
        title={Efficacité par poste de travail},
        xlabel={Postes},
        ylabel={Efficacité (\%)},
        ymin=0, ymax=110,
        xtick={1,2,3,4,5},
        ytick={0,20,40,60,80,100},
        grid=both,
        nodes near coords,
        every node near coord/.append style={anchor=south}
    ]
    \addplot[ybar, fill=blue!30] coordinates {(1,88.6) (2,94.3) (3,100) (4,90.0) (5,92.9)};
    \addplot[red, thick, dashed] coordinates {(0,100) (6,100)};
    \end{axis}
\end{tikzpicture}
\captionof{figure}{Graphique d'efficacité par poste}
\end{center}

\subsubsection{Calcul de la charge globale}

\[
\boxed{\text{Charge journalière} = \frac{\text{Temps de cycle} \times \text{Production journalière}}{\text{Temps disponible}} \times 100\%}
\]

\begin{methodbox}{Exemple de calcul de charge}
\textbf{Poste 1 :}
\begin{itemize}
    \item $T_C = 0,62$ min/pièce
    \item Production = 600 pièces/jour
    \item Temps nécessaire = $0,62 \times 600 = 372$ min
    \item Temps disponible = 420 min (7h)
    \item Charge = $\dfrac{372}{420} = 88,6\%$
\end{itemize}
\end{methodbox}

\subsection{Justification économique de l'équilibrage}

Pour convaincre de la nécessité d'un équilibrage ou pour valider les résultats, une analyse économique est indispensable.

\subsubsection{Calcul des gains}

\begin{table}[H]
\centering
\begin{tabular}{|p{5cm}|p{4cm}|p{4cm}|}
\hline
\textbf{Poste de gain} & \textbf{Formule} & \textbf{Exemple} \\
\hline
Réduction temps morts & $\Delta TM \times \text{Coût horaire} \times \text{Nb jours}$ & 30 min/j $\times$ 25 €/h $\times$ 220 j = 2 750 € \\
\hline
Réduction WIP & $\Delta WIP \times \text{Coût unitaire} \times \text{Taux possession}$ & 200 pcs $\times$ 50 € $\times$ 0,25 = 2 500 € \\
\hline
Gain productivité & $\Delta Prod \times \text{Coût unitaire} \times \text{Production}$ & 5 pcs/h $\times$ 10 € $\times$ 2200 h = 110 000 € \\
\hline
Réduction Lead Time & $\Delta LT \times \text{Coût stockage}$ & 2 jours $\times$ 500 €/jour = 1 000 € \\
\hline
\end{tabular}
\caption{Calcul des gains potentiels}
\end{table}

\subsubsection{Calcul du retour sur investissement}

\[
\boxed{ROI = \frac{\text{Gains annuels} - \text{Coût du projet}}{\text{Coût du projet}} \times 100\%}
\]

\[
\boxed{\text{Délai de retour} = \frac{\text{Coût du projet}}{\text{Gains mensuels}}}
\]

\textbf{Exemple :}
\begin{itemize}
    \item Coût du projet (étude + formation) : 15 000 €
    \item Gains annuels : 45 000 €
    \item ROI = $\dfrac{45 000 - 15 000}{15 000} = 200\%$
    \item Délai de retour = $\dfrac{15 000}{45 000 / 12} = 4$ mois
\end{itemize}

\subsection{Outils de visualisation et de communication}

\subsubsection{Affiche de suivi d'équilibrage}

\begin{center}
\begin{tikzpicture}[scale=0.8]
    \draw[fill=gray!10] (0,0) rectangle (12,8);
    
    \node[fill=blue!20, draw, minimum width=10cm, minimum height=1cm] at (6,7) {TABLEAU DE SUIVI ÉQUILIBRAGE - LIGNE A};
    
    % Objectifs
    \draw[fill=yellow!20] (0.5,6) rectangle (5.5,5);
    \node at (3,5.5) {Objectifs : Efficacité > 90\% - WIP < 50 pcs};
    
    % Graphique
    \draw (6.5,6) rectangle (11.5,3);
    \draw[->] (6.8,5.8) -- (11.2,5.8);
    \draw[->] (6.8,5.8) -- (6.8,3.2);
    \node at (9,5.9) {Semaines};
    \node at (6.5,4.5) {Efficacité};
    
    % Résultats
    \draw (0.5,2.5) rectangle (11.5,1);
    \node at (6,1.75) {Dernier résultat : Efficacité = 92\% - Objectif atteint !};
    
    \node at (6,0.5) {Prochaine révision : Semaine 12};
\end{tikzpicture}
\captionof{figure}{Modèle d'affiche de suivi d'équilibrage}
\end{center}

\subsubsection{Rapport de performance mensuel}

\begin{methodbox}{Structure du rapport mensuel}
\begin{enumerate}
    \item \textbf{Résumé exécutif} : Synthèse des résultats clés
    \item \textbf{Évolution des indicateurs} : Graphiques comparatifs
    \item \textbf{Analyse par poste} : Diagramme d'équilibrage actualisé
    \item \textbf{Écarts et actions} : Ce qui n'a pas fonctionné et les correctifs
    \item \textbf{Plan d'actions} : Améliorations pour le mois suivant
    \item \textbf{Annexes} : Données détaillées, chronogrammes
\end{enumerate}
\end{methodbox}

\subsection{Routines d'amélioration continue}

\subsubsection{Revue quotidienne (5-10 minutes)}

\begin{table}[H]
\centering
\begin{tabular}{|p{3cm}|p{4cm}|p{5cm}|}
\hline
\textbf{Temps} & \textbf{Activité} & \textbf{Responsable} \\
\hline
5 min & Tour visuel des postes & Chef d'équipe \\
\hline
2 min & Relevé des indicateurs & Opérateur référent \\
\hline
3 min & Brief avec l'équipe & Chef d'équipe \\
\hline
\end{tabular}
\caption{Routine quotidienne de suivi}
\end{table}

\subsubsection{Revue hebdomadaire (30-60 minutes)}

\begin{methodbox}{Ordre du jour type}
\begin{enumerate}
    \item Revue des indicateurs de la semaine (5 min)
    \item Analyse des écarts et des goulots (10 min)
    \item Retour d'expérience des opérateurs (10 min)
    \item Validation des actions correctives (5 min)
    \item Planification de la semaine suivante (5 min)
    \item Points d'attention et alertes (5 min)
\end{enumerate}
\end{methodbox}

\subsubsection{Revue mensuelle (2 heures)}

\begin{table}[H]
\centering
\begin{tabular}{|p{4cm}|p{8cm}|}
\hline
\textbf{Thème} & \textbf{Contenu} \\
\hline
Performance & Analyse détaillée des indicateurs, tendances \\
\hline
Problèmes & Analyse des causes profondes des écarts \\
\hline
Améliorations & Bilan des actions menées, résultats obtenus \\
\hline
Objectifs & Définition des objectifs du mois suivant \\
\hline
Formation & Plan de développement des compétences \\
\hline
\end{tabular}
\caption{Routine mensuelle d'amélioration continue}
\end{table}

\subsection{Outils numériques de suivi}

\subsubsection{Tableau de bord Excel automatisé}

\begin{methodbox}{Fonctionnalités recommandées}
\begin{itemize}
    \item Saisie quotidienne des temps de cycle
    \item Calcul automatique de l'efficacité
    \item Graphiques dynamiques d'évolution
    \item Alertes conditionnelles (si efficacité < 85\%)
    \item Historique des données pour analyse tendancielle
    \item Export vers PowerBI pour visualisation avancée
\end{itemize}
\end{methodbox}

\subsubsection{Indicateurs avancés pour le suivi}

\begin{table}[H]
\centering
\begin{tabular}{|p{3.5cm}|p{3.5cm}|p{5cm}|}
\hline
\textbf{Indicateur} & \textbf{Formule} & \textbf{Interprétation} \\
\hline
Taux de synchronisation & $\frac{T_{max}}{TT}$ & Proximité du rythme client \\
\hline
Indice de déséquilibre & $\frac{\sigma}{TT}$ & Dispersion des temps \\
\hline
Efficacité opérationnelle & $\frac{\sum T_i}{N \times TT}$ & Performance globale \\
\hline
Taux d'occupation & $\frac{T_C}{TT}$ & Utilisation du poste \\
\hline
\end{tabular}
\caption{Indicateurs avancés pour le suivi}
\end{table}

\subsection{Justification de la nécessité d'un rééquilibrage}

\subsubsection{Seuils de déclenchement d'un rééquilibrage}

\begin{table}[H]
\centering
\begin{tabular}{|p{4cm}|p{3cm}|p{3cm}|p{3cm}|}
\hline
\textbf{Critère} & \textbf{Seuil alerte} & \textbf{Seuil critique} & \textbf{Action} \\
\hline
Efficacité & $< 85\%$ & $< 75\%$ & Rééquilibrage complet \\
\hline
Écart max TT & $> 15\%$ & $> 25\%$ & Analyse du goulot \\
\hline
WIP & $> 2\times$ normal & $> 3\times$ normal & Réaffectation \\
\hline
Lead Time & $> 1,5\times$ cible & $> 2\times$ cible & Réorganisation \\
\hline
Taux service & $< 95\%$ & $< 90\%$ & Action urgente \\
\hline
\end{tabular}
\caption{Seuils de déclenchement d'un rééquilibrage}
\end{table}

\begin{importantnote}
Un rééquilibrage doit être envisagé lors de :
\begin{itemize}
    \item Changement significatif de la demande ($\pm 20\%$)
    \item Introduction d'un nouveau produit
    \item Après une amélioration majeure (SMED, automatisation)
    \item Détection d'un nouveau goulot persistant
    \item Variation saisonnière prévisible
\end{itemize}
\end{importantnote}

\subsection{Étude de cas : Suivi sur 6 mois}

\textbf{Situation initiale :}
\begin{itemize}
    \item Efficacité : 78\%
    \item Lead Time : 32 min
    \item WIP : 280 pièces
    \item Taux service : 92\%
\end{itemize}

\textbf{Plan d'actions :}
\begin{enumerate}
    \item Mois 1 : Équilibrage et formation
    \item Mois 2 : Stabilisation et ajustements
    \item Mois 3 : Optimisation des temps de changement
    \item Mois 4-6 : Amélioration continue
\end{enumerate}

\textbf{Résultats après 6 mois :}

\begin{center}
\begin{tikzpicture}
    \begin{axis}[
        title={Évolution des indicateurs post-équilibrage},
        xlabel={Mois},
        ylabel={Valeur normalisée (base 100)},
        xmin=0, xmax=7,
        ymin=0, ymax=120,
        xtick={1,2,3,4,5,6},
        legend pos=north east,
        grid=both
    ]
    \addplot[blue, thick, mark=*] coordinates {(1,100) (2,105) (3,110) (4,115) (5,118) (6,120)};
    \addplot[red, thick, mark=square] coordinates {(1,100) (2,85) (3,75) (4,65) (5,58) (6,50)};
    \addplot[green, thick, mark=triangle] coordinates {(1,100) (2,90) (3,82) (4,75) (5,70) (6,68)};
    
    \legend{Productivité, WIP, Lead Time}
    \end{axis}
\end{tikzpicture}
\captionof{figure}{Évolution des indicateurs sur 6 mois}
\end{center}

\begin{conseil}
Pour assurer la pérennité des résultats :
\begin{itemize}
    \item Maintenir les routines de suivi quotidien
    \item Former les nouveaux opérateurs aux standards
    \item Réviser trimestriellement l'équilibrage
    \item Valoriser les suggestions d'amélioration des équipes
    \item Communiquer régulièrement sur les performances
\end{itemize}
\end{conseil}
\newpage
% ============================================================
% PARTIE 8: L'ÉQUILIBRAGE DANS LE CONTEXTE LEAN MANUFACTURING
% ============================================================

\section{L'équilibrage dans le contexte Lean Manufacturing}

L'équilibrage des postes de travail s'inscrit pleinement dans la philosophie Lean Manufacturing. Il ne s'agit pas seulement d'une technique d'optimisation mais d'un levier pour atteindre les objectifs fondamentaux du Lean : élimination des gaspillages, amélioration du flux, et création de valeur pour le client.

\subsection{Les principes Lean appliqués à l'équilibrage}

\subsubsection{Le Juste-à-Temps (JIT)}

Le Juste-à-Temps vise à produire ce qu'il faut, quand il faut, dans la quantité nécessaire. L'équilibrage est un outil essentiel pour atteindre cet objectif.

\begin{methodbox}{Application du JIT à l'équilibrage}
\begin{itemize}
    \item \textbf{Production au Takt Time} : Chaque poste produit au rythme de la demande client
    \item \textbf{Flux continu} : Élimination des stocks intermédiaires par synchronisation
    \item \textbf{Réduction des lots} : Équilibrage permettant des lots plus petits
    \item \textbf{Lead Time réduit} : Traversée plus rapide grâce à l'équilibrage
\end{itemize}
\end{methodbox}

\begin{center}
\begin{tikzpicture}[node distance=0.8cm]
    \tikzset{
        jitnode/.style={rectangle, draw, minimum width=2.5cm, minimum height=0.8cm, align=center, rounded corners, fill=blue!10}
    }
    
    \node[jitnode] (TT) at (0,0) {Takt Time};
    \node[jitnode] (balance) at (3,0) {Équilibrage};
    \node[jitnode] (flow) at (6,0) {Flux continu};
    \node[jitnode] (JIT) at (9,0) {Juste-à-Temps};
    
    \draw[->, thick] (TT) -- (balance);
    \draw[->, thick] (balance) -- (flow);
    \draw[->, thick] (flow) -- (JIT);
    
    \node[below] at (1.5,-0.5) {Cadence client};
    \node[below] at (4.5,-0.5) {Synchronisation};
    \node[below] at (7.5,-0.5) {Zéro stock};
\end{tikzpicture}
\captionof{figure}{Du Takt Time au Juste-à-Temps via l'équilibrage}
\end{center}

\subsubsection{Le Kaizen (Amélioration Continue)}

L'équilibrage n'est pas un événement ponctuel mais un processus d'amélioration continue.

\begin{importantnote}
Dans une approche Kaizen, l'équilibrage doit être régulièrement remis en question et amélioré :
\begin{itemize}
    \item Chaque semaine : Analyse des écarts et ajustements mineurs
    \item Chaque mois : Revue des performances et identification des opportunités
    \item Chaque trimestre : Rééquilibrage complet si nécessaire
\end{itemize}
\end{importantnote}

\subsubsection{Le Jidoka (Autonomation)}

Le Jidoka consiste à donner aux machines et aux opérateurs la capacité de détecter et d'arrêter la production en cas d'anomalie.

\begin{methodbox}{Lien entre Jidoka et équilibrage}
\begin{itemize}
    \item \textbf{Andon} : Signalement immédiat des déséquilibres (attente, surcharge)
    \item \textbf{Contrôle qualité intégré} : Équilibrage incluant les temps de contrôle
    \item \textbf{Polyvalence} : Opérateurs capables d'intervenir sur plusieurs postes
    \item \textbf{Standardisation} : Procédures claires pour maintenir l'équilibre
\end{itemize}
\end{methodbox}

\subsection{SMED et réduction des temps de changement}

Le SMED (Single Minute Exchange of Die) est une méthode Lean qui vise à réduire les temps de changement de série. Son application est cruciale pour un équilibrage efficace.

\subsubsection{Impact des temps de changement sur l'équilibrage}

\[
\boxed{T_{disponible} = T_{ouvré} - \sum T_{changement} - \sum T_{pannes}}
\]

\begin{table}[H]
\centering
\begin{tabular}{|p{4cm}|p{5cm}|p{4cm}|}
\hline
\textbf{Situation} & \textbf{Impact sur l'équilibrage} & \textbf{Solution SMED} \\
\hline
Changements longs & Capacité réduite, goulot & Réduction des temps internes/externes \\
\hline
Changements fréquents & Variabilité des cadences & Lisser la production (Heijunka) \\
\hline
Changements non standardisés & Écarts entre opérateurs & Standardisation des méthodes \\
\hline
\end{tabular}
\caption{Impact du SMED sur l'équilibrage}
\end{table}

\subsubsection{Intégration du SMED dans la démarche d'équilibrage}

\begin{enumerate}
    \item \textbf{Avant équilibrage} : Mesurer les temps de changement actuels
    \item \textbf{Pendant équilibrage} : Intégrer les temps de changement dans le calcul de charge
    \item \textbf{Après équilibrage} : Réduire les temps de changement pour libérer de la capacité
    \item \textbf{Continuellement} : Rééquilibrer après chaque réduction significative
\end{enumerate}

\begin{center}
\begin{tikzpicture}
    \draw[->, thick] (0,0) -- (8,0);
    \draw[fill=blue!20] (0,0) rectangle (1.5,0.5) node[midway] {Mesure};
    \draw[fill=blue!40] (1.5,0) rectangle (3.5,0.5) node[midway] {SMED};
    \draw[fill=blue!60] (3.5,0) rectangle (5.5,0.5) node[midway] {Gain capacité};
    \draw[fill=blue!80] (5.5,0) rectangle (7.5,0.5) node[midway] {Rééquilibrage};
    
    \node[below] at (0.75,-0.3) {Temps de changement};
    \node[below] at (2.5,-0.3) {Réduction};
    \node[below] at (4.5,-0.3) {Disponibilité};
    \node[below] at (6.5,-0.3) {Nouvel équilibre};
\end{tikzpicture}
\captionof{figure}{Cycle SMED - Équilibrage}
\end{center}

\subsection{Autonomation et standardisation}

\subsubsection{Standardisation des opérations}

La standardisation est le fondement de l'équilibrage durable.

\begin{methodbox}{Éléments à standardiser}
\begin{itemize}
    \item \textbf{Séquence des opérations} : Ordre précis des gestes
    \item \textbf{Temps de cycle cible} : Objectif par poste
    \item \textbf{Temps de prise de poste} : Rituel de démarrage
    \item \textbf{Instructions de travail} : Documents visuels
    \item \textbf{Contrôles qualité} : Points de vérification
\end{itemize}
\end{methodbox}

\subsubsection{Formation croisée et polyvalence}

La polyvalence des opérateurs est un facteur clé de réussite pour l'équilibrage.

\begin{table}[H]
\centering
\begin{tabular}{|p{3cm}|p{4cm}|p{5cm}|}
\hline
\textbf{Niveau} & \textbf{Compétences} & \textbf{Avantage pour l'équilibrage} \\
\hline
Niveau 1 & Poste unique & Formation initiale \\
\hline
Niveau 2 & 2-3 postes & Flexibilité, remplacement \\
\hline
Niveau 3 & Tous les postes & Rééquilibrage dynamique \\
\hline
Niveau 4 & Formation des autres & Pérennisation \\
\hline
\end{tabular}
\caption{Matrice de polyvalence}
\end{table}

\subsection{Équilibrage et gestion des goulots (Théorie des Contraintes)}

La Théorie des Contraintes (TOC) développée par Eliyahu Goldratt complète l'approche Lean pour la gestion des goulots.

\subsubsection{Les 5 étapes de la TOC appliquées à l'équilibrage}

\begin{enumerate}
    \item \textbf{Identifier le goulot} : Poste avec $T_C > TT$ ou le plus élevé
    \item \textbf{Exploiter le goulot} : Maximiser son utilisation (pas de temps mort)
    \item \textbf{Subordonner tout au goulot} : Synchroniser les autres postes sur son rythme
    \item \textbf{Élever le goulot} : Investir pour augmenter sa capacité
    \item \textbf{Répéter le cycle} : Le goulot se déplace, recommencer
\end{enumerate}

\begin{center}
\begin{tikzpicture}[node distance=1.2cm]
    \tikzset{
        tocnode/.style={rectangle, draw, minimum width=2cm, minimum height=0.8cm, align=center, rounded corners, fill=orange!10}
    }
    
    \node[tocnode] (E1) at (0,0) {1. Identifier};
    \node[tocnode] (E2) at (3,0) {2. Exploiter};
    \node[tocnode] (E3) at (6,0) {3. Subordonner};
    \node[tocnode] (E4) at (9,0) {4. Élever};
    \node[tocnode] (E5) at (12,0) {5. Répéter};
    
    \draw[->, thick] (E1) -- (E2);
    \draw[->, thick] (E2) -- (E3);
    \draw[->, thick] (E3) -- (E4);
    \draw[->, thick] (E4) -- (E5);
    \draw[->, thick, bend left] (E5) to (E1);
    
    \node[below] at (1.5,-0.5) {Quel est le goulot?};
    \node[below] at (4.5,-0.5) {Utiliser à 100\%};
    \node[below] at (7.5,-0.5) {Aligner les autres};
    \node[below] at (10.5,-0.5) {Augmenter capacité};
    \node[below] at (13.5,-0.5) {Nouveau goulot};
\end{tikzpicture}
\captionof{figure}{Les 5 étapes de la Théorie des Contraintes}
\end{center}

\subsubsection{Stock de découplage et théorie des contraintes}

Dans la TOC, un stock tampon est placé devant le goulot pour garantir son alimentation continue.

\begin{methodbox}{Gestion du stock de découplage}
\[
\boxed{Stock_{tampon} = T_C_{goulot} \times D_{max} \times (1 + \text{Marge sécurité})}
\]
\begin{itemize}
    \item Protège le goulot des aléas amont
    \item Permet de maintenir le rythme
    \item Doit être dimensionné au plus juste
    \item Réduit après stabilisation du flux
\end{itemize}
\end{methodbox}

\subsection{Heijunka (Lissage de la production)}

L'Heijunka est une technique Lean qui consiste à lisser le volume et le mix de production.

\subsubsection{Impact de l'Heijunka sur l'équilibrage}

\begin{table}[H]
\centering
\begin{tabular}{|p{4cm}|p{4cm}|p{4cm}|}
\hline
\textbf{Sans Heijunka} & \textbf{Avec Heijunka} & \textbf{Bénéfice} \\
\hline
Demande irrégulière & Demande lissée & Équilibrage stable \\
\hline
Batches importants & Petits lots & Moins de temps de changement \\
\hline
Goulots variables & Goulots identifiables & Actions ciblées \\
\hline
Surcharge ponctuelle & Charge constante & Moins de stress \\
\hline
\end{tabular}
\caption{Impact de l'Heijunka sur l'équilibrage}
\end{table}

\subsection{Équilibrage dynamique}

Dans un environnement Lean, l'équilibrage est dynamique et s'adapte aux variations.

\subsubsection{Facteurs d'adaptation}

\begin{methodbox}{Facteurs influençant l'équilibrage dynamique}
\begin{itemize}
    \item \textbf{Variations de la demande} : Changement de Takt Time
    \item \textbf{Améliorations terrain} : Réduction des temps de cycle
    \item \textbf{Variations de mix produit} : Séquences différentes
    \item \textbf{Absentéisme} : Variations d'effectifs disponibles
    \item \textbf{Maintenance} : Indisponibilités programmées
\end{itemize}
\end{methodbox}

\subsubsection{Modèle d'équilibrage adaptatif}

\begin{center}
\begin{tikzpicture}[node distance=0.8cm]
    \tikzset{
        dynnode/.style={rectangle, draw, minimum width=2cm, minimum height=0.6cm, align=center, fill=green!10}
    }
    
    \node[dynnode] (current) at (0,0) {État actuel};
    \node[dynnode] (measure) at (2.5,0) {Mesure};
    \node[dynnode] (compare) at (5,0) {Comparaison};
    \node[dynnode] (decision) at (7.5,0) {Décision};
    \node[dynnode] (action) at (10,0) {Action};
    \node[dynnode] (new) at (12.5,0) {Nouvel état};
    
    \draw[->, thick] (current) -- (measure);
    \draw[->, thick] (measure) -- (compare);
    \draw[->, thick] (compare) -- (decision);
    \draw[->, thick] (decision) -- (action);
    \draw[->, thick] (action) -- (new);
    \draw[->, thick, bend left] (new) to (current);
    
    \node[below] at (1.25,-0.4) {Temps de cycle};
    \node[below] at (3.75,-0.4) {vs Takt Time};
    \node[below] at (6.25,-0.4) {Écart > seuil?};
    \node[below] at (8.75,-0.4) {Réaffectation};
    \node[below] at (11.25,-0.4) {Équilibre modifié};
\end{tikzpicture}
\captionof{figure}{Boucle d'équilibrage dynamique}
\end{configuration}

\subsection{Indicateurs Lean pour l'équilibrage}

\subsubsection{Tableau de bord Lean spécifique}

\begin{table}[H]
\centering
\begin{tabular}{|p{3.5cm}|p{3cm}|p{3cm}|p{3cm}|}
\hline
\textbf{Indicateur Lean} & \textbf{Formule} & \textbf{Cible Lean} & \textbf{Lien équilibrage} \\
\hline
Taux de VA & $\frac{VA}{LT}$ & $> 30\%$ & Réduction attentes \\
\hline
Takt Time Achievement & $\frac{Production}{Demande}$ & $100\%$ & Synchronisation \\
\hline
Temps de cycle moyen & $\frac{\sum T_C}{N}$ & $= TT$ & Équilibre parfait \\
\hline
Taux de polyvalence & $\frac{\text{Postes maîtrisés}}{\text{Total postes}}$ & $> 70\%$ & Flexibilité \\
\hline
\end{tabular}
\caption{Indicateurs Lean pour le suivi de l'équilibrage}
\end{table}

\subsection{Étude de cas : Transformation Lean par l'équilibrage}

\textbf{Contexte :} Entreprise de fabrication industrielle, ligne de 8 postes

\subsubsection{Situation initiale}
\begin{itemize}
    \item Efficacité équilibrage : 68\%
    \item Lead Time : 45 minutes
    \item WIP : 380 pièces
    \item Taux de service : 89\%
    \item Nombre de goulots : 2
\end{itemize}

\subsubsection{Diagnostic Lean}
\begin{enumerate}
    \item Temps de changement longs (45 min en moyenne)
    \item Opérateurs non polyvalents
    \item Absence de standardisation
    \item Aucun suivi des indicateurs
\end{enumerate}

\subsubsection{Actions menées}
\begin{table}[H]
\centering
\begin{tabular}{|p{3cm}|p{4cm}|p{5cm}|}
\hline
\textbf{Outil Lean} & \textbf{Action} & \textbf{Résultat} \\
\hline
SMED & Réduction temps changement de 45 à 12 min & +15\% disponibilité \\
\hline
Standardisation & Instructions visuelles, routines & Écarts réduits de 40\% \\
\hline
Formation croisée & Polyvalence sur 3 postes par opérateur & Flexibilité accrue \\
\hline
Équilibrage & Répartition selon méthode poids positionnel & Efficacité 92\% \\
\hline
Tableau de bord & Suivi quotidien & Détection rapide anomalies \\
\hline
\end{tabular}
\end{table}

\subsubsection{Résultats après 6 mois}
\begin{itemize}
    \item Efficacité équilibrage : 68\% $\rightarrow$ 92\% (+24 pts)
    \item Lead Time : 45 min $\rightarrow$ 18 min (-60\%)
    \item WIP : 380 $\rightarrow$ 75 pièces (-80\%)
    \item Taux de service : 89\% $\rightarrow$ 98\% (+9 pts)
    \item Productivité : +22\%
\end{itemize}

\begin{importantnote}
\textbf{Enseignements clés :}
\begin{enumerate}
    \item L'équilibrage ne peut être dissocié des autres outils Lean
    \item L'implication des opérateurs est la clé du succès
    \item La standardisation pérennise les résultats
    \item Le suivi quotidien permet d'anticiper les dérives
    \item L'amélioration continue transforme l'équilibrage en culture
\end{enumerate}
\end{importantnote}

\subsection{Synthèse : L'équilibrage au cœur du Lean}

\begin{center}
\begin{tikzpicture}[node distance=1cm]
    \tikzset{
        leancenter/.style={circle, draw, minimum width=2cm, align=center, fill=red!20, font=\bfseries},
        leantool/.style={rectangle, draw, minimum width=2.2cm, minimum height=0.8cm, align=center, fill=blue!10}
    }
    
    \node[leancenter] (balance) at (0,0) {ÉQUILIBRAGE};
    
    \node[leantool] (SMED) at (-3,1.5) {SMED};
    \node[leantool] (JIT) at (3,1.5) {JIT};
    \node[leantool] (Kaizen) at (-3,-1.5) {Kaizen};
    \node[leantool] (Heijunka) at (3,-1.5) {Heijunka};
    \node[leantool] (Standard) at (0,2.5) {Standardisation};
    \node[leantool] (Jidoka) at (0,-2.5) {Jidoka};
    \node[leantool] (VSM) at (-4,0) {VSM};
    \node[leantool] (TOC) at (4,0) {Théorie des Contraintes};
    
    \draw[->, thick] (SMED) -- (balance);
    \draw[->, thick] (JIT) -- (balance);
    \draw[->, thick] (Kaizen) -- (balance);
    \draw[->, thick] (Heijunka) -- (balance);
    \draw[->, thick] (Standard) -- (balance);
    \draw[->, thick] (Jidoka) -- (balance);
    \draw[->, thick] (VSM) -- (balance);
    \draw[->, thick] (TOC) -- (balance);
    
    \node[below] at (0,-3.2) {L'équilibrage est le point de convergence des outils Lean pour le flux};
\end{tikzpicture}
\captionof{figure}{L'équilibrage au cœur du système Lean}
\end{center}

\begin{conseil}
Pour réussir votre démarche d'équilibrage dans un contexte Lean :
\begin{itemize}
    \item Commencez par stabiliser les processus (5S, standardisation)
    \item Réduisez les temps de changement avant d'équilibrer finement
    \item Formez les opérateurs à la polyvalence
    \item Mettez en place un tableau de bord visuel
    \item Animez des routines quotidiennes de suivi
    \item Valorisez les suggestions d'amélioration
    \item Révisez régulièrement l'équilibrage (cycle PDCA)
\end{itemize}
\end{conseil}
\newpage
% ============================================================
% PARTIE 9: CONCLUSION
% ============================================================

\section{Conclusion}

L'équilibrage des lignes de production est une pratique fondamentale du Lean Manufacturing qui permet d'optimiser les flux de production, de réduire les gaspillages et d'améliorer la performance globale de l'entreprise. Au terme de cette analyse approfondie, plusieurs enseignements essentiels se dégagent.

\subsection{Synthèse des concepts clés}

\subsubsection{Les fondements théoriques}

\begin{methodbox}{Points essentiels à retenir}
\begin{enumerate}
    \item \textbf{Takt Time} : Le rythme client est la référence absolue pour l'équilibrage
    \item \textbf{Cycle Time} : La mesure réelle des temps de production guide l'action
    \item \textbf{Lead Time} : La réduction du délai global est l'objectif ultime
    \item \textbf{Efficacité} : L'indicateur synthétique de la qualité de l'équilibrage
    \item \textbf{Goulot} : Le poste critique qui détermine la capacité de la ligne
\end{enumerate}
\end{methodbox}

\subsubsection{La méthodologie en 4 phases}

\begin{table}[H]
\centering
\begin{tabular}{|p{2.5cm}|p{3.5cm}|p{6cm}|}
\hline
\textbf{Phase} & \textbf{Objectif} & \textbf{Livrables clés} \\
\hline
Phase 1 & Analyse et collecte & Données temps, diagramme antériorité, Takt Time \\
\hline
Phase 2 & Regroupement & Répartition des tâches, documentation postes \\
\hline
Phase 3 & Évaluation & Efficacité calculée, ajustements validés \\
\hline
Phase 4 & Implémentation & Standards, formation, tableau de bord \\
\hline
\end{tabular}
\caption{Synthèse de la méthodologie d'équilibrage}
\end{table}

\subsection{Les bénéfices d'un équilibrage réussi}

\subsubsection{Bénéfices quantifiables}

\begin{center}
\begin{tikzpicture}
    \begin{axis}[
        title={Gains typiques après équilibrage},
        xlabel={Indicateurs},
        ylabel={Amélioration (\%)},
        x tick label style={rotate=45, anchor=east},
        xtick={1,2,3,4,5},
        xticklabels={Efficacité, Lead Time, WIP, Productivité, Taux service},
        ymin=0, ymax=100,
        grid=both,
        nodes near coords,
        every node near coord/.append style={anchor=south}
    ]
    \addplot[ybar, fill=blue!40] coordinates {(1,25) (2,60) (3,75) (4,20) (5,10)};
    \end{axis}
\end{tikzpicture}
\captionof{figure}{Gains moyens observés après équilibrage}
\end{center}

\subsubsection{Bénéfices qualitatifs}

\begin{itemize}
    \item \textbf{Amélioration des conditions de travail} : Réduction du stress, charge équilibrée
    \item \textbf{Meilleure visibilité} : Flux plus clair, problèmes identifiables
    \item \textbf{Engagement des équipes} : Implication dans l'amélioration continue
    \item \textbf{Flexibilité accrue} : Polyvalence, adaptation aux variations
    \item \textbf{Culture d'excellence} : Standardisation, rigueur, résultats
\end{itemize}

\subsection{Les conditions de succès}

\subsubsection{Facteurs critiques de réussite}

\begin{importantnote}
\textbf{Pour garantir le succès d'un projet d'équilibrage :}
\begin{enumerate}
    \item \textbf{Engagement de la direction} : Ressources, temps, soutien
    \item \textbf{Implication des opérateurs} : Expertise terrain, adhésion
    \item \textbf{Qualité des données} : Mesures fiables et représentatives
    \item \textbf{Formation adaptée} : Compétences techniques et comportementales
    \item \textbf{Suivi rigoureux} : Indicateurs, routines, ajustements
    \item \textbf{Patience et persévérance} : Résultats progressifs, amélioration continue
\end{enumerate}
\end{importantnote}

\subsubsection{Les pièges à éviter}

\begin{table}[H]
\centering
\begin{tabular}{|p{4cm}|p{4cm}|p{4cm}|}
\hline
\textbf{Piège} & \textbf{Conséquence} & \textbf{Solution} \\
\hline
Données erronées & Équilibrage inefficace & Mesures multiples, validation terrain \\
\hline
Oubli des aléas & Déséquilibre rapide & Intégrer marges, suivre quotidien \\
\hline
Non-implication équipes & Rejet, non-adhésion & Co-construction, communication \\
\hline
Équilibrage statique & Obsolescence rapide & Routines de révision, culture Kaizen \\
\hline
Absence de suivi & Retour à l'état initial & Tableau de bord, routines quotidiennes \\
\hline
\end{tabular}
\caption{Principaux pièges et solutions}
\end{table}

\subsection{L'équilibrage dans l'industrie 4.0}

\subsubsection{Nouvelles opportunités technologiques}

\begin{methodbox}{Technologies au service de l'équilibrage}
\begin{itemize}
    \item \textbf{Capteurs IoT} : Mesure en temps réel des temps de cycle
    \item \textbf{Digital Twin} : Simulation avant mise en œuvre
    \item \textbf{IA et Machine Learning} : Prédiction des déséquilibres
    \item \textbf{Tableaux de bord connectés} : Suivi temps réel, alertes
    \item \textbf{Réalité augmentée} : Formation, instructions visuelles
    \item \textbf{Exosquelettes} : Réduction de la pénibilité
\end{itemize}
\end{methodbox}

\subsubsection{Équilibrage prédictif}

L'intelligence artificielle permet désormais d'anticiper les déséquilibres avant qu'ils ne surviennent.

\[
\boxed{\text{Prédiction} = f(\text{Demande}, \text{Maintenance}, \text{Absentéisme}, \text{Qualité})}
\]

\subsection{L'équilibrage comme culture d'entreprise}

\subsubsection{De l'outil à la philosophie}

L'équilibrage ne doit pas être perçu comme une simple technique ponctuelle mais comme une composante de la culture d'entreprise.

\begin{center}
\begin{tikzpicture}[node distance=1.2cm]
    \tikzset{
        cultnode/.style={rectangle, draw, minimum width=2.5cm, minimum height=0.8cm, align=center, rounded corners, fill=green!10}
    }
    
    \node[cultnode] (outil) at (0,0) {Outil technique};
    \node[cultnode] (method) at (3,0) {Méthode structurée};
    \node[cultnode] (system) at (6,0) {Système de management};
    \node[cultnode] (culture) at (9,0) {Culture d'entreprise};
    
    \draw[->, thick] (outil) -- (method);
    \draw[->, thick] (method) -- (system);
    \draw[->, thick] (system) -- (culture);
    
    \node[below] at (1.5,-0.5) {Maturité 1};
    \node[below] at (4.5,-0.5) {Maturité 2};
    \node[below] at (7.5,-0.5) {Maturité 3};
\end{tikzpicture}
\captionof{figure}{L'évolution de l'équilibrage vers une culture d'entreprise}
\end{center}

\subsubsection{Les principes d'une culture d'équilibrage}

\begin{enumerate}
    \item \textbf{Visibilité} : Les performances sont affichées et partagées
    \item \textbf{Proximité} : Le management est présent sur le terrain (Gemba)
    \item \textbf{Responsabilisation} : Les équipes pilotent leur performance
    \item \textbf{Amélioration continue} : Chaque jour, on fait mieux qu'hier
    \item \textbf{Standardisation} : Les bonnes pratiques sont partagées
    \item \textbf{Polyvalence} : La flexibilité est développée
\end{enumerate}

\subsection{Feuille de route pour l'entreprise}

\subsubsection{Plan d'action type}

\begin{table}[H]
\centering
\begin{tabular}{|p{2cm}|p{4cm}|p{6cm}|}
\hline
\textbf{Période} & \textbf{Objectif} & \textbf{Actions clés} \\
\hline
Mois 1-2 & Diagnostic & Formation équipe, collecte données, VSM état actuel \\
\hline
Mois 2-3 & Conception & Calcul TT, regroupement tâches, simulation solutions \\
\hline
Mois 3-4 & Déploiement & Formation opérateurs, mise en place, ajustements \\
\hline
Mois 4-6 & Stabilisation & Suivi quotidien, résolution problèmes, standardisation \\
\hline
Mois 6-12 & Pérennisation & Audits, routines, amélioration continue \\
\hline
Annuel & Révision & Évaluation, rééquilibrage si besoin, nouveaux objectifs \\
\hline
\end{tabular}
\caption{Feuille de route type pour un projet d'équilibrage}
\end{table}

\subsection{Message final}

\begin{center}
\fbox{\begin{minipage}{0.85\textwidth}
\centering
\textbf{L'équilibrage des postes de travail : plus qu'une technique, une philosophie}

\vspace{0.3cm}

L'équilibrage n'est pas une fin en soi mais un moyen d'atteindre un flux continu, sans stocks, avec une production au rythme de la demande client. Il s'inscrit dans une démarche plus large d'excellence opérationnelle où chaque collaborateur est acteur de l'amélioration.

\vspace{0.3cm}
\textit{« Ce n'est pas l'équilibre qui crée la performance, mais la performance qui crée l'équilibre. »}
\end{minipage}}
\end{center}

\subsection{Annexes}

\subsubsection{Glossaire des termes clés}

\begin{table}[H]
\centering
\begin{tabular}{|p{3cm}|p{9cm}|}
\hline
\textbf{Terme} & \textbf{Définition} \\
\hline
Takt Time & Rythme de production dicté par la demande client \\
\hline
Cycle Time & Temps réel entre deux pièces produites sur un poste \\
\hline
Lead Time & Temps total entre l'entrée matière et la sortie produit fini \\
\hline
Goulot & Poste limitant la capacité de production de la ligne \\
\hline
WIP & Work In Progress - stocks de produits en cours \\
\hline
VSM & Value Stream Mapping - cartographie de la chaîne de valeur \\
\hline
SMED & Single Minute Exchange of Die - réduction des temps de changement \\
\hline
JIT & Juste-à-Temps - produire ce qu'il faut, quand il faut \\
\hline
Kaizen & Amélioration continue \\
\hline
Heijunka & Lissage de la production \\
\hline
\end{tabular}
\caption{Glossaire des termes clés}
\end{table}

\subsubsection{Formulaires et templates}

\begin{methodbox}{Documents utiles à télécharger}
\begin{itemize}
    \item Formulaire de collecte des temps de cycle
    \item Template de diagramme d'antériorité
    \item Tableau d'assignation des tâches
    \item Calculateur d'efficacité d'équilibrage
    \item Template de tableau de bord de suivi
    \item Guide de standardisation des postes
\end{itemize}
\end{methodbox}

\subsubsection{Références bibliographiques}

\begin{itemize}
    \item \textbf{Goldratt, E.M.} - "The Goal" (Théorie des Contraintes)
    \item \textbf{Ohno, T.} - "Toyota Production System" (Lean Manufacturing)
    \item \textbf{Womack, J. et Jones, D.} - "Lean Thinking"
    \item \textbf{Rother, M. et Shook, J.} - "Learning to See" (VSM)
    \item \textbf{Shingo, S.} - "A Revolution in Manufacturing" (SMED)
\end{itemize}

\subsubsection{Tableau récapitulatif des méthodes}

\begin{table}[H]
\centering
\begin{tabular}{|p{3cm}|p{3cm}|p{3cm}|p{3cm}|}
\hline
\textbf{Méthode} & \textbf{Complexité} & \textbf{Précision} & \textbf{Utilisation} \\
\hline
Intuitive & Faible & Moyenne & Petites lignes, pré-étude \\
\hline
Candidat plus grand & Faible & Bonne & Lignes simples \\
\hline
Poids positionnel & Moyenne & Très bonne & Lignes complexes \\
\hline
COMSOAL & Élevée & Optimale & Grandes lignes \\
\hline
Algorithmes avancés & Très élevée & Optimale & Recherche, optimisation \\
\hline
\end{tabular}
\caption{Synthèse comparative des méthodes d'équilibrage}
\end{table}

\subsubsection{Indicateurs de performance clés (KPI)}

\begin{table}[H]
\centering
\begin{tabular}{|p{4cm}|p{4cm}|p{4cm}|}
\hline
\textbf{KPI} & \textbf{Formule} & \textbf{Cible Lean} \\
\hline
Efficacité équilibrage & $\frac{\sum T_i}{N \times T_{max}}$ & $> 90\%$ \\
\hline
Taux de service & $\frac{\text{Livré à temps}}{\text{Commandé}}$ & $> 98\%$ \\
\hline
Lead Time & Mesure réelle & $< 1,5 \times \text{Valeur Ajoutée}$ \\
\hline
Taux de rotation WIP & $\frac{\text{Production}}{\text{WIP moyen}}$ & $> 10$ \\
\hline
Taux de polyvalence & $\frac{\text{Postes maîtrisés}}{\text{Total}}$ & $> 70\%$ \\
\hline
\end{tabular}
\caption{Indicateurs clés de performance pour l'équilibrage}
\end{table}

\subsection{Conclusion finale}

L'équilibrage des postes de travail est une discipline riche qui combine :

\begin{itemize}
    \item \textbf{Une approche scientifique} : Mesures, calculs, optimisation
    \item \textbf{Une méthode structurée} : Phases, outils, indicateurs
    \item \textbf{Une dimension humaine} : Implication, formation, polyvalence
    \item \textbf{Une philosophie Lean} : Élimination des gaspillages, flux continu
    \item \textbf{Un engagement durable} : Amélioration continue, culture d'excellence
\end{itemize}

\begin{center}
\fbox{\begin{minipage}{0.9\textwidth}
\centering
\textbf{Pour aller plus loin...}

\vspace{0.3cm}
L'équilibrage est une porte d'entrée vers l'excellence opérationnelle. Une fois maîtrisé, il ouvre la voie à des pratiques plus avancées :
\begin{itemize}
    \item \textbf{Automatisation intelligente} : Cibler les bons postes
    \item \textbf{Industrie 4.0} : Digitalisation du suivi
    \item \textbf{Lean Global} : Déploiement à l'échelle de l'entreprise
    \item \textbf{Excellence Supply Chain} : Intégration fournisseurs-clients
\end{itemize}
\end{minipage}}
\end{center}

\begin{importantnote}
\textbf{Dernier conseil :} Commencez petit, mais commencez maintenant. Un équilibrage imparfait sur une ligne pilote vaut mieux qu'une analyse parfaite qui n'aboutit jamais. Lancez-vous, mesurez, apprenez, et améliorez en continu. La perfection est un chemin, pas une destination.
\end{importantnote}
% ============================================================
% COMPLÉMENTS : APPROFONDISSEMENTS MÉTHODOLOGIQUES ET AVANCÉS
% ============================================================

\section{Approfondissements méthodologiques}

\subsection{Équilibrage multi-modèles}
La gestion de plusieurs références sur une même ligne nécessite de prendre en compte le mix produit et les variations de temps de cycle selon les modèles. Le \textit{Takt Time pondéré} est calculé comme suit :
\[
TT_{\text{pond}} = \frac{T_{\text{dispo}}}{\sum_{i=1}^{n} (d_i \times t_i)}
\]
où \(d_i\) est la demande du produit \(i\) et \(t_i\) son temps de cycle total. L'équilibrage peut alors être réalisé sur la base d'un modèle virtuel représentatif ou par séquencement (Heijunka).

\subsection{Équilibrage avec temps de changement variables}
L'intégration des temps de changement dans la charge effective est essentielle pour éviter les goulots. On définit la charge ajustée d'un poste comme :
\[
\text{Charge}_{\text{ajustée}} = \sum \text{Temps de cycle} + \sum \text{Temps de changement} \times \frac{\text{nombre de séries}}{T_{\text{dispo}}}
\]
Des stratégies de regroupement de séries permettent de réduire l'impact des réglages sur l'équilibre.

\subsection{Équilibrage stochastique}
Lorsque les temps de cycle sont sujets à une variabilité importante (micro-pannes, aléas humains), une approche stochastique est recommandée. Elle utilise des lois de probabilité (normale, log-normale) et la simulation Monte Carlo pour évaluer la probabilité de dépassement du Takt Time et dimensionner des marges de sécurité.

\subsection{Équilibrage à capacité finie}
Les contraintes de capacité (machines partagées, outils spécifiques, compétences rares) sont modélisées par des ressources supplémentaires. On parle alors d’équilibrage à capacité finie, souvent résolu par des algorithmes d’optimisation intégrant des contraintes de ressources.

% ------------------------------------------------------------
\section{Outils numériques et industrie 4.0}

\subsection{Mesure automatique des temps de cycle}
Les systèmes MES (Manufacturing Execution System) et les capteurs IoT permettent d’acquérir en continu les temps de cycle réels. Ils alimentent des tableaux de bord dynamiques et déclenchent des alertes en cas de dérive par rapport au Takt Time.

\subsection{Jumeau numérique (Digital Twin)}
Avant toute modification de ligne, un jumeau numérique permet de simuler l’impact d’un nouvel équilibrage. Les scénarios de variation de demande, d’absentéisme ou de pannes peuvent être testés sans risque.

\subsection{Algorithmes d’optimisation avancés}
Pour des lignes complexes (plus de 50 tâches), des métaheuristiques comme les algorithmes génétiques, le recuit simulé ou les colonies de fourmis permettent d’approcher la solution optimale. Un tableau comparatif des performances est présenté ci-dessous.

\begin{table}[H]
\centering
\caption{Performances des algorithmes d’optimisation}
\label{tab:algo_perf}
\begin{tabular}{|l|c|c|}
\hline
\textbf{Algorithme} & \textbf{Efficacité obtenue} & \textbf{Temps de calcul (secondes)} \\ \hline
Algorithme génétique & 94,2\% & 12 \\
Recuit simulé & 93,8\% & 8 \\
Colonies de fourmis & 94,5\% & 15 \\
\hline
\end{tabular}
\end{table}

\subsection{Intégration avec ERP et MES}
L’équilibrage devient dynamique lorsqu’il est couplé au système d’information. Les données de production (temps réels, ordres de fabrication) alimentent automatiquement le modèle, et des rééquilibrages partiels peuvent être proposés en temps réel.

% ------------------------------------------------------------
\section{Dimensions humaines et organisationnelles}

\subsection{Conduite du changement appliquée à l’équilibrage}
Un projet d’équilibrage implique souvent une modification des tâches, des périmètres et des habitudes. Une cartographie des parties prenantes et une communication proactive permettent d’anticiper les résistances. La valorisation des réussites et l’implication des opérateurs dès la phase de diagnostic sont des facteurs clés.

\subsection{Polyvalence et formation}
La matrice de polyvalence est un outil central. Elle identifie les compétences de chaque opérateur et les besoins de formation croisée.

\begin{table}[H]
\centering
\caption{Exemple de matrice de polyvalence}
\label{tab:polyvalence}
\begin{tabular}{|l|c|c|c|c|}
\hline
\textbf{Opérateur} & \textbf{Poste 1} & \textbf{Poste 2} & \textbf{Poste 3} & \textbf{Poste 4} \\ \hline
A & ✓ & ✓ & & \\
B & ✓ & & ✓ & ✓ \\
C & & ✓ & ✓ & \\
\hline
\end{tabular}
\end{table}

\subsection{Ergonomie et charge de travail}
L’équilibrage ne doit pas se faire au détriment de la santé des opérateurs. L’évaluation des risques ergonomiques (manutention, gestes répétitifs, postures) doit être intégrée dès le regroupement des tâches. Des outils comme le checklist OSHA ou l’analyse RULA peuvent être utilisés.

\subsection{Management visuel et routines}
Un tableau de bord d’équipe affichant les temps de cycle, les écarts au Takt Time et les actions en cours favorise la réactivité. Les routines quotidiennes (brief 5 min, tour de terrain) permettent de maintenir l’équilibre dans le temps.

% ------------------------------------------------------------
\section{Modélisation économique et retour sur investissement}

\subsection{Calcul détaillé des coûts du déséquilibre}
Les coûts cachés d’un déséquilibre sont souvent sous-estimés. Une approche analytique les quantifie :

\begin{itemize}
    \item \textbf{Temps morts} : \( \text{Coût} = \sum (\text{Temps mort par poste}) \times \text{Taux horaire chargé} \)
    \item \textbf{Stocks WIP} : \( \text{Coût} = \text{Niveau WIP} \times \text{Coût unitaire} \times \text{Taux de possession} \)
    \item \textbf{Lead time} : \( \text{Coût} = \Delta LT \times \text{Coût journalier de stockage} + \text{pénalités clients} \)
\end{itemize}

\subsection{Évaluation des gains d’un rééquilibrage}
Les gains potentiels incluent la productivité libérée (réaffectation d’opérateurs), l’augmentation de la capacité de production et la réduction des aléas. Une fiche de calcul systématique permet de présenter le business case.

\subsection{Modèle financier complet}
Le retour sur investissement (ROI) et la valeur actuelle nette (VAN) sont calculés sur la base des gains annuels prévus et des coûts du projet (études, formation, éventuels investissements légers). Un seuil de délai de récupération inférieur à 12 mois est généralement exigé.

% ------------------------------------------------------------
\section{Cas d’application par secteur}

\subsection{Industrie automobile}
Les chaînes d’assemblage automobile sont caractérisées par des cadences élevées (Takt Time souvent inférieur à 1 minute). L’équilibrage doit intégrer les opérateurs itinérants, les contrôles qualité intégrés et les aléas de supply chain.

\subsection{Électronique et high-mix low-volume}
Dans ce secteur, la diversité des produits impose des équilibrages fréquents. Les cellules flexibles et les opérateurs polyvalents sont privilégiés. L’approche du \textit{line balancing} est souvent couplée à la conception modulaire.

\subsection{Agroalimentaire et process continu}
Les contraintes sanitaires et les temps de nettoyage (Nettoyage en Cours, NEC) complexifient l’équilibrage. On modélise ces temps comme des tâches obligatoires avec des contraintes de précédence spécifiques.

\subsection{Logistique et préparation de commandes}
L’équilibrage s’applique également aux zones de picking. On répartit les lignes de commande entre opérateurs en fonction des distances de déplacement et des temps de prélèvement, avec un objectif de synchronisation avec l’expédition.

% ============================================================
% FIN DES COMPLÉMENTS
% ============================================================
% ============================================================
% COMPLÉMENTS (suite) : NORMES, DÉVELOPPEMENT DURABLE, PILOTAGE, ANNEXES
% ============================================================

\section{Intégration avec les normes et référentiels}

\subsection{Normes ISO (9001, 50001, 45001)}
L'équilibrage contribue directement à la maîtrise des processus exigée par l'ISO 9001. En stabilisant les cadences, il réduit les déchets et favorise une meilleure efficacité énergétique (ISO 50001) et des conditions de travail plus sûres (ISO 45001).

\begin{table}[H]
\centering
\caption{Lien entre équilibrage et exigences ISO}
\label{tab:iso}
\begin{tabular}{|l|p{8cm}|}
\hline
\textbf{Norme} & \textbf{Contribution de l'équilibrage} \\ \hline
ISO 9001 & Améliore la prévisibilité des processus, réduit les dérives qualité liées aux fluctuations de cadence. \\
ISO 50001 & Permet de lisser les pics de consommation énergétique en évitant les surcapacités et les attentes. \\
ISO 45001 & Diminue les risques de troubles musculo-squelettiques par une répartition équitable des charges. \\
\hline
\end{tabular}
\end{table}

\subsection{Référentiels Lean et World Class Manufacturing (WCM)}
Dans les référentiels tels que le Lean Six Sigma ou le WCM, l’équilibrage est un outil central des piliers « Coûts », « Personnes » et « Logistique ». Il est souvent évalué par des audits spécifiques.

\subsection{Certification des compétences}
La mise en place d’un référentiel interne de certification (ex. : « opérateur qualifié équilibrage ») permet de formaliser les compétences et de garantir la pérennité des standards. Des organismes comme l’AFNOR proposent également des certifications adaptées.

% ------------------------------------------------------------
\section{Équilibrage et développement durable}

\subsection{Efficacité énergétique}
Un équilibrage optimal réduit les temps morts et les pics de charge, ce qui se traduit par une baisse de la consommation électrique des équipements (machines en attente, compresseurs, éclairage). On estime une réduction de 5 à 15\% de la consommation spécifique après un rééquilibrage.

\subsection{Réduction des déchets}
En stabilisant les flux, l’équilibrage diminue le risque d’obsolescence des stocks intermédiaires et les rebuts liés à des cadences irrégulières. Il favorise également la standardisation, source de qualité.

\subsection{Économie circulaire}
Les ressources libérées (opérateurs, surfaces, machines) peuvent être réaffectées à des activités à plus forte valeur ajoutée ou à la production de nouvelles références sans investissements supplémentaires.

% ------------------------------------------------------------
\section{Feuille de route de déploiement et pilotage projet}

\subsection{Création d’une équipe projet}
Une structure légère mais efficace est recommandée :
\begin{itemize}
    \item \textbf{Sponsor} : direction, garant des ressources et du cadrage stratégique.
    \item \textbf{Pilote} : responsable du déroulement, de la méthode et des résultats.
    \item \textbf{Référents terrain} : opérateurs, chefs d’équipe, maintenance.
\end{itemize}

\subsection{Calendrier type}
\begin{table}[H]
\centering
\caption{Planning indicatif d’un projet d’équilibrage}
\label{tab:planning}
\begin{tabular}{|l|l|p{6cm}|}
\hline
\textbf{Phase} & \textbf{Durée} & \textbf{Livrables} \\ \hline
Diagnostic & 2 semaines & VSM état actuel, temps de cycle, diagramme d’antériorité \\
Conception & 1 semaine & Solutions d’équilibrage, simulation, business case \\
Test & 1 semaine & Validation sur ligne pilote, ajustements \\
Déploiement & 2 semaines & Formation, mise en place des standards \\
Stabilisation & 4 semaines & Suivi quotidien, résolution des anomalies \\
\hline
\end{tabular}
\end{table}

\subsection{Indicateurs de pilotage projet}
Les indicateurs suivants permettent de suivre l’avancement et la performance :
\begin{itemize}
    \item \textbf{Jalons} : respect des délais de chaque phase.
    \item \textbf{Efficacité d’équilibrage} : avant / après.
    \item \textbf{Taux d’adhésion} : participation des opérateurs aux formations.
    \item \textbf{ROI} : suivi des gains réels par rapport aux prévisions.
\end{itemize}

\subsection{Capitalisation et essaimage}
Une fois la ligne pilote stabilisée, un retour d’expérience (REX) est formalisé et diffusé. Les bonnes pratiques sont intégrées dans un guide interne. L’essaimage consiste à déployer la méthodologie sur d’autres lignes ou usines.

% ------------------------------------------------------------
\section{Annexes pratiques}

\subsection{Bibliographie enrichie}
\begin{itemize}
    \item Goldratt, E.M., Cox, J. – \textit{Le But} (The Goal), 1984.
    \item Ohno, T. – \textit{Toyota Production System}, 1988.
    \item Womack, J., Jones, D. – \textit{Lean Thinking}, 1996.
    \item Rother, M., Shook, J. – \textit{Learning to See}, 1999.
    \item Shingo, S. – \textit{A Revolution in Manufacturing: The SMED System}, 1985.
    \item Baudin, M. – \textit{Lean Logistics}, 2004.
    \item Hopp, W.J., Spearman, M.L. – \textit{Factory Physics}, 3e éd., 2008.
\end{itemize}

\subsection{Templates téléchargeables}
Les documents suivants sont disponibles sur le site d’accompagnement du cours (codes QR à insérer) :
\begin{itemize}
    \item \texttt{Fiche\_relevé\_temps.xlsx} : collecte des temps avec calcul automatique de la moyenne et de l’écart-type.
    \item \texttt{Diagramme\_antériorité.pptx} : gabarit visuel pour construire le graphe de précédence.
    \item \texttt{Calculateur\_équilibrage.xlsx} : outil Excel avec algorithmes intuitif, poids positionnel et COMSOAL.
    \item \texttt{Rapport\_rééquilibrage.docx} : modèle de rapport structuré (diagnostic, solution, résultats, plan d’actions).
\end{itemize}

\subsection{QCM et évaluations}
\textbf{QCM auto-évaluatif} (extrait) :

\begin{enumerate}
    \item Le Takt Time se calcule par : 
    \begin{enumerate}
        \item Temps de production disponible / Demande client
        \item Demande client / Temps de production disponible
        \item Somme des temps de cycle / Nombre de postes
    \end{enumerate}
    \item Quelle méthode d’équilibrage est la plus adaptée pour une ligne de plus de 50 tâches ?
    \begin{enumerate}
        \item Méthode intuitive
        \item Méthode du poids positionnel
        \item Algorithme génétique ou logiciel spécialisé
    \end{enumerate}
    \item Un poste avec \(T_C > T_T\) est :
    \begin{enumerate}
        \item Un goulot
        \item Un poste en surcapacité
        \item Un poste à équilibrer en priorité (les deux réponses précédentes)
    \end{enumerate}
\end{enumerate}
\textit{(Corrigés disponibles dans le fichier annexe\_corrige.pdf)}

% ============================================================
% FIN DES COMPLÉMENTS
% ============================================================
% ============================================================
% COMPLÉMENTS (suite) : EXERCICES PRATIQUES, CORRIGÉS ET JEUX PÉDAGOGIQUES
% ============================================================

\section{Exercices pratiques et corrigés}

\subsection{Exercice 1 : Calcul du Takt Time et diagnostic préliminaire}
\subsubsection{Énoncé}
Une ligne de production fonctionne 7h30 par jour (pause déjeuner incluse : 45 min). La demande client est de 450 pièces par jour. Les temps de cycle mesurés sur les 5 postes sont les suivants :

\begin{table}[H]
\centering
\begin{tabular}{|l|c|}
\hline
\textbf{Poste} & \textbf{Temps de cycle (min)} \\ \hline
Poste 1 & 0,85 \\
Poste 2 & 0,92 \\
Poste 3 & 0,78 \\
Poste 4 & 0,88 \\
Poste 5 & 0,72 \\ \hline
\end{tabular}
\caption{Temps de cycle initiaux}
\label{tab:ex1_temps}
\end{table}

\textbf{Questions :}
\begin{enumerate}
    \item Calculer le Takt Time.
    \item Calculer le temps total de traitement.
    \item Déterminer le nombre théorique de postes.
    \item Calculer l'efficacité d'équilibrage actuelle.
    \item Identifier le goulot et justifier.
\end{enumerate}

\subsubsection{Corrigé}
\begin{enumerate}
    \item Temps disponible = 7h30 - 0h45 = 6h45 = 405 min \\
    \[
    TT = \frac{405}{450} = 0,90 \text{ min/pièce} = 54 \text{ secondes/pièce}
    \]
    
    \item Somme des temps : \\
    \[
    \sum T_i = 0,85 + 0,92 + 0,78 + 0,88 + 0,72 = 4,15 \text{ min}
    \]
    
    \item Nombre théorique de postes : \\
    \[
    N_{th} = \left\lceil \frac{4,15}{0,90} \right\rceil = \lceil 4,61 \rceil = 5 \text{ postes}
    \]
    
    \item Efficacité actuelle : \\
    \[
    Eff = \frac{\sum T_i}{N \times T_{max}} = \frac{4,15}{5 \times 0,92} = \frac{4,15}{4,60} = 0,902 \text{ soit } 90,2\%
    \]
    
    \item Le goulot est le poste 2 avec \(T_C = 0,92\) min > TT (0,90 min). \\
    C'est le poste le plus lent, il limite la capacité de toute la ligne.
\end{enumerate}

\subsection{Exercice 2 : Construction du diagramme d’antériorité}
\subsubsection{Énoncé}
Une ligne d’assemblage comporte 7 tâches avec les temps et contraintes suivantes :

\begin{table}[H]
\centering
\begin{tabular}{|l|c|l|}
\hline
\textbf{Tâche} & \textbf{Temps (min)} & \textbf{Prédécesseurs} \\ \hline
A & 0,25 & - \\
B & 0,30 & A \\
C & 0,20 & A \\
D & 0,35 & B, C \\
E & 0,15 & D \\
F & 0,25 & D \\
G & 0,20 & E, F \\ \hline
\end{tabular}
\caption{Données des tâches}
\label{tab:ex2_taches}
\end{table}

Le Takt Time est de 0,80 min.

\textbf{Questions :}
\begin{enumerate}
    \item Construire le diagramme d’antériorité.
    \item Appliquer la méthode du poids positionnel.
    \item Proposer une répartition des tâches en postes.
    \item Calculer l’efficacité obtenue.
\end{enumerate}

\subsubsection{Corrigé}

\textbf{1. Diagramme d’antériorité :}

\begin{center}
\begin{tikzpicture}[node distance=1.5cm, auto]
    % Nodes
    \node[draw, rectangle, minimum width=1cm, minimum height=0.8cm] (A) {A (0,25)};
    \node[draw, rectangle, below left of=A, xshift=-1cm, yshift=-1cm] (B) {B (0,30)};
    \node[draw, rectangle, below right of=A, xshift=1cm, yshift=-1cm] (C) {C (0,20)};
    \node[draw, rectangle, below of=B, yshift=-0.5cm] (D) {D (0,35)};
    \node[draw, rectangle, below left of=D, xshift=-0.8cm, yshift=-0.8cm] (E) {E (0,15)};
    \node[draw, rectangle, below right of=D, xshift=0.8cm, yshift=-0.8cm] (F) {F (0,25)};
    \node[draw, rectangle, below of=E, yshift=-0.5cm] (G) {G (0,20)};
    
    % Arrows
    \draw[->] (A) -- (B);
    \draw[->] (A) -- (C);
    \draw[->] (B) -- (D);
    \draw[->] (C) -- (D);
    \draw[->] (D) -- (E);
    \draw[->] (D) -- (F);
    \draw[->] (E) -- (G);
    \draw[->] (F) -- (G);
\end{tikzpicture}
\end{center}

\textbf{2. Calcul des poids positionnels :}

\[
\begin{aligned}
PP_G &= 0,20 \\
PP_E &= 0,15 + 0,20 = 0,35 \\
PP_F &= 0,25 + 0,20 = 0,45 \\
PP_D &= 0,35 + 0,15 + 0,25 + 0,20 = 0,95 \\
PP_B &= 0,30 + 0,35 + 0,15 + 0,25 + 0,20 = 1,25 \\
PP_C &= 0,20 + 0,35 + 0,15 + 0,25 + 0,20 = 1,15 \\
PP_A &= 0,25 + 1,25 + 1,15 + 0,95 + 0,35 + 0,45 + 0,20 = 4,60
\end{aligned}
\]

\textbf{3. Assignation (TT = 0,80 min) :}

\begin{itemize}
    \item \textbf{Poste 1} : A (0,25) + B (0,30) = 0,55 ; C (0,20) = 0,75 (≤ TT) \\
    \textit{Total poste 1 : 0,75 min}
    
    \item \textbf{Poste 2} : D (0,35) + E (0,15) = 0,50 ; F (0,25) = 0,75 (≤ TT) \\
    \textit{Total poste 2 : 0,75 min}
    
    \item \textbf{Poste 3} : G (0,20) \\
    \textit{Total poste 3 : 0,20 min}
\end{itemize}

\textbf{4. Efficacité obtenue :}
\[
Eff = \frac{4,15}{3 \times 0,75} = \frac{4,15}{2,25} = 1,84 > 1
\]
(Attention : ce résultat >1 indique une erreur de calcul car la somme des temps est 4,15 et le temps max par poste est 0,75, ce qui donne un rapport >1. Vérifions : somme des temps = 0,25+0,30+0,20+0,35+0,15+0,25+0,20 = 1,70 min ?)

\textit{Correction :} La somme réelle est 1,70 min (et non 4,15). \\
\[
Eff = \frac{1,70}{3 \times 0,75} = \frac{1,70}{2,25} = 0,756 \text{ soit } 75,6\%
\]

\subsection{Exercice 3 : Équilibrage multi-modèles}
\subsubsection{Énoncé}
Une ligne doit produire deux références avec les temps de cycle suivants :

\begin{table}[H]
\centering
\begin{tabular}{|l|c|c|}
\hline
\textbf{Poste} & \textbf{Temps Réf A (min)} & \textbf{Temps Réf B (min)} \\ \hline
Poste 1 & 0,45 & 0,55 \\
Poste 2 & 0,50 & 0,42 \\
Poste 3 & 0,48 & 0,53 \\
Poste 4 & 0,52 & 0,45 \\ \hline
\end{tabular}
\caption{Temps par poste pour deux références}
\label{tab:ex3_multi}
\end{table}

La demande journalière est de 250 pièces pour A et 200 pièces pour B. Le temps disponible est de 420 min.

\textbf{Questions :}
\begin{enumerate}
    \item Calculer le Takt Time pondéré.
    \item Déterminer le nombre théorique de postes.
    \item Proposer un équilibrage multi-modèles.
\end{enumerate}

\subsubsection{Corrigé}

\textbf{1. Takt Time pondéré :}
\[
TT_{\text{pond}} = \frac{420}{(250 \times 1) + (200 \times 1)} = \frac{420}{450} = 0,933 \text{ min/pièce}
\]

\textbf{2. Nombre théorique :}
\[
\sum T_i \text{ moyen } = \frac{(0,45+0,50+0,48+0,52) + (0,55+0,42+0,53+0,45)}{2} = \frac{1,95 + 1,95}{2} = 1,95 \text{ min}
\]
\[
N_{th} = \left\lceil \frac{1,95}{0,933} \right\rceil = \lceil 2,09 \rceil = 3 \text{ postes}
\]

\textbf{3. Proposition d’équilibrage :}
En regroupant les postes 1+2 en un poste, 3+4 en un autre, on obtient pour la référence A : 0,45+0,50=0,95 min et 0,48+0,52=1,00 min ; pour B : 0,55+0,42=0,97 min et 0,53+0,45=0,98 min. Tous inférieurs à TTpond (0,933 ? Attention : 0,95 > 0,933 → goulot). Un dédoublement est nécessaire.

\subsection{Exercice 4 : Calcul du retour sur investissement}
\subsubsection{Énoncé}
Un projet d’équilibrage coûte 25 000 € (étude, formation, communication). Les gains estimés sont :
\begin{itemize}
    \item Réduction des temps morts : 12 000 €/an
    \item Réduction des stocks WIP : 8 000 €/an
    \item Gain de productivité : 18 000 €/an
    \item Réduction du lead time : 5 000 €/an
\end{itemize}

\textbf{Questions :}
\begin{enumerate}
    \item Calculer le gain annuel total.
    \item Calculer le ROI.
    \item Calculer le délai de retour sur investissement.
\end{enumerate}

\subsubsection{Corrigé}

\textbf{1. Gain annuel total :}
\[
G = 12 000 + 8 000 + 18 000 + 5 000 = 43 000 \text{ €/an}
\]

\textbf{2. ROI :}
\[
ROI = \frac{43 000 - 25 000}{25 000} \times 100 = \frac{18 000}{25 000} \times 100 = 72\%
\]

\textbf{3. Délai de retour :}
\[
DR = \frac{25 000}{43 000 / 12} = \frac{25 000}{3 583,33} \approx 7 \text{ mois}
\]

% ------------------------------------------------------------
\section{Jeux pédagogiques et simulations}

\subsection{Jeu de rôle : Équilibrage sur ligne LEGO®}
\subsubsection{Objectif}
Sensibiliser les participants aux principes de l’équilibrage par une simulation concrète et ludique.

\subsubsection{Matériel nécessaire}
\begin{itemize}
    \item Briques LEGO® (ou tout autre objet d’assemblage)
    \item Chronomètres
    \item Postes de travail identifiés (tables)
    \item Fiches d’instruction
\end{itemize}

\subsubsection{Déroulement}
\begin{enumerate}
    \item \textbf{Phase 1 – Situation déséquilibrée} (15 min) \\
    Les participants sont répartis sur 4 postes avec des tâches inégalement réparties. On mesure le flux, les stocks intermédiaires et le lead time.
    
    \item \textbf{Phase 2 – Diagnostic} (10 min) \\
    Collecte des temps, calcul du Takt Time, identification des goulots.
    
    \item \textbf{Phase 3 – Équilibrage} (20 min) \\
    Les participants proposent et testent une nouvelle répartition des tâches.
    
    \item \textbf{Phase 4 – Mesure et comparaison} (10 min) \\
    On mesure à nouveau les indicateurs et on compare avec la situation initiale.
\end{enumerate}

\subsubsection{Grille d’évaluation}

\begin{table}[H]
\centering
\begin{tabular}{|l|c|c|}
\hline
\textbf{Indicateur} & \textbf{Avant équilibrage} & \textbf{Après équilibrage} \\ \hline
Production (pièces/5 min) & & \\
Stocks WIP (nb pièces) & & \\
Lead Time (min) & & \\
Efficacité équilibrage & & \\ \hline
\end{tabular}
\caption{Grille de suivi du jeu LEGO®}
\label{tab:jeu_lego}
\end{table}

\subsection{Simulation Excel interactive}
Une feuille de calcul dynamique permet de visualiser en temps réel l’impact des regroupements de tâches sur l’efficacité d’équilibrage. Les participants peuvent :
\begin{itemize}
    \item Saisir les temps de cycle
    \item Définir le Takt Time
    \item Proposer des assignations de tâches
    \item Visualiser instantanément le diagramme d’équilibrage et le calcul d’efficacité
\end{itemize}

\subsection{Jeu de cartes « Balance Challenge »}
\subsubsection{Règles}
\begin{itemize}
    \item Chaque carte représente une tâche avec son temps et ses contraintes de précédence.
    \item Les joueurs (en équipes) doivent constituer des postes de travail (mains) en respectant les contraintes et en ne dépassant pas le Takt Time.
    \item L’équipe avec la meilleure efficacité d’équilibrage gagne.
    \item Variantes : introduction d’aléas (micro-pannes, absentéisme) pour simuler le réel.
\end{itemize}

% ------------------------------------------------------------
\section{Glossaire illustré et fiches mémo}

\subsection{Cartes mentales}
\subsubsection{Les 3 temps fondamentaux}

\begin{center}
\begin{tikzpicture}[scale=0.8, every node/.style={draw, rectangle, rounded corners, align=center}]
    \node (takt) at (0,0) {Takt Time\\Rythme client\\$TT = T_{dispo}/D$};
    \node (cycle) at (4,0) {Cycle Time\\Temps réel\\$TC = T_{prod}/Q$};
    \node (lead) at (8,0) {Lead Time\\Temps total\\$LT = \sum TC + \sum attente$};
    \draw[->, thick] (takt) -- (cycle) node[midway, above] {référence};
    \draw[->, thick] (cycle) -- (lead) node[midway, above] {contribue};
\end{tikzpicture}
\end{center}

\subsubsection{Les 4 méthodes d’équilibrage}

\begin{center}
\begin{tikzpicture}[scale=0.7]
    \node (intuitive) at (0,0) [draw, rectangle] {Intuitive\\Simple, rapide\\< 20 tâches};
    \node (candidat) at (4,0) [draw, rectangle] {Candidat le + grand\\Priorité aux tâches longues\\< 30 tâches};
    \node (poids) at (8,0) [draw, rectangle] {Poids positionnel\\Somme des successeurs\\< 50 tâches};
    \node (logiciel) at (12,0) [draw, rectangle] {Logiciel / IA\\Algorithmes avancés\\> 50 tâches};
\end{tikzpicture}
\end{center}

\subsection{Aide-mémoire des formules clés}

\begin{table}[H]
\centering
\begin{tabular}{|l|c|}
\hline
\textbf{Formule} & \textbf{Expression} \\ \hline
Takt Time & \(TT = \dfrac{T_{\text{dispo}}}{D}\) \\
Efficacité d’équilibrage & \(Eff = \dfrac{\sum T_i}{N \times T_{\text{max}}}\) \\
Nombre théorique de postes & \(N_{th} = \left\lceil \dfrac{\sum T_i}{TT} \right\rceil\) \\
Poids positionnel & \(PP_i = T_i + \sum_{j \in Succ(i)} T_j\) \\
Écart au Takt Time & \(\Delta_i = T_{C_i} - TT\) \\
Taux de service & \(TS = \dfrac{\text{Livré à temps}}{\text{Commande}} \times 100\) \\
\hline
\end{tabular}
\caption{Formules essentielles}
\label{tab:formules}
\end{table}

\subsection{Fiches synthétiques par phase}

\textbf{Phase 1 – Analyse et collecte :}
\begin{itemize}
    \item Calculer TT
    \item Mesurer les temps (min 10 mesures)
    \item Construire le diagramme d’antériorité
    \item Valider les données avec les opérateurs
\end{itemize}

\textbf{Phase 2 – Regroupement :}
\begin{itemize}
    \item Choisir la méthode adaptée
    \item Respecter précédences et TT
    \item Maximiser le remplissage de chaque poste
    \item Documenter l’assignation
\end{itemize}

\textbf{Phase 3 – Évaluation :}
\begin{itemize}
    \item Calculer l’efficacité
    \item Identifier le goulot
    \item Ajuster si efficacité < 85\%
\end{itemize}

\textbf{Phase 4 – Implémentation :}
\begin{itemize}
    \item Former les opérateurs
    \item Standardiser les postes
    \item Mettre en place le suivi quotidien
    \item Animer les routines d’amélioration
\end{itemize}

% ============================================================
% FIN DES COMPLÉMENTS
% ============================================================
% ============================================================
% COMPLÉMENTS (suite) : EXERCICES PRATIQUES, CORRIGÉS ET JEUX PÉDAGOGIQUES
% ============================================================

\section{Exercices pratiques et corrigés}

\subsection{Exercice 1 : Calcul du Takt Time et diagnostic préliminaire}
\subsubsection{Énoncé}
Une ligne de production fonctionne 7h30 par jour (pause déjeuner incluse : 45 min). La demande client est de 450 pièces par jour. Les temps de cycle mesurés sur les 5 postes sont les suivants :

\begin{table}[H]
\centering
\begin{tabular}{|l|c|}
\hline
\textbf{Poste} & \textbf{Temps de cycle (min)} \\ \hline
Poste 1 & 0,85 \\
Poste 2 & 0,92 \\
Poste 3 & 0,78 \\
Poste 4 & 0,88 \\
Poste 5 & 0,72 \\ \hline
\end{tabular}
\caption{Temps de cycle initiaux}
\label{tab:ex1_temps}
\end{table}

\textbf{Questions :}
\begin{enumerate}
    \item Calculer le Takt Time.
    \item Calculer le temps total de traitement.
    \item Déterminer le nombre théorique de postes.
    \item Calculer l'efficacité d'équilibrage actuelle.
    \item Identifier le goulot et justifier.
\end{enumerate}

\subsubsection{Corrigé}
\begin{enumerate}
    \item Temps disponible = 7h30 - 0h45 = 6h45 = 405 min \\
    \[
    TT = \frac{405}{450} = 0,90 \text{ min/pièce} = 54 \text{ secondes/pièce}
    \]
    
    \item Somme des temps : \\
    \[
    \sum T_i = 0,85 + 0,92 + 0,78 + 0,88 + 0,72 = 4,15 \text{ min}
    \]
    
    \item Nombre théorique de postes : \\
    \[
    N_{th} = \left\lceil \frac{4,15}{0,90} \right\rceil = \lceil 4,61 \rceil = 5 \text{ postes}
    \]
    
    \item Efficacité actuelle : \\
    \[
    Eff = \frac{\sum T_i}{N \times T_{max}} = \frac{4,15}{5 \times 0,92} = \frac{4,15}{4,60} = 0,902 \text{ soit } 90,2\%
    \]
    
    \item Le goulot est le poste 2 avec \(T_C = 0,92\) min > TT (0,90 min). \\
    C'est le poste le plus lent, il limite la capacité de toute la ligne.
\end{enumerate}

\subsection{Exercice 2 : Construction du diagramme d’antériorité}
\subsubsection{Énoncé}
Une ligne d’assemblage comporte 7 tâches avec les temps et contraintes suivantes :

\begin{table}[H]
\centering
\begin{tabular}{|l|c|l|}
\hline
\textbf{Tâche} & \textbf{Temps (min)} & \textbf{Prédécesseurs} \\ \hline
A & 0,25 & - \\
B & 0,30 & A \\
C & 0,20 & A \\
D & 0,35 & B, C \\
E & 0,15 & D \\
F & 0,25 & D \\
G & 0,20 & E, F \\ \hline
\end{tabular}
\caption{Données des tâches}
\label{tab:ex2_taches}
\end{table}

Le Takt Time est de 0,80 min.

\textbf{Questions :}
\begin{enumerate}
    \item Construire le diagramme d’antériorité.
    \item Appliquer la méthode du poids positionnel.
    \item Proposer une répartition des tâches en postes.
    \item Calculer l’efficacité obtenue.
\end{enumerate}

\subsubsection{Corrigé}

\textbf{1. Diagramme d’antériorité :}

\begin{center}
\begin{tikzpicture}[node distance=1.5cm, auto]
    % Nodes
    \node[draw, rectangle, minimum width=1cm, minimum height=0.8cm] (A) {A (0,25)};
    \node[draw, rectangle, below left of=A, xshift=-1cm, yshift=-1cm] (B) {B (0,30)};
    \node[draw, rectangle, below right of=A, xshift=1cm, yshift=-1cm] (C) {C (0,20)};
    \node[draw, rectangle, below of=B, yshift=-0.5cm] (D) {D (0,35)};
    \node[draw, rectangle, below left of=D, xshift=-0.8cm, yshift=-0.8cm] (E) {E (0,15)};
    \node[draw, rectangle, below right of=D, xshift=0.8cm, yshift=-0.8cm] (F) {F (0,25)};
    \node[draw, rectangle, below of=E, yshift=-0.5cm] (G) {G (0,20)};
    
    % Arrows
    \draw[->] (A) -- (B);
    \draw[->] (A) -- (C);
    \draw[->] (B) -- (D);
    \draw[->] (C) -- (D);
    \draw[->] (D) -- (E);
    \draw[->] (D) -- (F);
    \draw[->] (E) -- (G);
    \draw[->] (F) -- (G);
\end{tikzpicture}
\end{center}

\textbf{2. Calcul des poids positionnels :}

\[
\begin{aligned}
PP_G &= 0,20 \\
PP_E &= 0,15 + 0,20 = 0,35 \\
PP_F &= 0,25 + 0,20 = 0,45 \\
PP_D &= 0,35 + 0,15 + 0,25 + 0,20 = 0,95 \\
PP_B &= 0,30 + 0,35 + 0,15 + 0,25 + 0,20 = 1,25 \\
PP_C &= 0,20 + 0,35 + 0,15 + 0,25 + 0,20 = 1,15 \\
PP_A &= 0,25 + 1,25 + 1,15 + 0,95 + 0,35 + 0,45 + 0,20 = 4,60
\end{aligned}
\]

\textbf{3. Assignation (TT = 0,80 min) :}

\begin{itemize}
    \item \textbf{Poste 1} : A (0,25) + B (0,30) = 0,55 ; C (0,20) = 0,75 (≤ TT) \\
    \textit{Total poste 1 : 0,75 min}
    
    \item \textbf{Poste 2} : D (0,35) + E (0,15) = 0,50 ; F (0,25) = 0,75 (≤ TT) \\
    \textit{Total poste 2 : 0,75 min}
    
    \item \textbf{Poste 3} : G (0,20) \\
    \textit{Total poste 3 : 0,20 min}
\end{itemize}

\textbf{4. Efficacité obtenue :}
\[
Eff = \frac{4,15}{3 \times 0,75} = \frac{4,15}{2,25} = 1,84 > 1
\]
(Attention : ce résultat >1 indique une erreur de calcul car la somme des temps est 4,15 et le temps max par poste est 0,75, ce qui donne un rapport >1. Vérifions : somme des temps = 0,25+0,30+0,20+0,35+0,15+0,25+0,20 = 1,70 min ?)

\textit{Correction :} La somme réelle est 1,70 min (et non 4,15). \\
\[
Eff = \frac{1,70}{3 \times 0,75} = \frac{1,70}{2,25} = 0,756 \text{ soit } 75,6\%
\]

\subsection{Exercice 3 : Équilibrage multi-modèles}
\subsubsection{Énoncé}
Une ligne doit produire deux références avec les temps de cycle suivants :

\begin{table}[H]
\centering
\begin{tabular}{|l|c|c|}
\hline
\textbf{Poste} & \textbf{Temps Réf A (min)} & \textbf{Temps Réf B (min)} \\ \hline
Poste 1 & 0,45 & 0,55 \\
Poste 2 & 0,50 & 0,42 \\
Poste 3 & 0,48 & 0,53 \\
Poste 4 & 0,52 & 0,45 \\ \hline
\end{tabular}
\caption{Temps par poste pour deux références}
\label{tab:ex3_multi}
\end{table}

La demande journalière est de 250 pièces pour A et 200 pièces pour B. Le temps disponible est de 420 min.

\textbf{Questions :}
\begin{enumerate}
    \item Calculer le Takt Time pondéré.
    \item Déterminer le nombre théorique de postes.
    \item Proposer un équilibrage multi-modèles.
\end{enumerate}

\subsubsection{Corrigé}

\textbf{1. Takt Time pondéré :}
\[
TT_{\text{pond}} = \frac{420}{(250 \times 1) + (200 \times 1)} = \frac{420}{450} = 0,933 \text{ min/pièce}
\]

\textbf{2. Nombre théorique :}
\[
\sum T_i \text{ moyen } = \frac{(0,45+0,50+0,48+0,52) + (0,55+0,42+0,53+0,45)}{2} = \frac{1,95 + 1,95}{2} = 1,95 \text{ min}
\]
\[
N_{th} = \left\lceil \frac{1,95}{0,933} \right\rceil = \lceil 2,09 \rceil = 3 \text{ postes}
\]

\textbf{3. Proposition d’équilibrage :}
En regroupant les postes 1+2 en un poste, 3+4 en un autre, on obtient pour la référence A : 0,45+0,50=0,95 min et 0,48+0,52=1,00 min ; pour B : 0,55+0,42=0,97 min et 0,53+0,45=0,98 min. Tous inférieurs à TTpond (0,933 ? Attention : 0,95 > 0,933 → goulot). Un dédoublement est nécessaire.

\subsection{Exercice 4 : Calcul du retour sur investissement}
\subsubsection{Énoncé}
Un projet d’équilibrage coûte 25 000 € (étude, formation, communication). Les gains estimés sont :
\begin{itemize}
    \item Réduction des temps morts : 12 000 €/an
    \item Réduction des stocks WIP : 8 000 €/an
    \item Gain de productivité : 18 000 €/an
    \item Réduction du lead time : 5 000 €/an
\end{itemize}

\textbf{Questions :}
\begin{enumerate}
    \item Calculer le gain annuel total.
    \item Calculer le ROI.
    \item Calculer le délai de retour sur investissement.
\end{enumerate}

\subsubsection{Corrigé}

\textbf{1. Gain annuel total :}
\[
G = 12 000 + 8 000 + 18 000 + 5 000 = 43 000 \text{ €/an}
\]

\textbf{2. ROI :}
\[
ROI = \frac{43 000 - 25 000}{25 000} \times 100 = \frac{18 000}{25 000} \times 100 = 72\%
\]

\textbf{3. Délai de retour :}
\[
DR = \frac{25 000}{43 000 / 12} = \frac{25 000}{3 583,33} \approx 7 \text{ mois}
\]

% ------------------------------------------------------------
\section{Jeux pédagogiques et simulations}

\subsection{Jeu de rôle : Équilibrage sur ligne LEGO®}
\subsubsection{Objectif}
Sensibiliser les participants aux principes de l’équilibrage par une simulation concrète et ludique.

\subsubsection{Matériel nécessaire}
\begin{itemize}
    \item Briques LEGO® (ou tout autre objet d’assemblage)
    \item Chronomètres
    \item Postes de travail identifiés (tables)
    \item Fiches d’instruction
\end{itemize}

\subsubsection{Déroulement}
\begin{enumerate}
    \item \textbf{Phase 1 – Situation déséquilibrée} (15 min) \\
    Les participants sont répartis sur 4 postes avec des tâches inégalement réparties. On mesure le flux, les stocks intermédiaires et le lead time.
    
    \item \textbf{Phase 2 – Diagnostic} (10 min) \\
    Collecte des temps, calcul du Takt Time, identification des goulots.
    
    \item \textbf{Phase 3 – Équilibrage} (20 min) \\
    Les participants proposent et testent une nouvelle répartition des tâches.
    
    \item \textbf{Phase 4 – Mesure et comparaison} (10 min) \\
    On mesure à nouveau les indicateurs et on compare avec la situation initiale.
\end{enumerate}

\subsubsection{Grille d’évaluation}

\begin{table}[H]
\centering
\begin{tabular}{|l|c|c|}
\hline
\textbf{Indicateur} & \textbf{Avant équilibrage} & \textbf{Après équilibrage} \\ \hline
Production (pièces/5 min) & & \\
Stocks WIP (nb pièces) & & \\
Lead Time (min) & & \\
Efficacité équilibrage & & \\ \hline
\end{tabular}
\caption{Grille de suivi du jeu LEGO®}
\label{tab:jeu_lego}
\end{table}

\subsection{Simulation Excel interactive}
Une feuille de calcul dynamique permet de visualiser en temps réel l’impact des regroupements de tâches sur l’efficacité d’équilibrage. Les participants peuvent :
\begin{itemize}
    \item Saisir les temps de cycle
    \item Définir le Takt Time
    \item Proposer des assignations de tâches
    \item Visualiser instantanément le diagramme d’équilibrage et le calcul d’efficacité
\end{itemize}

\subsection{Jeu de cartes « Balance Challenge »}
\subsubsection{Règles}
\begin{itemize}
    \item Chaque carte représente une tâche avec son temps et ses contraintes de précédence.
    \item Les joueurs (en équipes) doivent constituer des postes de travail (mains) en respectant les contraintes et en ne dépassant pas le Takt Time.
    \item L’équipe avec la meilleure efficacité d’équilibrage gagne.
    \item Variantes : introduction d’aléas (micro-pannes, absentéisme) pour simuler le réel.
\end{itemize}

% ------------------------------------------------------------
\section{Glossaire illustré et fiches mémo}

\subsection{Cartes mentales}
\subsubsection{Les 3 temps fondamentaux}

\begin{center}
\begin{tikzpicture}[scale=0.8, every node/.style={draw, rectangle, rounded corners, align=center}]
    \node (takt) at (0,0) {Takt Time\\Rythme client\\$TT = T_{dispo}/D$};
    \node (cycle) at (4,0) {Cycle Time\\Temps réel\\$TC = T_{prod}/Q$};
    \node (lead) at (8,0) {Lead Time\\Temps total\\$LT = \sum TC + \sum attente$};
    \draw[->, thick] (takt) -- (cycle) node[midway, above] {référence};
    \draw[->, thick] (cycle) -- (lead) node[midway, above] {contribue};
\end{tikzpicture}
\end{center}

\subsubsection{Les 4 méthodes d’équilibrage}

\begin{center}
\begin{tikzpicture}[scale=0.7]
    \node (intuitive) at (0,0) [draw, rectangle] {Intuitive\\Simple, rapide\\< 20 tâches};
    \node (candidat) at (4,0) [draw, rectangle] {Candidat le + grand\\Priorité aux tâches longues\\< 30 tâches};
    \node (poids) at (8,0) [draw, rectangle] {Poids positionnel\\Somme des successeurs\\< 50 tâches};
    \node (logiciel) at (12,0) [draw, rectangle] {Logiciel / IA\\Algorithmes avancés\\> 50 tâches};
\end{tikzpicture}
\end{center}

\subsection{Aide-mémoire des formules clés}

\begin{table}[H]
\centering
\begin{tabular}{|l|c|}
\hline
\textbf{Formule} & \textbf{Expression} \\ \hline
Takt Time & \(TT = \dfrac{T_{\text{dispo}}}{D}\) \\
Efficacité d’équilibrage & \(Eff = \dfrac{\sum T_i}{N \times T_{\text{max}}}\) \\
Nombre théorique de postes & \(N_{th} = \left\lceil \dfrac{\sum T_i}{TT} \right\rceil\) \\
Poids positionnel & \(PP_i = T_i + \sum_{j \in Succ(i)} T_j\) \\
Écart au Takt Time & \(\Delta_i = T_{C_i} - TT\) \\
Taux de service & \(TS = \dfrac{\text{Livré à temps}}{\text{Commande}} \times 100\) \\
\hline
\end{tabular}
\caption{Formules essentielles}
\label{tab:formules}
\end{table}

\subsection{Fiches synthétiques par phase}

\textbf{Phase 1 – Analyse et collecte :}
\begin{itemize}
    \item Calculer TT
    \item Mesurer les temps (min 10 mesures)
    \item Construire le diagramme d’antériorité
    \item Valider les données avec les opérateurs
\end{itemize}

\textbf{Phase 2 – Regroupement :}
\begin{itemize}
    \item Choisir la méthode adaptée
    \item Respecter précédences et TT
    \item Maximiser le remplissage de chaque poste
    \item Documenter l’assignation
\end{itemize}

\textbf{Phase 3 – Évaluation :}
\begin{itemize}
    \item Calculer l’efficacité
    \item Identifier le goulot
    \item Ajuster si efficacité < 85\%
\end{itemize}

\textbf{Phase 4 – Implémentation :}
\begin{itemize}
    \item Former les opérateurs
    \item Standardiser les postes
    \item Mettre en place le suivi quotidien
    \item Animer les routines d’amélioration
\end{itemize}

% ============================================================
% FIN DES COMPLÉMENTS
% ============================================================
% ============================================================
% COMPLÉMENTS (suite) : CAS PRATIQUES PAR SECTEUR, OUTILS AVANCÉS ET ÉTUDES DE CAS COMPLÈTES
% ============================================================

\section{Cas pratiques par secteur d'activité}

\subsection{Cas automobile : Équilibrage d'une chaîne de montage moteur}

\subsubsection{Contexte}
Une usine de fabrication de moteurs thermiques doit équilibrer sa ligne d'assemblage final. La demande est de 320 moteurs par jour, pour un temps disponible de 7h20 (440 minutes). Les temps de changement sont de 15 minutes par série, avec 2 séries par jour.

\subsubsection{Données collectées}

\begin{table}[H]
\centering
\caption{Temps de cycle des opérations - Ligne moteur}
\label{tab:auto_temps}
\begin{tabular}{|c|l|c|c|}
\hline
\textbf{Tâche} & \textbf{Description} & \textbf{Temps (s)} & \textbf{Prédécesseurs} \\ \hline
1 & Préparation bloc & 45 & - \\
2 & Montage vilebrequin & 52 & 1 \\
3 & Montage pistons & 48 & 1 \\
4 & Montage bielles & 55 & 2, 3 \\
5 & Contrôle jeu axial & 25 & 4 \\
6 & Montage culasse & 62 & 4 \\
7 & Montage arbre à cames & 58 & 6 \\
8 & Réglage distribution & 42 & 7 \\
9 & Contrôle étanchéité & 38 & 5, 8 \\
10 & Montage accessoires & 35 & 9 \\
11 & Test final & 48 & 10 \\
\hline
\end{tabular}
\end{table}

\subsubsection{Calculs préliminaires}

\begin{itemize}
    \item Temps disponible : 440 min/jour
    \item Demande : 320 moteurs/jour
    \item Takt Time : \(TT = \frac{440 \times 60}{320} = \frac{26400}{320} = 82,5 \text{ secondes/pièce}\)
    \item Temps total : \(\sum T_i = 45+52+48+55+25+62+58+42+38+35+48 = 508 \text{ secondes}\)
    \item Nombre théorique : \(N_{th} = \left\lceil \frac{508}{82,5} \right\rceil = \lceil 6,16 \rceil = 7 \text{ postes}\)
\end{itemize}

\subsubsection{Application de la méthode du poids positionnel}

Calcul des poids positionnels (extrait) :

\begin{itemize}
    \item Tâche 11 : PP = 48
    \item Tâche 10 : PP = 35 + 48 = 83
    \item Tâche 9 : PP = 38 + 48 = 86
    \item Tâche 8 : PP = 42 + 38 + 48 = 128
    \item Tâche 7 : PP = 58 + 42 + 38 + 48 = 186
    \item Tâche 6 : PP = 62 + 58 + 42 + 38 + 48 = 248
    \item Tâche 5 : PP = 25 + 38 + 48 = 111
    \item Tâche 4 : PP = 55 + 25 + 62 + 58 + 42 + 38 + 48 = 328
    \item Tâche 3 : PP = 48 + 55 + 25 + 62 + 58 + 42 + 38 + 48 = 376
    \item Tâche 2 : PP = 52 + 55 + 25 + 62 + 58 + 42 + 38 + 48 = 380
    \item Tâche 1 : PP = 45 + 52 + 48 + 55 + 25 + 62 + 58 + 42 + 38 + 35 + 48 = 508
\end{itemize}

\subsubsection{Résultats de l'équilibrage}

\begin{table}[H]
\centering
\caption{Répartition des tâches après équilibrage}
\label{tab:auto_repartition}
\begin{tabular}{|c|l|c|c|}
\hline
\textbf{Poste} & \textbf{Tâches} & \textbf{Temps cumulé (s)} & \textbf{Écart TT} \\ \hline
Poste 1 & 1, 2 & 45 + 52 = 97 & +14,5 \\
Poste 2 & 3, 5 & 48 + 25 = 73 & -9,5 \\
Poste 3 & 4 & 55 & -27,5 \\
Poste 4 & 6, 8 & 62 + 42 = 104 & +21,5 \\
Poste 5 & 7 & 58 & -24,5 \\
Poste 6 & 9, 10 & 38 + 35 = 73 & -9,5 \\
Poste 7 & 11 & 48 & -34,5 \\
\hline
\end{tabular}
\end{table}

\subsubsection{Analyse et ajustements}

Les postes 1 et 4 dépassent le Takt Time (97s et 104s > 82,5s). Des ajustements sont nécessaires :

\begin{itemize}
    \item Dédoublement du poste 4 : répartir les tâches 6 et 8 sur deux postes
    \item Transfert de la tâche 2 vers un autre poste
    \item Optimisation SMED pour réduire les temps de changement
\end{itemize}

\subsubsection{Nouvelle proposition d'équilibrage}

\begin{table}[H]
\centering
\caption{Équilibrage optimisé après ajustements}
\label{tab:auto_optimise}
\begin{tabular}{|c|l|c|c|}
\hline
\textbf{Poste} & \textbf{Tâches} & \textbf{Temps cumulé (s)} & \textbf{Écart TT} \\ \hline
Poste 1 & 1 & 45 & -37,5 \\
Poste 2 & 2, 3 & 52 + 48 = 100 & +17,5 \\
Poste 3 & 4, 5 & 55 + 25 = 80 & -2,5 \\
Poste 4 & 6 & 62 & -20,5 \\
Poste 5 & 7, 8 & 58 + 42 = 100 & +17,5 \\
Poste 6 & 9, 10 & 38 + 35 = 73 & -9,5 \\
Poste 7 & 11 & 48 & -34,5 \\
\hline
\end{tabular}
\end{table}

\subsubsection{Résultats obtenus}
\begin{itemize}
    \item Efficacité finale : \(Eff = \frac{508}{7 \times 100} = \frac{508}{700} = 72,6\%\)
    \item Lead Time réduit de 35 min à 12 min
    \item Productivité augmentée de 18\%
\end{itemize}

\subsection{Cas électronique : High-mix low-volume}

\subsubsection{Contexte}
Une unité de production de cartes électroniques doit traiter 15 références différentes avec des volumes variant de 50 à 500 pièces par jour. La ligne comprend 8 postes flexibles.

\subsubsection{Méthodologie spécifique}

\paragraph{Calcul du Takt Time pondéré}
\[
TT_{\text{pond}} = \frac{T_{\text{dispo}}}{\sum_{i=1}^{n} (d_i \times t_i)}
\]

Avec les données suivantes :

\begin{table}[H]
\centering
\caption{Mix produit et temps de cycle}
\label{tab:elec_mix}
\begin{tabular}{|c|c|c|}
\hline
\textbf{Référence} & \textbf{Demande (pcs/jour)} & \textbf{Temps total (min)} \\ \hline
REF A & 500 & 8,5 \\
REF B & 300 & 6,2 \\
REF C & 200 & 7,8 \\
REF D & 100 & 9,1 \\
REF E & 50 & 5,4 \\
\hline
\end{tabular}
\end{table}

Temps disponible : 420 min/jour

\[
TT_{\text{pond}} = \frac{420}{(500 \times 1) + (300 \times 1) + (200 \times 1) + (100 \times 1) + (50 \times 1)} = \frac{420}{1150} = 0,365 \text{ min/pièce} = 21,9 \text{ secondes/pièce}
\]

\paragraph{Stratégie d'équilibrage multi-modèles}
\begin{enumerate}
    \item Identification des opérations communes à toutes les références
    \item Regroupement des opérations spécifiques par famille de produits
    \item Mise en place de cellules flexibles avec opérateurs polyvalents
    \item Utilisation du séquencement Heijunka pour lisser la production
\end{enumerate}

\paragraph{Résultats}
\begin{itemize}
    \item Taux de service : 95\% → 98,5\%
    \item Lead Time moyen : 4,2 jours → 1,8 jour
    \item Flexibilité : capacité à changer de référence en moins de 5 minutes
\end{itemize}

\subsection{Cas agroalimentaire : Conditionnement de produits frais}

\subsubsection{Contraintes spécifiques}
\begin{itemize}
    \item Nettoyage obligatoire entre les lots (30 min)
    \item Traçabilité renforcée
    \item Température contrôlée (limitation du temps hors chaîne du froid)
    \item Cadences variables selon la saisonnalité
\end{itemize}

\subsubsection{Modélisation des temps de nettoyage}

Les temps de nettoyage sont intégrés comme des tâches avec contraintes de précédence :

\[
T_{\text{dispo net}} = T_{\text{dispo}} - \sum_{i=1}^{n} T_{\text{nettoyage}_i} \times Nb_{\text{séries}}
\]

\subsubsection{Équilibrage avec contraintes sanitaires}

\begin{figure}[H]
\centering
\begin{tikzpicture}[node distance=1cm, auto]
    % Production blocks
    \node[draw, rectangle, minimum width=2cm, minimum height=0.8cm] (prod1) {Lot A};
    \node[draw, rectangle, right of=prod1, xshift=1cm] (clean1) {Nettoyage};
    \node[draw, rectangle, right of=clean1, xshift=1cm] (prod2) {Lot B};
    \node[draw, rectangle, right of=prod2, xshift=1cm] (clean2) {Nettoyage};
    \node[draw, rectangle, right of=clean2, xshift=1cm] (prod3) {Lot C};
    
    \draw[->] (prod1) -- (clean1);
    \draw[->] (clean1) -- (prod2);
    \draw[->] (prod2) -- (clean2);
    \draw[->] (clean2) -- (prod3);
\end{tikzpicture}
\caption{Intégration des cycles de nettoyage dans la séquence de production}
\label{fig:clean_sequence}
\end{figure}

\subsubsection{Indicateurs spécifiques}
\begin{itemize}
    \item Taux de respect de la chaîne du froid : 100\%
    \item Temps de nettoyage / temps de production : 12\% → 8\%
    \item Conformité hygiène : 98\% → 99,5\%
\end{itemize}

% ------------------------------------------------------------
\section{Outils avancés et modélisation}

\subsection{Programmation linéaire pour l'équilibrage}

\subsubsection{Formulation mathématique complète}

Minimiser \(N\) (nombre de postes)

Sous les contraintes :

\[
\begin{aligned}
&\sum_{k=1}^{K} x_{ik} = 1 \quad \forall i = 1,\ldots,n \\
&\sum_{i=1}^{n} T_i \cdot x_{ik} \leq TT \quad \forall k = 1,\ldots,K \\
&\sum_{k=1}^{K} k \cdot x_{ik} \leq \sum_{k=1}^{K} k \cdot x_{jk} \quad \forall (i,j) \text{ où } i \text{ précède } j \\
&x_{ik} \in \{0,1\} \quad \forall i,k
\end{aligned}
\]

Avec :
\begin{itemize}
    \item \(n\) : nombre de tâches
    \item \(K\) : nombre maximum de postes (borne supérieure)
    \item \(x_{ik}\) : variable binaire = 1 si tâche i assignée au poste k
    \item \(T_i\) : temps de la tâche i
    \item \(TT\) : Takt Time
\end{itemize}

\subsubsection{Résolution avec Excel Solver}

\begin{enumerate}
    \item Créer une matrice d'assignation \(n \times K\)
    \item Définir la fonction objectif : minimiser le nombre de postes utilisés
    \item Ajouter les contraintes de précédence sous forme linéaire
    \item Utiliser le solveur avec méthode Simplex ou Evolutionary
\end{enumerate}

\subsection{Simulation Monte Carlo}

\subsubsection{Principe}
La simulation Monte Carlo permet de prendre en compte la variabilité des temps de cycle et d'évaluer la robustesse de l'équilibrage.

\subsubsection{Modélisation des temps de cycle}

Pour chaque tâche, on définit une loi de distribution :
\begin{itemize}
    \item Loi normale : \(T \sim \mathcal{N}(\mu, \sigma^2)\) pour les processus stables
    \item Loi triangulaire : \(T \sim \text{Triang}(a, b, c)\) pour les processus avec incertitudes
    \item Loi log-normale : \(T \sim \text{LogN}(\mu, \sigma)\) pour les temps avec asymétrie positive
\end{itemize}

\subsubsection{Algorithme de simulation}

\begin{verbatim}
Pour itération = 1 à N_simulations:
    Pour chaque tâche i:
        Générer t_i selon sa loi de distribution
    Pour chaque poste j:
        Calculer T_poste = somme(t_i) pour i dans poste_j
        Stocker T_poste
    Calculer efficacité = somme(t_i) / (N * max(T_poste))
    Stocker efficacité

Calculer:
    - Efficacité moyenne
    - Intervalle de confiance
    - Probabilité d'avoir T_poste > TT
\end{verbatim}

\subsubsection{Exemple de résultats}

\begin{table}[H]
\centering
\caption{Résultats de simulation Monte Carlo (1000 itérations)}
\label{tab:monte_carlo}
\begin{tabular}{|l|c|c|}
\hline
\textbf{Indicateur} & \textbf{Valeur moyenne} & \textbf{Intervalle 95\%} \\ \hline
Efficacité équilibrage & 86,4\% & [82,1\% ; 90,7\%] \\
Probabilité goulot (poste 3) & 32\% & - \\
Taux de service atteint & 94\% & [91\% ; 97\%] \\
\hline
\end{tabular}
\end{table}

\subsection{Analyse de sensibilité}

\subsubsection{Objectif}
Identifier les paramètres les plus influents sur la performance de l'équilibrage.

\subsubsection{Méthode}
\begin{enumerate}
    \item Faire varier chaque paramètre de ±20\%
    \item Mesurer l'impact sur l'efficacité et le lead time
    \item Classer les paramètres par ordre d'influence
\end{enumerate}

\begin{table}[H]
\centering
\caption{Analyse de sensibilité - Impact sur l'efficacité}
\label{tab:sensibilite}
\begin{tabular}{|l|c|c|}
\hline
\textbf{Paramètre} & \textbf{Variation} & \textbf{Impact efficacité} \\ \hline
Temps de cycle poste goulot & +10\% & -15\% \\
Demande client & +20\% & -12\% \\
Temps de changement & +50\% & -8\% \\
Taux de disponibilité & -10\% & -6\% \\
\hline
\end{tabular}
\end{table}

% ------------------------------------------------------------
\section{Études de cas complètes}

\subsection{Étude de cas 1 : Transformation d'une ligne d'assemblage industrielle}

\subsubsection{Situation initiale}
Une entreprise de fabrication de composants mécaniques dispose d'une ligne d'assemblage de 12 postes produisant 800 pièces par jour sur 8h. Les indicateurs montrent des signes de déséquilibre :

\begin{itemize}
    \item Efficacité : 71\%
    \item Lead Time : 3,2 heures
    \item WIP total : 450 pièces
    \item Taux de service : 87\%
    \item 3 goulots identifiés
\end{itemize}

\subsubsection{Démarche suivie}

\paragraph{Phase 1 : Diagnostic approfondi (2 semaines)}
\begin{itemize}
    \item Réalisation d'un VSM complet
    \item Mesure des temps sur 3 semaines (5000 mesures)
    \item Construction du diagramme d'antériorité (42 tâches)
    \item Analyse des causes des goulots
\end{itemize}

\paragraph{Phase 2 : Conception des solutions (1 semaine)}
\begin{itemize}
    \item Calcul du nouveau Takt Time après validation demande (850 pièces/jour)
    \item Utilisation de la méthode du poids positionnel
    \item Simulation de 5 scénarios d'équilibrage
    \item Validation des solutions avec les équipes terrain
\end{itemize}

\paragraph{Phase 3 : Mise en œuvre (2 semaines)}
\begin{itemize}
    \item Réorganisation des postes de 12 à 10 postes
    \item Formation des 24 opérateurs (2 équipes)
    \item Mise en place des instructions de travail standardisées
    \item Installation du tableau de bord visuel
\end{itemize}

\paragraph{Phase 4 : Suivi et ajustements (1 mois)}
\begin{itemize}
    \item Suivi quotidien des indicateurs
    \item 15 ajustements mineurs
    \item Optimisation SMED sur les 2 postes goulots
\end{itemize}

\subsubsection{Résultats obtenus après 3 mois}

\begin{table}[H]
\centering
\caption{Évolution des indicateurs - Étude de cas 1}
\label{tab:cas1_resultats}
\begin{tabular}{|l|c|c|c|}
\hline
\textbf{Indicateur} & \textbf{Initial} & \textbf{Objectif} & \textbf{Réalisé} \\ \hline
Efficacité équilibrage & 71\% & > 85\% & 89\% \\
Lead Time & 3,2 h & < 2,0 h & 1,8 h \\
WIP total & 450 pcs & < 200 pcs & 165 pcs \\
Taux de service & 87\% & > 95\% & 96\% \\
Productivité & 100 pcs/h & 110 pcs/h & 118 pcs/h \\
\hline
\end{tabular}
\end{table}

\subsubsection{Retour sur investissement}
\begin{itemize}
    \item Coût du projet : 35 000 €
    \item Gains annuels : 98 000 €
    \item ROI : 180\%
    \item Délai de retour : 4,3 mois
\end{itemize}

\subsection{Étude de cas 2 : Équilibrage dynamique en contexte de forte variabilité}

\subsubsection{Contexte}
Une entreprise de sous-traitance automobile fait face à une demande très fluctuante (variations de ±30\% selon les semaines) et à une diversité de produits importante (25 références actives).

\subsubsection{Solution déployée}

\paragraph{Mise en place d'un équilibrage adaptatif}
\begin{itemize}
    \item Calcul hebdomadaire du Takt Time basé sur les commandes fermes
    \item Matrice de polyvalence des opérateurs (3 postes maîtrisés en moyenne)
    \item Postes modulaires avec flexibilité de configuration
    \item Stock tampon contrôlé devant les goulots potentiels
\end{itemize}

\paragraph{Modèle d'équilibrage adaptatif}

\begin{figure}[H]
\centering
\begin{tikzpicture}[node distance=1.5cm, auto]
    \node[draw, rectangle, align=center] (demande) {Demande\\hebdomadaire};
    \node[draw, rectangle, right of=demande, xshift=2cm, align=center] (takt) {Calcul\\TT};
    \node[draw, rectangle, right of=takt, xshift=2cm, align=center] (postes) {Détermination\\nb postes};
    \node[draw, rectangle, below of=postes, yshift=-1cm, align=center] (affectation) {Affectation\\ressources};
    \node[draw, rectangle, left of=affectation, xshift=-2cm, align=center] (polyvalence) {Matrice\\polyvalence};
    \node[draw, rectangle, left of=polyvalence, xshift=-2cm, align=center] (execution) {Exécution\\suivi};
    
    \draw[->] (demande) -- (takt);
    \draw[->] (takt) -- (postes);
    \draw[->] (postes) -- (affectation);
    \draw[->] (polyvalence) -- (affectation);
    \draw[->] (affectation) -- (execution);
    \draw[->] (execution) edge[loop below] (execution);
\end{tikzpicture}
\caption{Cycle d'équilibrage adaptatif hebdomadaire}
\label{fig:adaptatif}
\end{figure}

\subsubsection{Résultats sur 12 mois}

\begin{table}[H]
\centering
\caption{Performance du système adaptatif}
\label{tab:cas2_resultats}
\begin{tabular}{|l|c|c|}
\hline
\textbf{Indicateur} & \textbf{Avant} & \textbf{Après} \\ \hline
Efficacité moyenne & 73\% & 86\% \\
Écart-type efficacité & 12\% & 5\% \\
Taux de service & 91\% & 97\% \\
Temps d'adaptation (changement TT) & 3 jours & 0,5 jour \\
Heures supplémentaires & 8\% & 3\% \\
\hline
\end{tabular}
\end{table}

\subsubsection{Enseignements clés}
\begin{enumerate}
    \item La polyvalence est un facteur clé de réussite (investissement formation)
    \item Les routines hebdomadaires d'équilibrage deviennent un standard
    \item Les outils numériques (tableau de bord temps réel) sont indispensables
    \item L'implication des équipes dans le calcul hebdomadaire du TT favorise l'adhésion
\end{enumerate}

% ------------------------------------------------------------
\section{Exercices de synthèse et évaluation finale}

\subsection{Exercice de synthèse 1 : Équilibrage complet}

\subsubsection{Énoncé}
Une ligne de production comporte 15 tâches dont les données sont fournies dans le tableau ci-dessous. La demande client est de 450 pièces par jour. Le temps disponible est de 7h30 (pauses déduites).

\begin{table}[H]
\centering
\caption{Données des tâches pour l'exercice de synthèse}
\label{tab:synth1_donnees}
\begin{tabular}{|c|c|l|}
\hline
\textbf{Tâche} & \textbf{Temps (min)} & \textbf{Prédécesseurs} \\ \hline
1 & 0,35 & - \\
2 & 0,42 & - \\
3 & 0,28 & 1 \\
4 & 0,31 & 1, 2 \\
5 & 0,25 & 3 \\
6 & 0,38 & 4 \\
7 & 0,22 & 5 \\
8 & 0,45 & 6, 7 \\
9 & 0,33 & 8 \\
10 & 0,29 & 8 \\
11 & 0,27 & 9, 10 \\
12 & 0,24 & 11 \\
13 & 0,36 & 12 \\
14 & 0,21 & 13 \\
15 & 0,32 & 14 \\
\hline
\end{tabular}
\end{table}

\textbf{Travail demandé :}
\begin{enumerate}
    \item Calculer le Takt Time
    \item Construire le diagramme d'antériorité
    \item Calculer le nombre théorique de postes
    \item Appliquer la méthode du poids positionnel
    \item Proposer une répartition des tâches en postes
    \item Calculer l'efficacité d'équilibrage obtenue
    \item Identifier le goulot et proposer des améliorations
    \item Calculer le lead time théorique après équilibrage
\end{enumerate}

\subsection{Évaluation finale - Questionnaire à choix multiples}

\subsubsection{Partie 1 : Concepts fondamentaux (10 questions)}

\begin{enumerate}
    \item Le Takt Time est défini comme :
    \begin{enumerate}
        \item Le temps total de fabrication d'une pièce
        \item Le rythme de production dicté par la demande client
        \item Le temps réel entre deux pièces produites
        \item Le temps de changement de série
    \end{enumerate}
    
    \item Un goulot d'étranglement se caractérise par :
    \begin{enumerate}
        \item Un temps de cycle inférieur au Takt Time
        \item Un temps de cycle supérieur au Takt Time
        \item Un temps de cycle égal au Takt Time
        \item Un stock nul en amont
    \end{enumerate}
    
    \item La formule de l'efficacité d'équilibrage est :
    \begin{enumerate}
        \item \(Eff = \frac{N \times T_{max}}{\sum T_i}\)
        \item \(Eff = \frac{\sum T_i}{N \times T_{max}}\)
        \item \(Eff = \frac{\sum T_i}{N \times TT}\)
        \item \(Eff = \frac{TT}{T_{max}}\)
    \end{enumerate}
    
    \item La méthode du poids positionnel consiste à :
    \begin{enumerate}
        \item Prioriser les tâches les plus longues
        \item Prioriser les tâches avec le plus de successeurs pondérés par les temps
        \item Prioriser les tâches sans prédécesseur
        \item Prioriser les tâches avec le plus de prédécesseurs
    \end{enumerate}
    
    \item Dans un diagramme d'antériorité, une flèche indique :
    \begin{enumerate}
        \item Une relation de précédence
        \item Un flux de matière
        \item Un déplacement d'opérateur
        \item Un temps d'attente
    \end{enumerate}
    
    \item Le Lead Time correspond à :
    \begin{enumerate}
        \item La somme des temps de cycle
        \item La somme des temps de cycle et des temps d'attente
        \item Le temps entre deux pièces produites
        \item Le temps de changement de série
    \end{enumerate}
    
    \item Un équilibrage parfait est atteint lorsque :
    \begin{enumerate}
        \item Tous les postes ont le même temps de cycle
        \item Tous les postes ont un temps de cycle égal au Takt Time
        \item Le nombre de postes est minimal
        \item Les stocks intermédiaires sont nuls
    \end{enumerate}
    
    \item Le SMED est une méthode qui vise à :
    \begin{enumerate}
        \item Équilibrer les lignes de production
        \item Réduire les temps de changement de série
        \item Standardiser les opérations
        \item Former les opérateurs à la polyvalence
    \end{enumerate}
    
    \item Dans le VSM, un goulot se caractérise par :
    \begin{enumerate}
        \item Un WIP important en amont
        \item Un WIP important en aval
        \item Un taux de disponibilité faible
        \item Un temps de changement long
    \end{enumerate}
    
    \item La polyvalence des opérateurs permet :
    \begin{enumerate}
        \item De réduire le nombre de postes
        \item D'augmenter la flexibilité de l'équilibrage
        \item De diminuer le Takt Time
        \item De supprimer les contraintes de précédence
    \end{enumerate}
\end{enumerate}

\subsubsection{Partie 2 : Exercice de calcul (20 points)}

Une ligne de production a les caractéristiques suivantes après équilibrage :

\begin{table}[H]
\centering
\begin{tabular}{|c|c|}
\hline
\textbf{Poste} & \textbf{Temps de cycle (min)} \\ \hline
1 & 0,68 \\
2 & 0,72 \\
3 & 0,65 \\
4 & 0,70 \\
5 & 0,66 \\
\hline
\end{tabular}
\end{table}

Le Takt Time est de 0,70 min. La somme des temps de toutes les tâches est de 3,41 min.

\textbf{Questions :}
\begin{enumerate}
    \item Calculer l'efficacité d'équilibrage. (5 pts)
    \item Identifier le ou les goulots. (3 pts)
    \item Calculer le nombre théorique de postes. (4 pts)
    \item Proposer une amélioration simple pour optimiser l'équilibrage. (4 pts)
    \item Calculer le nouveau temps de cycle du poste goulot après amélioration. (4 pts)
\end{enumerate}

\subsection{Corrigé de l'exercice de calcul}

\begin{enumerate}
    \item Efficacité d'équilibrage :
    \[
    Eff = \frac{3,41}{5 \times 0,72} = \frac{3,41}{3,60} = 0,947 \text{ soit } 94,7\%
    \]
    
    \item Goulot : poste 2 avec \(T_C = 0,72 > TT = 0,70\) min
    
    \item Nombre théorique de postes :
    \[
    N_{th} = \left\lceil \frac{3,41}{0,70} \right\rceil = \lceil 4,87 \rceil = 5 \text{ postes}
    \]
    
    \item Amélioration : transférer une tâche du poste 2 (0,72 min) vers le poste 3 (0,65 min) pour équilibrer les charges.
    
    \item Après transfert de 0,04 min de tâche :
    \begin{itemize}
        \item Poste 2 : 0,72 - 0,04 = 0,68 min
        \item Poste 3 : 0,65 + 0,04 = 0,69 min
        \item Nouveau goulot : poste 3 (0,69 min) < TT (0,70 min) → plus de goulot
    \end{itemize}
\end{enumerate}

% ============================================================
% FIN DES COMPLÉMENTS
% ============================================================
% ============================================================
% COMPLÉMENTS (suite) : OUTILS DE VISUALISATION AVANCÉS, TEMPLATES ET RESSOURCES MULTIMÉDIA
% ============================================================

\section{Outils de visualisation avancés}

\subsection{Diagramme de Gantt pour l'équilibrage}

Le diagramme de Gantt permet de visualiser la séquence des tâches au sein de chaque poste et d'identifier les temps morts.

\subsubsection{Construction du diagramme de Gantt d'équilibrage}

\begin{figure}[H]
\centering
\begin{tikzpicture}[scale=0.8, every node/.style={font=\small}]
    % Axes
    \draw[->] (0,0) -- (12,0) node[right] {Temps (minutes)};
    \draw[->] (0,0) -- (0,5) node[above] {Postes};
    
    % Takt Time line
    \draw[dashed, red, thick] (8,0) -- (8,5) node[above, red] {TT = 0,80 min};
    
    % Poste 1
    \draw[fill=blue!30] (0,4) rectangle (2.5,4.5) node[pos=0.5, above] {A (0,25)};
    \draw[fill=blue!30] (2.5,4) rectangle (5.5,4.5) node[pos=0.5, above] {B (0,30)};
    \draw[fill=blue!30] (5.5,4) rectangle (7.8,4.5) node[pos=0.5, above] {C (0,23)};
    \draw[fill=gray!50] (7.8,4) rectangle (8,4.5) node[pos=0.5, above] {};
    \node at (9,4.25) {Total: 0,78 min};
    
    % Poste 2
    \draw[fill=green!30] (0,3) rectangle (4.2,3.5) node[pos=0.5, above] {D (0,42)};
    \draw[fill=green!30] (4.2,3) rectangle (7.2,3.5) node[pos=0.5, above] {E (0,30)};
    \draw[fill=gray!50] (7.2,3) rectangle (8,3.5) node[pos=0.5, above] {};
    \node at (9,3.25) {Total: 0,72 min};
    
    % Poste 3
    \draw[fill=yellow!30] (0,2) rectangle (5.0,2.5) node[pos=0.5, above] {F (0,50)};
    \draw[fill=yellow!30] (5.0,2) rectangle (7.5,2.5) node[pos=0.5, above] {G (0,25)};
    \draw[fill=gray!50] (7.5,2) rectangle (8,2.5) node[pos=0.5, above] {};
    \node at (9,2.25) {Total: 0,75 min};
    
    % Poste 4
    \draw[fill=orange!30] (0,1) rectangle (2.8,1.5) node[pos=0.5, above] {H (0,28)};
    \draw[fill=orange!30] (2.8,1) rectangle (5.3,1.5) node[pos=0.5, above] {I (0,25)};
    \draw[fill=orange!30] (5.3,1) rectangle (8,1.5) node[pos=0.5, above] {J (0,27)};
    \node at (9,1.25) {Total: 0,80 min};
    
    % Labels postes
    \node at (-0.5,4.25) {Poste 1};
    \node at (-0.5,3.25) {Poste 2};
    \node at (-0.5,2.25) {Poste 3};
    \node at (-0.5,1.25) {Poste 4};
    
\end{tikzpicture}
\caption{Diagramme de Gantt d'équilibrage avec identification des temps morts}
\label{fig:gantt}
\end{figure}

\subsection{Carte de contrôle des temps de cycle}

Les cartes de contrôle permettent de surveiller la stabilité des temps de cycle dans le temps et de détecter les dérives avant qu'elles n'impactent l'équilibrage.

\subsubsection{Construction d'une carte de contrôle}

\begin{figure}[H]
\centering
\begin{tikzpicture}[scale=0.9]
    % Axes
    \draw[->] (0,0) -- (12,0) node[right] {Échantillons (jours)};
    \draw[->] (0,0) -- (0,5) node[above] {Temps de cycle (min)};
    
    % Ligne centrale
    \draw[blue, thick] (0,2.5) -- (12,2.5) node[right, blue] {Moyenne = 0,70 min};
    
    % Limites de contrôle
    \draw[dashed, red] (0,3.2) -- (12,3.2) node[right, red] {UCL = 0,85 min};
    \draw[dashed, red] (0,1.8) -- (12,1.8) node[right, red] {LCL = 0,55 min};
    
    % Données simulées
    \draw[thick, black] (1,2.6) -- (2,2.7) -- (3,2.5) -- (4,2.8) -- (5,2.4) -- (6,3.0) -- (7,2.6) -- (8,2.5) -- (9,2.7) -- (10,2.4) -- (11,3.1) -- (12,2.6);
    
    % Points hors limites
    \fill[red] (6,3.0) circle (0.08);
    \fill[red] (11,3.1) circle (0.08);
    
    % Annotations
    \node[red, above] at (6,3.1) {Hors limite - Alerte};
    \node[red, above] at (11,3.2) {Dérive détectée};
    
\end{tikzpicture}
\caption{Carte de contrôle des temps de cycle - Poste goulot}
\label{fig:control_chart}
\end{figure}

\subsection{Heatmap des charges par poste}

La heatmap permet de visualiser instantanément la répartition des charges et d'identifier visuellement les déséquilibres.

\begin{table}[H]
\centering
\caption{Heatmap des charges par poste et par produit}
\label{tab:heatmap}
\begin{tabular}{|l|c|c|c|c|c|}
\hline
\textbf{Poste} & \textbf{Produit A} & \textbf{Produit B} & \textbf{Produit C} & \textbf{Produit D} & \textbf{Charge moyenne} \\ \hline
Poste 1 & \cellcolor{green!80} 0,45 & \cellcolor{green!60} 0,52 & \cellcolor{yellow!80} 0,68 & \cellcolor{red!60} 0,85 & 0,63 \\
Poste 2 & \cellcolor{green!70} 0,48 & \cellcolor{green!50} 0,55 & \cellcolor{yellow!60} 0,72 & \cellcolor{red!40} 0,78 & 0,63 \\
Poste 3 & \cellcolor{green!90} 0,42 & \cellcolor{green!80} 0,48 & \cellcolor{green!60} 0,58 & \cellcolor{yellow!80} 0,68 & 0,54 \\
Poste 4 & \cellcolor{red!80} 0,92 & \cellcolor{red!70} 0,88 & \cellcolor{red!60} 0,82 & \cellcolor{red!50} 0,75 & 0,84 \\
Poste 5 & \cellcolor{green!60} 0,52 & \cellcolor{green!50} 0,58 & \cellcolor{yellow!70} 0,65 & \cellcolor{yellow!60} 0,70 & 0,61 \\
\hline
\multicolumn{6}{|c|}{\cellcolor{green!50} Sous-charge \cellcolor{yellow!50} Charge nominale \cellcolor{red!50} Surcharge} \\
\end{tabular}
\end{table}

% ------------------------------------------------------------
\section{Templates et formulaires opérationnels}

\subsection{Template 1 : Fiche de collecte des temps de cycle}

\begin{table}[H]
\centering
\caption{Fiche de collecte des temps de cycle - Version terrain}
\label{tab:template_temps}
\begin{tabular}{|c|l|c|c|c|c|c|c|c|c|c|c|}
\hline
\textbf{Tâche} & \textbf{Opérateur} & \textbf{M1} & \textbf{M2} & \textbf{M3} & \textbf{M4} & \textbf{M5} & \textbf{M6} & \textbf{M7} & \textbf{M8} & \textbf{M9} & \textbf{M10} \\ \hline
A & & & & & & & & & & & \\
B & & & & & & & & & & & \\
C & & & & & & & & & & & \\
D & & & & & & & & & & & \\
E & & & & & & & & & & & \\
\hline
\multicolumn{12}{|c|}{Moyenne : \hspace{2cm} Min : \hspace{2cm} Max : \hspace{2cm} Écart-type :} \\
\hline
\end{tabular}
\end{table}

\subsection{Template 2 : Diagramme d'antériorité vierge}

\begin{figure}[H]
\centering
\begin{tikzpicture}[node distance=1.8cm, auto]
    % Grille de positionnement
    \draw[step=1.5, gray, very thin] (0,0) grid (9,6);
    
    % Exemple de nœuds à compléter
    \node[draw, circle, minimum size=0.8cm] (t1) at (1.5,5) {T1};
    \node[draw, circle, minimum size=0.8cm] (t2) at (3,5) {T2};
    \node[draw, circle, minimum size=0.8cm] (t3) at (4.5,5) {T3};
    \node[draw, circle, minimum size=0.8cm] (t4) at (6,5) {T4};
    \node[draw, circle, minimum size=0.8cm] (t5) at (7.5,5) {T5};
    
    \node at (4.5,2) {À compléter avec les tâches et leurs précédences};
\end{tikzpicture}
\caption{Template de diagramme d'antériorité}
\label{fig:template_precedence}
\end{figure}

\subsection{Template 3 : Tableau d'assignation des tâches}

\begin{table}[H]
\centering
\caption{Tableau d'assignation des tâches par poste}
\label{tab:template_assignation}
\begin{tabular}{|c|c|l|c|c|c|}
\hline
\textbf{Poste} & \textbf{Tâches} & \textbf{Description} & \textbf{Temps unitaire} & \textbf{Temps cumulé} & \textbf{Observations} \\ \hline
1 & & & & & \\
 & & & & & \\
\hline
2 & & & & & \\
 & & & & & \\
\hline
3 & & & & & \\
 & & & & & \\
\hline
4 & & & & & \\
 & & & & & \\
\hline
\multicolumn{4}{|c|}{Somme des temps} & \multicolumn{2}{c|}{} \\
\multicolumn{4}{|c|}{Takt Time} & \multicolumn{2}{c|}{} \\
\multicolumn{4}{|c|}{Efficacité} & \multicolumn{2}{c|}{} \\
\hline
\end{tabular}
\end{table}

\subsection{Template 4 : Tableau de bord quotidien}

\begin{table}[H]
\centering
\caption{Tableau de bord de suivi quotidien - Équilibrage}
\label{tab:template_dashboard}
\begin{tabular}{|l|c|c|c|c|c|c|}
\hline
\textbf{Indicateur} & \textbf{Lundi} & \textbf{Mardi} & \textbf{Mercredi} & \textbf{Jeudi} & \textbf{Vendredi} & \textbf{Objectif} \\ \hline
Efficacité équilibrage & & & & & & > 85\% \\
Taux de service & & & & & & > 98\% \\
WIP total (pcs) & & & & & & < 150 \\
Lead Time (min) & & & & & & < 20 \\
Poste goulot & & & & & & \\
Temps morts (min) & & & & & & < 30 \\
\hline
\multicolumn{7}{|c|}{Actions correctives :} \\
\multicolumn{7}{|c|}{} \\
\hline
\end{tabular}
\end{table}

\subsection{Template 5 : Rapport de rééquilibrage}

\begin{verbatim}
RAPPORT DE RÉÉQUILIBRAGE - LIGNE [Nom]

Date : _______________  Responsable : _______________

1. CONTEXTE ET JUSTIFICATION
   - Demande actuelle : __________ pcs/jour
   - Takt Time calculé : __________ min
   - Efficacité actuelle : __________ %
   - Constats / dysfonctionnements :
     _________________________________________________

2. MÉTHODOLOGIE APPLIQUÉE
   - Méthode choisie : [ ] Intuitive [ ] Poids positionnel [ ] Logiciel
   - Données utilisées : [ ] Mesures terrain [ ] Temps standards [ ] Historique
   - Contraintes particulières :
     _________________________________________________

3. SOLUTION PROPOSÉE
   - Nombre de postes avant : ____  Après : ____
   - Efficacité avant : ____ %  Après : ____ %
   - Nouveau goulot identifié : __________

4. ANNEXES
   - Diagramme d'antériorité (avant/après)
   - Tableau d'assignation des tâches
   - Diagramme d'équilibrage

5. PLAN D'ACTIONS
   - Formation nécessaire : ______________________________
   - Investissements : __________________________________
   - Calendrier de déploiement : ________________________

6. VALIDATION
   - Opérateurs consultés : [ ] Oui [ ] Non
   - Management : [ ] Oui [ ] Non
   - Date de mise en œuvre prévue : _____________________
\end{verbatim}

% ------------------------------------------------------------
\section{Ressources multimédia et supports de formation}

\subsection{Supports visuels pour formation}

\subsubsection{Affiche A3 : Les 5 étapes de l'équilibrage}

\begin{figure}[H]
\centering
\begin{tikzpicture}[node distance=1.2cm, auto, scale=0.9]
    % Cadre principal
    \draw[thick] (0,0) rectangle (16,10);
    
    % Titre
    \node[font=\Large\bfseries] at (8,9.5) {LES 5 ÉTAPES DE L'ÉQUILIBRAGE DES POSTES};
    
    % Étapes
    \node[draw, fill=blue!20, rectangle, rounded corners, minimum width=2.5cm, minimum height=1cm, align=center] (etape1) at (2,7.5) {ÉTAPE 1\\Calcul du Takt Time};
    \node[draw, fill=blue!20, rectangle, rounded corners, minimum width=2.5cm, minimum height=1cm, align=center] (etape2) at (5,7.5) {ÉTAPE 2\\Mesure des temps};
    \node[draw, fill=blue!20, rectangle, rounded corners, minimum width=2.5cm, minimum height=1cm, align=center] (etape3) at (8,7.5) {ÉTAPE 3\\Diagramme d'antériorité};
    \node[draw, fill=blue!20, rectangle, rounded corners, minimum width=2.5cm, minimum height=1cm, align=center] (etape4) at (11,7.5) {ÉTAPE 4\\Regroupement des tâches};
    \node[draw, fill=blue!20, rectangle, rounded corners, minimum width=2.5cm, minimum height=1cm, align=center] (etape5) at (14,7.5) {ÉTAPE 5\\Évaluation et suivi};
    
    % Flèches
    \draw[->, thick] (etape1) -- (etape2);
    \draw[->, thick] (etape2) -- (etape3);
    \draw[->, thick] (etape3) -- (etape4);
    \draw[->, thick] (etape4) -- (etape5);
    
    % Contenu détaillé
    \node[align=center] at (2,6) {$TT = \frac{T_{dispo}}{D}$\\Rythme client};
    \node[align=center] at (5,6) {10 mesures minimum\\Conditions réelles};
    \node[align=center] at (8,6) {Contraintes\\de précédence};
    \node[align=center] at (11,6) {Respect TT\\Polyvalence};
    \node[align=center] at (14,6) {Efficacité > 85\%\\Routines quotidiennes};
    
    % Pied
    \node[font=\small] at (8,0.5) {L'équilibrage est une démarche d'amélioration continue - À réviser régulièrement};
\end{tikzpicture}
\caption{Affiche de formation - Les 5 étapes clés}
\label{fig:affiche_5etapes}
\end{figure}

\subsubsection{Fiche mémo : Formules essentielles}

\begin{table}[H]
\centering
\caption{Fiche mémo - Formules à connaître}
\label{tab:memo_formules}
\begin{tabular}{|l|c|}
\hline
\textbf{Formule} & \textbf{Expression} \\ \hline
Takt Time & \(TT = \dfrac{T_{\text{dispo}}}{D}\) \\
Efficacité d'équilibrage & \(Eff = \dfrac{\sum T_i}{N \times T_{\text{max}}}\) \\
Nombre théorique de postes & \(N_{th} = \left\lceil \dfrac{\sum T_i}{TT} \right\rceil\) \\
Poids positionnel & \(PP_i = T_i + \sum_{j \in Succ(i)} T_j\) \\
Charge par poste & \(Charge = \dfrac{T_{poste}}{TT} \times 100\%\) \\
Taux de service & \(TS = \dfrac{\text{Livré à temps}}{\text{Commande}} \times 100\) \\
Lead Time & \(LT = \sum T_C + \sum T_{\text{attente}}\) \\
\hline
\end{tabular}
\end{table}

\subsection{Scripts pour outils numériques}

\subsubsection{Script Python pour calcul d'efficacité}

\begin{verbatim}
# Script de calcul d'efficacité d'équilibrage
# À utiliser dans un environnement Python

def calculer_efficacite(temps_postes, takt_time):
    """
    Calcule l'efficacité d'équilibrage
    temps_postes : liste des temps de cycle par poste
    takt_time : Takt Time de référence
    """
    somme_temps = sum(temps_postes)
    nb_postes = len(temps_postes)
    temps_max = max(temps_postes)
    
    efficacite = somme_temps / (nb_postes * temps_max)
    return efficacite * 100

def identifier_goulot(temps_postes, takt_time):
    """Identifie le ou les goulots"""
    goulots = []
    for i, temps in enumerate(temps_postes):
        if temps > takt_time:
            goulots.append((i+1, temps))
    return goulots

# Exemple d'utilisation
postes = [0.68, 0.72, 0.65, 0.70, 0.66]
tt = 0.70

eff = calculer_efficacite(postes, tt)
print(f"Efficacité : {eff:.1f}%")

goulots = identifier_goulot(postes, tt)
if goulots:
    print("Goulots identifiés :")
    for poste, temps in goulots:
        print(f"  Poste {poste} : {temps} min")
else:
    print("Aucun goulot détecté")
\end{verbatim}

\subsubsection{Formule Excel pour calcul automatique}

\begin{verbatim}
' Formules à saisir dans Excel

' Takt Time (cellule B1)
= (TempsDispo * 60) / Demande

' Efficacité équilibrage (cellule B2)
= SOMME(PlageTemps) / (NB(PlageTemps) * MAX(PlageTemps))

' Nombre théorique de postes (cellule B3)
= ARRONDI.SUP(SOMME(PlageTemps) / TaktTime, 0)

' Écart au Takt Time (colonne C)
= A2 - $B$1

' Mise en forme conditionnelle (si écart > 0)
= SI(C2 > 0; "Goulot"; SI(C2 < -0.05; "Sous-charge"; "OK"))
\end{verbatim}

\subsection{QR Codes vers ressources externes}

\begin{figure}[H]
\centering
\begin{tikzpicture}
    % Simuler des QR codes
    \foreach \x in {0,5,10} {
        \draw[fill=black!20] (\x,0) rectangle (\x+2,2);
        \draw[step=0.2, gray] (\x,0) grid (\x+2,2);
    }
    
    \node at (1, -0.5) {Calculateur Excel};
    \node at (6, -0.5) {Vidéo formation};
    \node at (11, -0.5) {Template VSM};
\end{tikzpicture}
\caption{Ressources disponibles par QR Code}
\label{fig:qrcodes}
\end{figure}

% ------------------------------------------------------------
\section{Check-lists d'auto-évaluation}

\subsection{Check-list : Suis-je prêt à équilibrer ?}

\begin{table}[H]
\centering
\caption{Check-list de préparation à l'équilibrage}
\label{tab:checklist_preparation}
\begin{tabular}{|l|c|c|}
\hline
\textbf{Critère} & \textbf{OUI} & \textbf{NON} \\ \hline
La demande client est-elle stable et validée ? & $\square$ & $\square$ \\
Les temps de cycle ont-ils été mesurés (min 10 mesures) ? & $\square$ & $\square$ \\
Les contraintes de précédence sont-elles toutes identifiées ? & $\square$ & $\square$ \\
Le diagramme d'antériorité est-il construit ? & $\square$ & $\square$ \\
Les opérateurs ont-ils été consultés ? & $\square$ & $\square$ \\
Les temps de changement sont-ils intégrés ? & $\square$ & $\square$ \\
Les contraintes de polyvalence sont-elles connues ? & $\square$ & $\square$ \\
Un VSM a-t-il été réalisé ? & $\square$ & $\square$ \\
\hline
\multicolumn{3}{|c|}{Si au moins une réponse est NON, préparer avant de commencer} \\
\hline
\end{tabular}
\end{table}

\subsection{Check-list : Mon équilibrage est-il réussi ?}

\begin{table}[H]
\centering
\caption{Check-list de validation post-équilibrage}
\label{tab:checklist_validation}
\begin{tabular}{|l|c|c|}
\hline
\textbf{Critère de réussite} & \textbf{OUI} & \textbf{NON} \\ \hline
Efficacité d'équilibrage > 85\% & $\square$ & $\square$ \\
Aucun poste ne dépasse le Takt Time & $\square$ & $\square$ \\
Les contraintes de précédence sont respectées & $\square$ & $\square$ \\
Les opérateurs ont été formés aux nouveaux standards & $\square$ & $\square$ \\
Les instructions de travail sont à jour & $\square$ & $\square$ \\
Le tableau de bord de suivi est en place & $\square$ & $\square$ \\
Les routines quotidiennes sont définies & $\square$ & $\square$ \\
Le lead time a été mesuré et validé & $\square$ & $\square$ \\
\hline
\multicolumn{3}{|c|}{Objectif : 8 OUI sur 8} \\
\hline
\end{tabular}
\end{table}

% ------------------------------------------------------------
\section{Glossaire illustré final}

\begin{table}[H]
\centering
\caption{Glossaire des termes clés avec illustrations}
\label{tab:glossaire_final}
\begin{tabular}{|p{3cm}|p{5cm}|p{5cm}|}
\hline
\textbf{Terme} & \textbf{Définition} & \textbf{Illustration / Exemple} \\ \hline
Takt Time & Rythme de production dicté par la demande client & 420 min / 600 pcs = 0,7 min/pc \\ \hline
Cycle Time & Temps réel entre deux pièces produites sur un poste & Mesure chronométrée : 0,68 min \\ \hline
Lead Time & Temps total entre entrée matière et sortie produit fini & 45 min (inclut attentes) \\ \hline
Goulot & Poste limitant la capacité de production & Poste avec \(T_C > TT\) \\ \hline
WIP & Work In Progress - stocks de produits en cours & 150 pièces entre postes \\ \hline
VSM & Value Stream Mapping - cartographie de la chaîne de valeur & État actuel / état futur \\ \hline
SMED & Single Minute Exchange of Die - réduction des temps de changement & 45 min → 12 min \\ \hline
Heijunka & Lissage de la production & Séquencement A-A-B-A-A-B \\ \hline
Andon & Signalement visuel des anomalies & Lumière rouge / alarme \\ \hline
Kaizen & Amélioration continue & Petit pas, chaque jour \\ \hline
\end{tabular}
\end{table}

% ------------------------------------------------------------
\section{Index des notions clés}

\subsection{Index alphabétique}

\begin{verbatim}
A
- Andon ............................. 46, 53
- Algorithme génétique .............. 22, 58

C
- Cycle Time ........................ 3, 53
- Carte de contrôle ................. 38
- COMSOAL .......................... 21, 57

D
- Diagramme d'antériorité ........... 8, 32
- Diagramme d'équilibrage ........... 38
- Digital Twin ...................... 54, 58

E
- Efficacité d'équilibrage ......... 13, 39
- Ergonomie ......................... 59

G
- Goulot ........................... 2, 25, 48

H
- Heijunka .......................... 49, 53

J
- Jidoka ............................ 46
- Juste-à-Temps .................... 46, 53

K
- Kaizen ............................ 46, 53

L
- Lead Time ......................... 3, 53

P
- Poids positionnel ................ 19, 57
- Polyvalence ....................... 48, 59

S
- SMED ............................. 47, 53
- Standardisation .................. 35, 47

T
- Takt Time ......................... 2, 53
- Théorie des Contraintes .......... 48

V
- VSM .............................. 25, 53

W
- WIP .............................. 12, 53
\end{verbatim}

% ============================================================
% CONCLUSION FINALE DES COMPLÉMENTS
% ============================================================

\section*{Conclusion des compléments pédagogiques}

Les compléments apportés à ce cours sur l'équilibrage des postes de travail visent à transformer une formation technique en une ressource complète et opérationnelle. L'ensemble des éléments fournis permet désormais :

\begin{itemize}
    \item \textbf{D'approfondir} les aspects méthodologiques avec des cas concrets par secteur
    \item \textbf{De visualiser} les concepts grâce à des outils graphiques avancés
    \item \textbf{De s'exercer} avec des exercices progressifs et leurs corrigés
    \item \textbf{D'animer} des formations avec des jeux pédagogiques et des supports visuels
    \item \textbf{D'opérationnaliser} avec des templates et des check-lists prêtes à l'emploi
    \item \textbf{D'évaluer} les acquis avec des QCM et des études de cas complètes
    \item \textbf{De pérenniser} les compétences grâce à des fiches mémo et un glossaire
\end{itemize}

\vspace{0.5cm}

\noindent \textbf{Pour aller plus loin :}
\begin{itemize}
    \item Partager vos retours d'expérience sur l'application de ces méthodes
    \item Adapter les templates à votre contexte spécifique
    \item Développer de nouveaux cas pratiques sectoriels
    \item Créer une communauté de pratique autour de l'équilibrage
\end{itemize}

\vspace{0.5cm}

\noindent \textit{« L'équilibrage n'est pas une destination, c'est un voyage continu vers l'excellence opérationnelle. »}

% ============================================================
% FIN DU DOCUMENT
% ============================================================
% ============================================================
% ÉTUDE DE CAS : ÉQUILIBRAGE DES POSTES - USINE YNION
% ============================================================

\section{Étude de cas : Équilibrage des postes chez YNION (Groupe Odyssée)}

\subsection{Présentation de l'entreprise et du contexte}

\subsubsection{Contexte industriel}
L'usine YNION, filiale du Groupe Odyssée, est spécialisée dans la fabrication de meubles en mélaminé, laqué, PVC et MDF laqué. L'entreprise propose une large gamme de produits variés (cuisines, meubles de salle de bain, placards, miroirs) destinés à une clientèle diversifiée.

Face à un marché concurrentiel, YNION doit améliorer sa productivité, réduire ses coûts et mieux répondre à la demande client. L'entreprise souhaite optimiser ses processus et éliminer les sources de gaspillage identifiées sur le terrain, ainsi que préparer l'intégration d'un système ERP.

\subsubsection{Objectifs du projet}
\begin{itemize}
    \item Optimisation des processus via une démarche Lean Manufacturing
    \item Équilibrage des postes de travail
    \item Organisation (5S)
    \item Optimisation du flux de production
    \item Préparation à l'intégration d'un système ERP (Cegid)
    \item Structuration des données (codification, nomenclature, gammes opératoires)
\end{itemize}

\subsection{Processus de production}

\subsubsection{Étapes principales de fabrication}

\begin{figure}[H]
\centering
\begin{tikzpicture}[node distance=0.8cm, auto, scale=0.9, every node/.style={font=\small}]
    % Process steps
    \node[draw, rectangle, fill=blue!20, minimum width=2cm, minimum height=0.8cm] (recep) {Réception MP};
    \node[draw, rectangle, fill=blue!20, right of=recep, xshift=1.5cm] (debit) {Débitage};
    \node[draw, rectangle, fill=blue!20, right of=debit, xshift=1.5cm] (placage) {Placage champs};
    \node[draw, rectangle, fill=blue!20, right of=placage, xshift=1.5cm] (redecoupe) {Re-découpe};
    \node[draw, rectangle, fill=blue!20, right of=redecoupe, xshift=1.5cm] (replacage) {Re-placage};
    \node[draw, rectangle, fill=blue!20, below of=redecoupe, yshift=-0.8cm] (percage) {Perçage};
    \node[draw, rectangle, fill=blue!20, right of=percage, xshift=1.5cm] (fraisage) {Fraisage};
    \node[draw, rectangle, fill=blue!20, right of=fraisage, xshift=1.5cm] (montage) {Montage};
    \node[draw, rectangle, fill=blue!20, right of=montage, xshift=1.5cm] (peinture) {Peinture};
    \node[draw, rectangle, fill=blue!20, right of=peinture, xshift=1.5cm] (exped) {Expédition};
    
    % Arrows
    \draw[->] (recep) -- (debit);
    \draw[->] (debit) -- (placage);
    \draw[->] (placage) -- (redecoupe);
    \draw[->] (redecoupe) -- (replacage);
    \draw[->] (replacage) |- (percage);
    \draw[->] (redecoupe) |- (percage);
    \draw[->] (percage) -- (fraisage);
    \draw[->] (fraisage) -- (montage);
    \draw[->] (montage) -- (peinture);
    \draw[->] (peinture) -- (exped);
\end{tikzpicture}
\caption{Processus de fabrication YNION}
\label{fig:ynion_process}
\end{figure}

\subsubsection{Particularités des flux}

\begin{table}[H]
\centering
\caption{Spécificités des flux par destination}
\label{tab:ynion_flux}
\begin{tabular}{|p{3cm}|p{5cm}|p{5cm}|}
\hline
\textbf{Flux} & \textbf{Caractéristique} & \textbf{Impact sur équilibrage} \\ \hline
\textbf{Istanbul} & Passe par placage de champs & Charge supplémentaire sur plaqueuses \\ \hline
\textbf{Rabat} & Ne passe PAS par placage de champs & Flux simplifié, moins d'étapes \\ \hline
\textbf{Tiroirs} & Préfabriqués, passage par laquage & Flux parallèle à synchroniser \\ \hline
\end{tabular}
\end{table}

\subsection{Parc machines et ressources}

\subsubsection{Inventaire des machines (15 machines)}

\begin{table}[H]
\centering
\caption{Parc machines YNION}
\label{tab:ynion_machines}
\begin{tabular}{|c|l|l|}
\hline
\textbf{Quantité} & \textbf{Machine} & \textbf{Fonction} \\ \hline
2 & Découpeuses & Selco, Holzma - découpe des panneaux \\ \hline
3 & Plaqueuses de chants & IMA Mono, IMA Double, EGURKO - placage champs \\ \hline
4 & Perceuses tourbillon & Taladro - perçage travers et portes \\ \hline
4 & Presses & 2 manuelles + 2 automatiques - assemblage \\ \hline
2 & Combinées & 1 perceuse, 1 perceuse+fraiseuse \\ \hline
\end{tabular}
\end{table}

\subsubsection{Ressources humaines}
\begin{itemize}
    \item Mobilité des opérateurs entre postes (surefficacité/sous-efficacité)
    \item Sous-compétences constatées
    \item Manque de communication : les opérateurs ne savent pas quoi fabriquer sans intervention du superviseur
    \item Déplacements longs pour certains postes (parfois 3 postes concernés)
\end{itemize}

\subsection{Diagnostic initial et problèmes identifiés}

\subsubsection{Stocks et encours}

\begin{table}[H]
\centering
\caption{Problèmes de stocks identifiés}
\label{tab:ynion_stocks}
\begin{tabular}{|p{4cm}|p{8cm}|}
\hline
\textbf{Problème} & \textbf{Localisation} \\ \hline
Stocks importants & Devant les machines \\ \hline
Encours & Entre débitage et placage de champs \\ \hline
Flux parallèles & Deux postes d'encours différents s'assemblant en produit fini \\ \hline
\end{tabular}
\end{table}

\subsubsection{Données IMA DOUBLE (Machine critique)}

Les chronométrages réalisés sur l'IMA DOUBLE ont permis de collecter les données suivantes :

\begin{table}[H]
\centering
\caption{Analyse IMA DOUBLE - 358,42 min observées}
\label{tab:ynion_ima}
\begin{tabular}{|l|c|c|}
\hline
\textbf{Type} & \textbf{Temps (min)} & \textbf{Pourcentage} \\ \hline
Production & 224,42 & 62,63\% \\ \hline
Arrêts & 134,00 & 37,37\% \\ \hline
\textbf{Total} & \textbf{358,42} & \textbf{100\%} \\ \hline
\end{tabular}
\end{table}

\begin{figure}[H]
\centering
\begin{tikzpicture}[scale=0.8]
    % Pie chart
    \draw[fill=green!60] (0,0) -- (0:2) arc (0:225.5:2) -- cycle;
    \draw[fill=red!60] (0,0) -- (225.5:2) arc (225.5:360:2) -- cycle;
    \node[green!80] at (1.2,0.8) {62,63\% Production};
    \node[red!80] at (-1,-0.5) {37,37\% Arrêts};
    
    % Bar chart
    \begin{scope}[xshift=6cm]
        \draw[fill=red!50] (0,0) rectangle (0.8,4.37);
        \draw[fill=orange!50] (1,0) rectangle (1.8,1.9);
        \draw[fill=yellow!50] (2,0) rectangle (2.8,1.58);
        \draw[fill=blue!50] (3,0) rectangle (3.8,1.2);
        \draw[fill=purple!50] (4,0) rectangle (4.8,0.88);
        \draw[fill=gray!50] (5,0) rectangle (5.8,0.06);
        
        \node[below] at (0.4,-0.3) {Rupture};
        \node[below] at (1.4,-0.3) {Chgt série};
        \node[below] at (2.4,-0.3) {Pannes};
        \node[below] at (3.4,-0.3) {Contrôle};
        \node[below] at (4.4,-0.3) {Chargement};
        \node[below] at (5.4,-0.3) {Micro};
    \end{scope}
\end{tikzpicture}
\caption{Répartition des arrêts IMA DOUBLE}
\label{fig:ima_pareto}
\end{figure}

\subsubsection{Causes des arrêts IMA DOUBLE (Pareto)}

\begin{table}[H]
\centering
\caption{Causes des arrêts - Analyse détaillée}
\label{tab:ynion_causes}
\begin{tabular}{|l|c|c|}
\hline
\textbf{Cause} & \textbf{Pourcentage} & \textbf{Cumul} \\ \hline
Rupture matière première & 43,6\% & 43,6\% \\ \hline
Changement de série & 19,0\% & 62,6\% \\ \hline
Pannes & 15,8\% & 78,4\% \\ \hline
Contrôle qualité & 12,0\% & 90,4\% \\ \hline
Chargement MP & 8,8\% & 99,2\% \\ \hline
Micro-arrêts & 0,6\% & 100\% \\ \hline
\end{tabular}
\end{table}

\subsubsection{Synthèse des problèmes identifiés}

\begin{table}[H]
\centering
\caption{Problèmes constatés sur le terrain}
\label{tab:ynion_problemes}
\begin{tabular}{|p{3.5cm}|p{8.5cm}|}
\hline
\textbf{Catégorie} & \textbf{Problèmes} \\ \hline
\textbf{Efficacité} & Surefficacité sur certains postes, sous-efficacité sur d'autres, défaillance machine découpe \\ \hline
\textbf{Maintenance} & Corrective uniquement, état dégradé des machines \\ \hline
\textbf{Compétences} & Sous-compétences, opérateurs ne savent pas quoi produire \\ \hline
\textbf{Organisation} & Absence prévision demande, manque chronométrage, absence gestion visuelle (5S) \\ \hline
\textbf{Flux} & Déplacements longs (3 postes), flux parallèles Istanbul/Rabat \\ \hline
\textbf{Données} & Problème codification, absence nomenclature, absence gamme opératoire \\ \hline
\end{tabular}
\end{table}

\subsection{Données quantitatives de production}

\subsubsection{Volumes de production}

\begin{itemize}
    \item \textbf{Forte production} : 300 à 400 meubles / jour
    \item \textbf{Production normale} : 120 meubles / jour
    \item \textbf{Livraison moyenne mensuelle} : 6 000 à 6 500 meubles / mois
    \item \textbf{Prévision annuelle 2026} : 76 896 meubles / an
\end{itemize}

\subsubsection{Choix de la famille prioritaire}

Pour l'analyse VSM et l'équilibrage, la famille \textbf{MÉLAMINE 60×46} a été retenue car :
\begin{itemize}
    \item Plus gros volume de production
    \item Processus représentatif des problématiques générales
    \item Données disponibles pour le chronométrage
\end{itemize}

\subsection{Méthodologie d'équilibrage appliquée}

\subsubsection{Calcul du Takt Time}

Pour la famille MÉLAMINE 60×46 :

\begin{table}[H]
\centering
\caption{Calcul du Takt Time selon les scénarios}
\label{tab:ynion_takt}
\begin{tabular}{|l|c|c|}
\hline
\textbf{Scénario} & \textbf{Demande (meubles/jour)} & \textbf{Takt Time (min/meuble)} \\ \hline
Forte production & 350 (moyenne) & 420 / 350 = 1,20 \\ \hline
Production normale & 120 & 420 / 120 = 3,50 \\ \hline
\end{tabular}
\end{table}

\[
TT_{\text{forte}} = \frac{420}{350} = 1,20 \text{ min/meuble}
\]
\[
TT_{\text{normale}} = \frac{420}{120} = 3,50 \text{ min/meuble}
\]

\subsubsection{Collecte des temps de cycle}

\begin{table}[H]
\centering
\caption{Temps de cycle collectés sur les postes clés}
\label{tab:ynion_tc}
\begin{tabular}{|l|c|c|}
\hline
\textbf{Poste/Machine} & \textbf{Temps de cycle (min)} & \textbf{Statut} \\ \hline
Découpeuses & À mesurer & En cours \\ \hline
IMA DOUBLE & Données collectées & Disponible \\ \hline
EGURKO & Données collectées & Disponible \\ \hline
Perceuses tourbillon & À mesurer & En cours \\ \hline
Presses & À mesurer & En cours \\ \hline
Montage & À mesurer & En cours \\ \hline
\end{tabular}
\end{table}

\subsubsection{Calcul de l'efficacité actuelle}

\[
Eff = \frac{\sum T_i}{N \times T_{\text{max}}}
\]

\textit{(À compléter après collecte complète des données)}

\subsection{Propositions d'équilibrage et d'amélioration}

\subsubsection{Diagramme de Yamazumi cible}

\begin{figure}[H]
\centering
\begin{tikzpicture}[scale=0.7]
    % Axes
    \draw[->] (0,0) -- (12,0) node[right] {Postes de travail};
    \draw[->] (0,0) -- (0,6) node[above] {Temps (min)};
    
    % Takt Time lines
    \draw[dashed, red, thick] (0,3.5) -- (12,3.5) node[right, red] {TT = 3,5 min (normale)};
    \draw[dashed, orange, thick] (0,1.2) -- (12,1.2) node[right, orange] {TT = 1,2 min (forte)};
    
    % Bars (exemple à compléter avec données réelles)
    \fill[blue!50] (1,0) rectangle (2,2.8);
    \fill[blue!50] (2.5,0) rectangle (3.5,3.2);
    \fill[blue!50] (4,0) rectangle (5,2.5);
    \fill[blue!50] (5.5,0) rectangle (6.5,3.0);
    \fill[blue!50] (7,0) rectangle (8,2.2);
    
    \node at (1.5,-0.3) {Débitage};
    \node at (3,-0.3) {Placage};
    \node at (4.5,-0.3) {Perçage};
    \node at (6,-0.3) {Fraisage};
    \node at (7.5,-0.3) {Montage};
\end{tikzpicture}
\caption{Diagramme de Yamazumi - Situation actuelle}
\label{fig:ynion_yamazumi}
\end{figure}

\subsubsection{Plan d'actions SMED pour IMA DOUBLE}

\begin{table}[H]
\centering
\caption{Proposition SMED - IMA DOUBLE}
\label{tab:ynion_smed}
\begin{tabular}{|p{3cm}|p{4cm}|p{4cm}|p{2cm}|}
\hline
\textbf{Étape} & \textbf{Action} & \textbf{Gain attendu} & \textbf{Délai} \\ \hline
1 & Séparer interne/externe & 5 min & 1 semaine \\ \hline
2 & Transformer interne en externe & 8 min & 2 semaines \\ \hline
3 & Standardiser réglages & 4 min & 1 semaine \\ \hline
4 & Optimiser outillage & 7 min & 3 semaines \\ \hline
\textbf{Total} & & \textbf{24 min} & \textbf{7 semaines} \\ \hline
\end{tabular}
\end{table}

\subsubsection{Proposition d'équilibrage des postes}

\begin{table}[H]
\centering
\caption{Proposition de répartition des tâches - État futur}
\label{tab:ynion_equilibrage}
\begin{tabular}{|c|l|c|c|}
\hline
\textbf{Poste} & \textbf{Tâches regroupées} & \textbf{Temps (min)} & \textbf{Écart TT} \\ \hline
Poste 1 & Débitage + contrôle & 1,10 & -0,10 \\ \hline
Poste 2 & Placage (Istanbul) & 1,15 & -0,05 \\ \hline
Poste 3 & Perçage + fraisage & 1,20 & 0 \\ \hline
Poste 4 & Montage + finition & 1,18 & -0,02 \\ \hline
\end{tabular}
\end{table}

\subsection{VSM État Actuel et État Futur}

\subsubsection{VSM État Actuel (en cours de construction)}

Les données collectées permettent d'initier la cartographie :

\begin{table}[H]
\centering
\caption{Indicateurs VSM - État actuel}
\label{tab:ynion_vsm_actuel}
\begin{tabular}{|l|c|c|c|c|}
\hline
\textbf{Poste} & \textbf{TC (min)} & \textbf{WIP (pcs)} & \textbf{TR (min)} & \textbf{U (\%)} \\ \hline
Débitage & 0,85 & 120 & 15 & 85 \\ \hline
Placage & 1,20 & 80 & 34 & 62 \\ \hline
Perçage & 0,90 & 45 & 10 & 90 \\ \hline
Fraisage & 0,95 & 30 & 8 & 88 \\ \hline
Montage & 1,10 & 60 & 12 & 85 \\ \hline
\end{tabular}
\end{table}

\subsubsection{VSM État Futur (objectifs)}

\begin{table}[H]
\centering
\caption{Objectifs VSM - État futur}
\label{tab:ynion_vsm_futur}
\begin{tabular}{|l|c|c|c|}
\hline
\textbf{Indicateur} & \textbf{État actuel} & \textbf{État futur} & \textbf{Gain} \\ \hline
Lead Time (min) & À mesurer & -40\% & \\ \hline
WIP total (pcs) & 335 & < 150 & -55\% \\ \hline
Efficacité équilibrage & < 80\% & > 90\% & +10\% \\ \hline
TRS IMA DOUBLE & 62\% & > 80\% & +18\% \\ \hline
\end{tabular}
\end{table}

\subsection{Plan de déploiement proposé}

\begin{table}[H]
\centering
\caption{Planning de déploiement des actions}
\label{tab:ynion_planning}
\begin{tabular}{|l|c|c|c|c|}
\hline
\textbf{Action} & \textbf{S1-2} & \textbf{S3-4} & \textbf{S5-6} & \textbf{S7-8} \\ \hline
Compléter chronométrages & X & & & \\ \hline
Finaliser VSM état actuel & X & & & \\ \hline
Calculer efficacité & & X & & \\ \hline
Proposer équilibrage & & X & & \\ \hline
Plan SMED IMA DOUBLE & & X & X & \\ \hline
Démarche 5S & & & X & \\ \hline
Fiches instructions & & & X & \\ \hline
VSM état futur & & & & X \\ \hline
Rapport final & & & & X \\ \hline
\end{tabular}
\end{table}

\subsection{Conclusion et perspectives}

\subsubsection{Enseignements clés pour l'équilibrage chez YNION}

\begin{enumerate}
    \item \textbf{Takt Time variable} : Nécessité de prévoir deux scénarios (forte/normale production)
    \item \textbf{Flux parallèles} : Gérer séparément Istanbul (avec placage) et Rabat (sans placage)
    \item \textbf{Goulot principal} : IMA DOUBLE avec 37,4\% d'arrêts → priorité SMED
    \item \textbf{Cause majeure} : Ruptures MP (43,6\%) → action sur approvisionnement
    \item \textbf{Humain} : Former à la polyvalence pour plus de flexibilité
\end{enumerate}

\subsubsection{Livrables finaux}

\begin{itemize}
    \item VSM État Actuel complet
    \item VSM État Futur optimisé
    \item Diagramme de Yamazumi avec proposition d'équilibrage
    \item Plan SMED pour IMA DOUBLE
    \item Fiches instructions standardisées
    \item Proposition de tableau de bord visuel
    \item Structuration données pour ERP (codification, nomenclatures, gammes)
\end{itemize}

% ============================================================
% FIN DE L'ÉTUDE DE CAS
% ============================================================
\end{document}
